
\clearpage
\chapter{Marco Teórico}

\addcontentsline{toc}{chapter}{Marco Teórico}

\section{Transferencia de calor y masa en secadores solares}

El funcionamiento de un secador solar de café tipo invernadero se fundamenta en la interacción simultánea de fenómenos de transferencia de calor y transferencia de masa. La energía necesaria para el secado proviene principalmente de la radiación solar incidente, la cual atraviesa la cubierta transparente del secador y es absorbida por las superficies internas y por el producto. Como consecuencia, la temperatura del aire contenido en el interior aumenta, reduciendo su humedad relativa e incrementando su capacidad para absorber vapor de agua proveniente de los granos de café.

Desde el punto de vista termodinámico, el proceso de secado puede interpretarse como un fenómeno acoplado de transferencia de calor y masa. Inicialmente, el aire recibe energía mediante radiación y convección, elevando su temperatura. Posteriormente, esta energía es transferida hacia el café principalmente por convección, suministrando el calor latente necesario para evaporar el agua contenida en el grano. Finalmente, el vapor de agua generado es transportado por el flujo de aire hacia el exterior del secador mediante procesos convectivos.

Los mecanismos dominantes de transferencia de calor presentes en un secador solar se resumen en la Tabla \ref{tab:heattransfer}.

\begin{table}[h!]
\centering
\begin{tabular}{p{4cm}p{10cm}}
\toprule
\textbf{Mecanismo} & \textbf{Descripción} \\
\midrule
\\
\textbf{Radiación} & La radiación solar constituye la principal fuente energética del sistema, calentando la cubierta, las paredes, las bandejas y el café.\\
\\
\textbf{Conducción} & El calor se propaga a través de los materiales sólidos del secador, incluyendo la cubierta, la estructura y las bandejas.\\
\\
\textbf{Convección natural} & Las diferencias de temperatura generan fuerzas de flotación que inducen el movimiento ascendente del aire caliente.\\
\\
\textbf{Convección forzada} & Puede producirse mediante ventiladores u otros actuadores destinados a incrementar el intercambio de calor y humedad.\\
\\
\bottomrule
\end{tabular}
\caption{Principales mecanismos de transferencia de calor presentes en un secador solar de café.}
\label{tab:heattransfer}
\end{table}

El flujo de calor por conducción en los materiales sólidos está gobernado por la Ley de Fourier,

\begin{equation}
\mathbf{q}=-k\nabla T,
\end{equation}

donde $\mathbf{q}$ representa el flujo de calor, $k$ la conductividad térmica del material y $\nabla T$ el gradiente de temperatura.

Por otra parte, el intercambio de calor entre el aire y la superficie del café puede describirse mediante la ley de enfriamiento de Newton,

\begin{equation}
\dot Q=hA(T_s-T_\infty),
\end{equation}

donde $h$ es el coeficiente convectivo, $A$ el área de intercambio, $T_s$ la temperatura superficial del producto y $T_\infty$ la temperatura del aire circundante.

La evaporación del agua contenida en el café constituye simultáneamente un fenómeno de transferencia de masa. El aire caliente reduce su humedad relativa, aumentando la diferencia de potencial de vapor entre la superficie del grano y el aire circundante, condición necesaria para favorecer la evaporación. De forma simplificada, el flujo másico evaporado puede expresarse como

\begin{equation}
\dot m_w=h_mA\left(C_s-C_\infty\right),
\end{equation}

donde $h_m$ representa el coeficiente convectivo de transferencia de masa, mientras que $C_s$ y $C_\infty$ corresponden a las concentraciones de vapor de agua en la superficie del producto y en el aire, respectivamente.

La circulación del aire dentro del secador ocurre principalmente por convección natural. Conforme el aire incrementa su temperatura, su densidad disminuye y aparecen fuerzas de flotación que inducen un movimiento ascendente. Como consecuencia, el aire caliente y húmedo abandona el secador a través de las aberturas superiores, mientras que aire más frío y seco ingresa por las aberturas inferiores, estableciendo un ciclo continuo de renovación de aire.

Desde la perspectiva del modelado matemático, esta interacción entre transferencia de calor, transferencia de masa y dinámica de fluidos convierte al secador solar en un sistema multifísico altamente acoplado. La temperatura, la humedad y la velocidad del aire evolucionan simultáneamente y dependen de múltiples perturbaciones externas, tales como la irradiancia solar, la temperatura ambiente, la humedad relativa, la velocidad del viento y la nubosidad. En consecuencia, el análisis del proceso requiere resolver de manera conjunta las ecuaciones de conservación de masa, cantidad de movimiento, energía y transporte de especies, aspecto que motiva el empleo de modelos basados en dinámica de fluidos computacional (CFD) y modelos térmicos reducidos desarrollados en las secciones siguientes.

Finalmente, una de las principales ventajas de los secadores solares tipo invernadero respecto al secado tradicional en patios consiste en un uso más eficiente del espacio disponible. Mientras que el secado a la intemperie requiere extender el café sobre grandes superficies para garantizar una exposición uniforme a la radiación solar, los secadores tipo invernadero permiten distribuir el producto en múltiples bandejas ubicadas en diferentes niveles. Esta configuración incrementa significativamente la capacidad de procesamiento por unidad de área, facilita la automatización del proceso y reduce la exposición del producto a contaminantes ambientales y eventos meteorológicos adversos.


\subsection{Conducción, convección y radiación}
La transferencia de calor en un secador solar de café ocurre mediante la acción simultánea de tres mecanismos fundamentales: conducción, convección y radiación. La contribución relativa de cada uno depende de la geometría del secador, las propiedades térmicas de los materiales, la intensidad de la radiación solar incidente y las condiciones ambientales. Debido a que estos mecanismos interactúan continuamente entre sí, el proceso de secado constituye un problema de transferencia de calor acoplada.

\subsubsection{Transferencia de calor por conducción}

La conducción corresponde al transporte microscópico de energía térmica a través de un medio material debido a la existencia de gradientes de temperatura. En un secador solar este fenómeno ocurre principalmente en la cubierta transparente, la estructura metálica, los materiales aislantes, las bandejas y el propio grano de café.

El flujo de calor por conducción se encuentra gobernado por la Ley de Fourier,

\begin{equation}
\mathbf{q}=-k\nabla T,
\label{eq:fourier}
\end{equation}

donde $\mathbf{q}$ representa el flujo de calor por unidad de área ($W/m^2$), $k$ es la conductividad térmica del material ($W/mK$) y $\nabla T$ corresponde al gradiente espacial de temperatura.

Cuando el fenómeno presenta dependencia temporal, la distribución de temperatura dentro de un sólido puede describirse mediante la ecuación general de difusión del calor,

\begin{equation}
\rho c_p\frac{\partial T}{\partial t}
=
\nabla\cdot(k\nabla T)+\dot q,
\label{eq:heatdiffusion}
\end{equation}

donde $\rho$ es la densidad del material, $c_p$ el calor específico y $\dot q$ representa una generación volumétrica de calor.

En el contexto de los secadores solares, este mecanismo resulta fundamental para modelar el almacenamiento de energía en la estructura y el intercambio térmico entre las superficies sólidas y el aire contenido en el interior.

\subsubsection{Transferencia de calor por convección}

La convección describe el intercambio de calor entre una superficie sólida y un fluido en movimiento. En un secador solar, este mecanismo gobierna la transferencia de energía entre el aire caliente y los granos de café, así como entre el aire y las paredes internas del secador.

El flujo convectivo suele modelarse mediante la ley de enfriamiento de Newton,

\begin{equation}
\dot Q=hA(T_s-T_\infty),
\label{eq:newton}
\end{equation}

donde $h$ es el coeficiente convectivo de transferencia de calor, $A$ el área de intercambio, $T_s$ la temperatura de la superficie y $T_\infty$ la temperatura del fluido.

Dependiendo del mecanismo responsable del movimiento del aire, la convección puede clasificarse en:

\begin{itemize}
\item \textbf{Convección natural}, originada por diferencias de densidad producidas por gradientes térmicos (efecto de flotación o \textit{buoyancy}).

\item \textbf{Convección forzada}, producida mediante ventiladores, extractores u otros dispositivos mecánicos.
\end{itemize}

En la presente investigación, el fenómeno dominante corresponde a la convección natural, ya que el movimiento del aire es generado principalmente por las diferencias de temperatura entre el interior y el exterior del secador. No obstante, también se evalúan estrategias de control que incorporan elementos de ventilación forzada como posibles actuadores.

\subsubsection{Transferencia de calor por radiación}

La radiación constituye la principal fuente de energía del secador solar. A diferencia de la conducción y la convección, este mecanismo no requiere un medio material para transportar energía, sino que ocurre mediante ondas electromagnéticas emitidas por el Sol.

La potencia radiativa emitida por una superficie puede aproximarse mediante la ley de Stefan-Boltzmann,

\begin{equation}
q_r=\varepsilon\sigma T^4,
\label{eq:stefan}
\end{equation}

donde $\varepsilon$ representa la emisividad de la superficie y $\sigma=5.67\times10^{-8}\;W/m^2K^4$ corresponde a la constante de Stefan-Boltzmann.

En aplicaciones de secadores solares resulta de especial interés la radiación solar incidente. Siguiendo la metodología propuesta por el \textit{ASHRAE Handbook}, la irradiancia total puede descomponerse en tres componentes principales:

\begin{enumerate}
\item Radiación solar directa.
\item Radiación difusa proveniente del cielo.
\item Radiación reflejada por el terreno circundante.
\end{enumerate}

Esta descomposición permite calcular con mayor precisión la carga térmica absorbida por las distintas superficies del secador y constituye la base de los modelos de carga solar implementados en herramientas CFD modernas. En la presente investigación, este aspecto resulta especialmente relevante debido a que la radiación solar representa la principal fuente energética del sistema y determina gran parte de la dinámica térmica observada. :contentReference[oaicite:0]{index=0}

\subsubsection{Interacción entre los mecanismos}

Los tres mecanismos descritos no actúan de manera independiente. La radiación solar incrementa la temperatura de la cubierta y de las superficies internas del secador; posteriormente, el calor es transferido por conducción a través de los materiales sólidos y finalmente intercambiado con el aire mediante convección. Este aire caliente reduce su humedad relativa y favorece la evaporación del agua contenida en el café, acoplando de manera directa la transferencia de calor con la transferencia de masa.

En consecuencia, la descripción matemática completa del secador requiere resolver simultáneamente las ecuaciones de conservación de masa, cantidad de movimiento, energía y transporte de especies, razón por la cual el empleo de modelos de dinámica de fluidos computacional constituye una herramienta apropiada para representar el comportamiento físico del sistema.
\subsection{Ecuación de energía en recintos cerrados}
\subsection{Revisión de literatura relevante}

Para aprovechar de manera óptima la energía solar en los secadores de café, se propone ajustar el diseño del secador en función de la posición del sol durante la temporada de cosecha. Esto implicaría adaptar el secador para maximizar la captación de radiación solar en los momentos más efectivos del día y del año. Además, considerar geometrías de secadores solares que sean fácilmente modificables para cambiar su orientación según la época o dirección del sol podría ofrecer una ventaja significativa frente a secadores con geometrías fijas y no ajustables. No obstante, esta adaptabilidad debe ser equilibrada con el diseño para minimizar la resistencia al viento. Por lo tanto, sería necesario analizar el viento como un factor de riesgo en función de las geometrías del secador, asegurando que las modificaciones en el diseño no incrementen la vulnerabilidad al viento, que podría afectar la eficiencia y estabilidad del secado.\\

Se recomienda dejar un espacio de 40 a 70 cm justo por encima del suelo sin cubrir y tener algunas aberturas cerca del techo para facilitar la circulación del aire. En el interior del secador se colocan plataformas de madera y malla de acero para esparcir los granos. Es posible tener uno, dos o tres niveles, según las necesidades. Se sugiere mantener un espacio de al menos 50 cm entre cada nivel para garantizar un acceso fácil y permitir una adecuada circulación de aire.\\

\section{Dinámica de fluidos en cavidades cerradas}
El comportamiento del aire dentro de un secador solar de café constituye el principal mecanismo responsable del transporte de energía térmica y vapor de agua durante el proceso de secado. Desde el punto de vista físico, el secador puede interpretarse como una cavidad parcialmente cerrada en cuyo interior el aire experimenta movimientos inducidos por diferencias de densidad generadas por el calentamiento solar.

A diferencia de aplicaciones convencionales de ventilación forzada, el flujo interno de un secador solar es gobernado principalmente por fenómenos de flotación (\textit{buoyancy}), consecuencia directa del calentamiento desigual del aire. Conforme la radiación solar incrementa la temperatura de las superficies internas, el aire cercano a estas regiones reduce su densidad y asciende bajo la acción de la gravedad. Simultáneamente, aire exterior más frío y denso ingresa por las aberturas inferiores, estableciendo una circulación continua capaz de transportar calor y humedad hacia el exterior.

Desde una perspectiva matemática, este fenómeno corresponde a un problema multifísico acoplado, donde la distribución espacial de velocidad, presión, temperatura y concentración de vapor de agua evolucionan simultáneamente. La descripción rigurosa del sistema requiere resolver las ecuaciones de conservación de masa, cantidad de movimiento, energía y transporte de especies.

Para formular el modelo físico empleado en esta investigación se adoptaron las siguientes hipótesis:

\begin{itemize}
\item El aire se modela como un fluido newtoniano.
\item Las velocidades del flujo son bajas ($Ma<0.3$), por lo que los efectos de compresibilidad pueden despreciarse.
\item La gravedad permanece constante durante todo el dominio.
\item Las propiedades termofísicas del aire se consideran aproximadamente constantes, excepto la densidad en el término de flotación.
\item La variación de densidad debida a diferencias de temperatura se modela mediante la aproximación de Boussinesq.
\item El transporte de humedad se representa mediante una ecuación adicional de transporte escalar.
\end{itemize}

Bajo estas hipótesis, el movimiento del aire puede describirse mediante las ecuaciones de Navier--Stokes.

%%%%%%%%%%%%%%%%%%%%%%%%%%%%%%%%%%%%%%%%%%%%%%%%%%%%%%%
\subsection{Ecuaciones de Navier--Stokes}

Las ecuaciones de Navier--Stokes constituyen la formulación matemática fundamental para describir el movimiento de fluidos viscosos. Estas ecuaciones representan directamente la aplicación de los principios de conservación de masa, conservación de cantidad de movimiento y conservación de energía sobre un volumen diferencial de control.

\subsubsection{Conservación de la masa}

La conservación de masa establece que la masa no puede crearse ni destruirse dentro del dominio. En su forma diferencial,

\begin{equation}
\frac{\partial \rho}{\partial t}
+\nabla\cdot(\rho\mathbf{u})=0,
\label{eq:continuidad}
\end{equation}

donde

\begin{itemize}
\item $\rho$ es la densidad del fluido,
\item $\mathbf{u}$ es el vector velocidad.
\end{itemize}

Para flujos con pequeñas variaciones de densidad, como ocurre en el interior del secador, esta ecuación se reduce a

\begin{equation}
\nabla\cdot\mathbf{u}=0.
\end{equation}

Esta condición implica que el flujo de aire puede considerarse prácticamente incompresible.

\subsubsection{Conservación del momento}

Aplicando la segunda ley de Newton sobre un volumen diferencial de fluido se obtiene

\begin{equation}
\rho
\left(
\frac{\partial\mathbf{u}}{\partial t}
+
(\mathbf{u}\cdot\nabla)\mathbf{u}
\right)
=
-\nabla p
+\mu\nabla^2\mathbf{u}
+\rho\mathbf{g}
+\mathbf{S},
\label{eq:navierstokes}
\end{equation}

donde

\begin{itemize}

\item $p$ representa la presión,

\item $\mu$ la viscosidad dinámica,

\item $\mathbf{g}$ el vector gravitacional,

\item $\mathbf{S}$ engloba posibles términos fuente adicionales.
\end{itemize}

Cada término posee una interpretación física bien definida. El primer término representa la aceleración local, el segundo la aceleración convectiva, mientras que los términos del lado derecho describen respectivamente las fuerzas de presión, los efectos viscosos y las fuerzas volumétricas.

En aplicaciones de secadores solares el término gravitacional adquiere especial importancia debido a que constituye el mecanismo responsable de la circulación natural del aire.

\subsubsection{Conservación de la energía}

La evolución de la temperatura del fluido puede describirse mediante la ecuación de energía,

\begin{equation}
\rho c_p
\left(
\frac{\partial T}{\partial t}
+\mathbf{u}\cdot\nabla T
\right)
=
\nabla\cdot(k\nabla T)
+
S_r
+
S_h,
\label{eq:energia}
\end{equation}

donde

\begin{itemize}

\item $c_p$ corresponde al calor específico,

\item $k$ es la conductividad térmica,

\item $S_r$ representa la contribución de la radiación solar,

\item $S_h$ agrupa otras fuentes de calor.
\end{itemize}

La ecuación anterior establece el balance entre el transporte convectivo de energía, la difusión térmica y las fuentes internas de calor. En el caso del secador solar, la radiación incidente constituye la principal fuente energética del sistema.

\subsubsection{Transporte de humedad}

Debido a que el objetivo principal del secado consiste en remover agua del producto agrícola, resulta necesario incorporar una ecuación adicional para describir el transporte del vapor de agua dentro del dominio.

Una formulación ampliamente utilizada consiste en resolver la ecuación de transporte de especies,

\begin{equation}
\frac{\partial(\rho Y)}
{\partial t}
+
\nabla\cdot(\rho\mathbf{u}Y)
=
\nabla\cdot
(\rho D\nabla Y)
+
S_Y,
\label{eq:species}
\end{equation}

donde

\begin{itemize}

\item $Y$ representa la fracción másica de vapor,

\item $D$ el coeficiente de difusión,

\item $S_Y$ el término fuente asociado a la evaporación del agua desde el café.
\end{itemize}

Esta ecuación permite estudiar simultáneamente la evolución espacial de la humedad relativa dentro del secador, aspecto indispensable para evaluar el desempeño del proceso de secado. Diversos trabajos recientes sobre secadores solares e invernaderos resuelven explícitamente esta ecuación junto con las ecuaciones de Navier--Stokes para representar el transporte de humedad. 

%%%%%%%%%%%%%%%%%%%%%%%%%%%%%%%%%%%%%%%%%%%%%%%%%%%%%%%%%%%%%%%
\subsection{Turbulencia}

Las ecuaciones de Navier--Stokes presentadas en la sección anterior describen de manera exacta el movimiento de un fluido newtoniano. Sin embargo, su resolución directa únicamente resulta factible para flujos laminares o para dominios de tamaño reducido. En la mayoría de las aplicaciones de ingeniería, incluyendo los secadores solares de café, el flujo de aire presenta fluctuaciones espaciales y temporales que caracterizan un régimen turbulento.

La turbulencia corresponde a un fenómeno altamente tridimensional, caótico y disipativo, caracterizado por la aparición de remolinos (\textit{eddies}) de múltiples escalas que interactúan continuamente entre sí. Estas estructuras incrementan considerablemente el transporte de cantidad de movimiento, energía y especies químicas respecto a un flujo estrictamente laminar.

Desde el punto de vista físico, la turbulencia produce un mezclado mucho más eficiente del aire dentro del secador, modificando directamente la distribución espacial de temperatura y humedad. En consecuencia, la correcta representación de este fenómeno resulta indispensable para predecir con precisión el comportamiento térmico del sistema.

\subsubsection{Número de Reynolds}

La transición entre un régimen laminar y uno turbulento suele caracterizarse mediante el número de Reynolds,

\begin{equation}
Re=\frac{\rho UL}{\mu},
\end{equation}

donde

\begin{itemize}
\item $\rho$ representa la densidad del fluido,
\item $U$ una velocidad característica,
\item $L$ una longitud característica,
\item $\mu$ la viscosidad dinámica.
\end{itemize}

Este número adimensional representa la razón entre las fuerzas inerciales y las fuerzas viscosas presentes en el flujo.

Cuando predominan los efectos viscosos ($Re$ pequeño), las trayectorias del fluido permanecen ordenadas y el flujo se mantiene laminar. Conforme aumenta el número de Reynolds, las fuerzas inerciales dominan el movimiento y aparecen inestabilidades que eventualmente evolucionan hacia un régimen turbulento.

En secadores solares tipo invernadero, el régimen de flujo depende simultáneamente de la geometría del secador, la intensidad de la radiación solar, la ventilación, las diferencias de temperatura y las condiciones meteorológicas exteriores, por lo que pueden coexistir regiones laminares y regiones localmente turbulentas dentro del mismo dominio.

%%%%%%%%%%%%%%%%%%%%%%%%%%%%%%%%%%%%%%%%%%%%%%%%%%%%%%%%%%%%%%%
\subsubsection{Descomposición de Reynolds}

Resolver directamente todas las escalas turbulentas presentes en las ecuaciones de Navier--Stokes (DNS) resulta computacionalmente inviable para aplicaciones de ingeniería debido al enorme número de grados de libertad requeridos.

Una alternativa consiste en descomponer cada variable instantánea como la suma de una componente media y una componente fluctuante,

\begin{equation}
\phi=\overline{\phi}+\phi',
\end{equation}

de manera que, para la velocidad,

\begin{equation}
u_i=\overline{u_i}+u_i'.
\end{equation}

Esta formulación, conocida como descomposición de Reynolds, permite separar el movimiento promedio del fluido de las fluctuaciones turbulentas.

Al sustituir esta descomposición dentro de las ecuaciones de Navier--Stokes y efectuar un promedio temporal, se obtienen las ecuaciones conocidas como Reynolds Averaged Navier--Stokes (RANS),

\begin{equation}
\rho
\left(
\frac{\partial\overline{u_i}}{\partial t}
+
\overline{u_j}
\frac{\partial\overline{u_i}}{\partial x_j}
\right)
=
-\frac{\partial\overline{p}}{\partial x_i}
+
\mu
\frac{\partial^2\overline{u_i}}
{\partial x_j^2}
-
\frac{\partial}
{\partial x_j}
\left(
\rho
\overline{u_i'u_j'}
\right)
+
\rho g_i.
\end{equation}

El término

\[
-\rho\overline{u_i'u_j'}
\]

recibe el nombre de tensor de esfuerzos de Reynolds y representa el transporte adicional de cantidad de movimiento ocasionado por las fluctuaciones turbulentas.

%%%%%%%%%%%%%%%%%%%%%%%%%%%%%%%%%%%%%%%%%%%%%%%%%%%%%%%%%%%%%%%
\subsubsection{Problema de cierre}

La aparición del tensor de Reynolds introduce seis incógnitas adicionales dentro del sistema de ecuaciones, fenómeno conocido como el problema de cierre de la turbulencia.

Como consecuencia, las ecuaciones RANS dejan de ser un sistema cerrado y requieren incorporar relaciones constitutivas adicionales que permitan modelar dichos esfuerzos turbulentos.

La aproximación más ampliamente utilizada consiste en asumir la hipótesis de viscosidad turbulenta propuesta por Boussinesq,

\begin{equation}
-\rho\overline{u_i'u_j'}
=
\mu_t
\left(
\frac{\partial u_i}{\partial x_j}
+
\frac{\partial u_j}{\partial x_i}
\right)
-
\frac{2}{3}
\rho k\delta_{ij},
\end{equation}

donde

\begin{itemize}
\item $\mu_t$ corresponde a la viscosidad turbulenta,

\item $k$ representa la energía cinética turbulenta,

\item $\delta_{ij}$ es el delta de Kronecker.
\end{itemize}

Esta formulación constituye la base de prácticamente todos los modelos RANS empleados en dinámica de fluidos computacional.

%%%%%%%%%%%%%%%%%%%%%%%%%%%%%%%%%%%%%%%%%%%%%%%%%%%%%%%%%%%%%%%
\subsubsection{Modelos de turbulencia}

Diversos modelos han sido desarrollados para estimar la viscosidad turbulenta.

Entre los más utilizados se encuentran:

\begin{table}[h!]
\centering
\caption{Modelos de turbulencia empleados en CFD.}
\begin{tabular}{p{3.8cm}p{10cm}}
\toprule
Modelo & Características principales\\
\midrule

Spalart--Allmaras &
Modelo de una ecuación. Presenta bajo costo computacional y buen desempeño para flujos externos aerodinámicos. \\

$k-\varepsilon$ estándar &
Modelo clásico de dos ecuaciones ampliamente utilizado en aplicaciones industriales con flujos completamente desarrollados. \\

Realizable $k-\varepsilon$ &
Introduce mejoras en la formulación de la viscosidad turbulenta y suele ofrecer mejores predicciones para flujos con recirculaciones y convección natural. \\

$k-\omega$ SST &
Adecuado para regiones próximas a paredes y flujos con separación. \\

LES &
Resuelve directamente las estructuras turbulentas de gran escala, modelando únicamente las escalas pequeñas. \\

DNS &
Resuelve todas las escalas de la turbulencia sin utilizar modelos empíricos. Su costo computacional resulta prohibitivo para aplicaciones industriales. \\

\bottomrule
\end{tabular}
\end{table}

En investigaciones recientes sobre secadores solares e invernaderos agrícolas predominan los modelos RANS debido a que proporcionan un adecuado compromiso entre precisión y costo computacional. Diversos trabajos emplean variantes del modelo $k-\varepsilon$ junto con la aproximación de Boussinesq para representar la convección natural en cavidades cerradas. 

Finalmente, el modelo de turbulencia seleccionado para una simulación debe responder al fenómeno físico dominante y no únicamente a criterios computacionales. En secadores solares, donde la circulación del aire está gobernada principalmente por diferencias de temperatura y fuerzas de flotación, resulta indispensable que el modelo escogido reproduzca adecuadamente los efectos de la convección natural y las regiones de recirculación presentes en el interior del dominio.

%%%%%%%%%%%%%%%%%%%%%%%%%%%%%%%%%%%%%%%%%%%%%%%%%%%%%%%%%%%%%%%
\subsection{Convección natural}

En los secadores solares tipo invernadero, el movimiento del aire no es generado principalmente por dispositivos mecánicos, sino por diferencias de densidad inducidas por el calentamiento solar. Este fenómeno, conocido como \textit{convección natural} o \textit{convección libre}, constituye el mecanismo dominante responsable de la circulación del aire, el transporte de calor y la remoción de humedad durante el proceso de secado.

Cuando la radiación solar incide sobre la cubierta y las superficies internas del secador, estas absorben energía y aumentan su temperatura. El aire en contacto con dichas superficies recibe calor por convección, experimentando una disminución de densidad debido a la expansión térmica. Bajo la acción de la gravedad, este aire menos denso asciende mientras que aire más frío y denso desciende o ingresa desde el exterior, estableciendo una circulación convectiva continua.

Desde el punto de vista físico, este mecanismo representa una conversión de energía térmica en energía mecánica. El calentamiento produce diferencias de densidad, las cuales generan fuerzas de flotación que aceleran el fluido y originan el movimiento convectivo.

%%%%%%%%%%%%%%%%%%%%%%%%%%%%%%%%%%%%%%%%%%%%%%%%%%%%%%%%%%%%%%%
\subsubsection{Aproximación de Boussinesq}

Aunque las variaciones de densidad producidas por el calentamiento del aire son relativamente pequeñas, dichas variaciones resultan suficientes para generar el movimiento del fluido. En consecuencia, despreciarlas completamente impediría representar la convección natural.

Para resolver este inconveniente se emplea la aproximación de Boussinesq, ampliamente utilizada en simulaciones de convección natural debido a su bajo costo computacional y elevada precisión cuando las diferencias de temperatura son moderadas.

Bajo esta aproximación, todas las propiedades del aire se consideran constantes excepto la densidad en el término gravitacional de las ecuaciones de cantidad de movimiento. La densidad se aproxima mediante

\begin{equation}
\rho=\rho_0\left[1-\beta(T-T_0)\right],
\label{eq:boussinesq}
\end{equation}

donde

\begin{itemize}

\item $\rho_0$ representa la densidad de referencia,

\item $T_0$ corresponde a la temperatura de referencia,

\item $\beta$ es el coeficiente de expansión térmica del aire.
\end{itemize}

Al sustituir esta expresión en el término gravitacional de las ecuaciones de Navier--Stokes se obtiene

\begin{equation}
\rho\mathbf g
\approx
\rho_0\mathbf g
-
\rho_0\beta(T-T_0)\mathbf g.
\end{equation}

El primer término es constante y puede absorberse dentro del gradiente de presión, mientras que el segundo constituye la fuerza responsable del movimiento convectivo.

En consecuencia, la aceleración del aire deja de depender únicamente de gradientes de presión y pasa a depender directamente de la distribución de temperatura dentro del dominio.

Esta característica convierte el problema en un sistema fuertemente acoplado, ya que la temperatura modifica el campo de velocidades y, simultáneamente, dicho campo redistribuye nuevamente la temperatura.

%%%%%%%%%%%%%%%%%%%%%%%%%%%%%%%%%%%%%%%%%%%%%%%%%%%%%%%%%%%%%%%
\subsubsection{Números adimensionales}

El comportamiento de la convección natural puede analizarse mediante diversos números adimensionales que representan la relación entre los mecanismos físicos involucrados.

\paragraph{Número de Grashof}

El número de Grashof mide la importancia relativa de las fuerzas de flotación respecto a las fuerzas viscosas,

\begin{equation}
Gr=
\frac
{g\beta\Delta TL^3}
{\nu^2},
\end{equation}

donde

\begin{itemize}

\item $g$ es la aceleración gravitacional,

\item $\Delta T$ representa la diferencia característica de temperatura,

\item $L$ corresponde a una longitud característica,

\item $\nu$ es la viscosidad cinemática.
\end{itemize}

Valores elevados de $Gr$ indican que los efectos de flotación dominan sobre los efectos viscosos, favoreciendo una circulación convectiva más intensa.

%%%%%%%%%%%%%%%%%%%%%%%%%%%%%%%%%%%%%%%%%%%%%%%%%%%%%%%%%%%%%%%

\paragraph{Número de Prandtl}

El número de Prandtl relaciona la difusión de cantidad de movimiento con la difusión térmica,

\begin{equation}
Pr=
\frac{\nu}{\alpha},
\end{equation}

donde $\alpha$ representa la difusividad térmica del fluido.

Para el aire atmosférico bajo condiciones normales,

\[
Pr\approx0.71,
\]

lo que indica que el calor y el momento difunden a velocidades comparables.

%%%%%%%%%%%%%%%%%%%%%%%%%%%%%%%%%%%%%%%%%%%%%%%%%%%%%%%%%%%%%%%

\paragraph{Número de Rayleigh}

Uno de los parámetros más importantes en problemas de convección natural corresponde al número de Rayleigh,

\begin{equation}
Ra=Gr\cdot Pr
=
\frac
{g\beta\Delta TL^3}
{\nu\alpha}.
\end{equation}

Este parámetro representa simultáneamente los efectos de flotación, viscosidad y difusión térmica.

Su magnitud determina el régimen de circulación presente dentro de la cavidad.

De manera general,

\begin{center}
\begin{tabular}{cc}
\toprule
Número de Rayleigh & Comportamiento esperado\\
\midrule
$Ra<10^3$ & Conducción dominante\\
$10^3<Ra<10^6$ & Convección natural estable\\
$Ra>10^6$ & Convección intensa con aparición de inestabilidades\\
$Ra\gg10^8$ & Régimen altamente turbulento\\
\bottomrule
\end{tabular}
\end{center}

En secadores solares agrícolas, los valores de Rayleigh suelen ser suficientemente elevados para producir celdas convectivas y zonas de recirculación, especialmente durante periodos de máxima irradiancia solar.

%%%%%%%%%%%%%%%%%%%%%%%%%%%%%%%%%%%%%%%%%%%%%%%%%%%%%%%%%%%%%%%

\paragraph{Número de Nusselt}

El número de Nusselt cuantifica la intensidad del intercambio convectivo respecto al intercambio puramente conductivo,

\begin{equation}
Nu=
\frac{hL}{k},
\end{equation}

donde $h$ representa el coeficiente convectivo de transferencia de calor.

Cuando

\[
Nu=1,
\]

la transferencia ocurre exclusivamente por conducción.

Conforme aumenta $Nu$, la convección adquiere una mayor participación en el transporte de energía.

%%%%%%%%%%%%%%%%%%%%%%%%%%%%%%%%%%%%%%%%%%%%%%%%%%%%%%%%%%%%%%%

\subsubsection{Convección natural en cavidades cerradas}

El interior de un secador solar puede modelarse como una cavidad tridimensional calentada diferencialmente, uno de los problemas clásicos de la mecánica de fluidos computacional.

En este tipo de sistemas, la distribución espacial de temperatura genera regiones ascendentes de aire caliente próximas a las superficies irradiadas y regiones descendentes de aire relativamente frío cerca de las zonas menos calentadas. Como consecuencia aparecen celdas convectivas, recirculaciones locales y gradientes de velocidad que modifican significativamente la distribución térmica dentro del secador.

La geometría de la cavidad, la posición de las bandejas, las aberturas de ventilación y la intensidad de la radiación solar alteran directamente la estructura de dichas celdas convectivas, produciendo distribuciones no uniformes de temperatura y humedad sobre el producto agrícola.

Debido a este comportamiento altamente acoplado, pequeñas variaciones en la irradiancia solar o en las condiciones meteorológicas exteriores pueden modificar considerablemente la circulación interna del aire y, por consiguiente, la eficiencia del proceso de secado.

En consecuencia, el estudio de la convección natural constituye el fundamento físico sobre el cual se desarrollan los modelos de dinámica de fluidos computacional empleados en esta investigación, ya que permite representar simultáneamente el transporte de cantidad de movimiento, energía y vapor de agua dentro del secador solar.

\section{Modelos solares y radiación}

La radiación solar constituye la principal fuente de energía en los secadores solares tipo invernadero. A diferencia de los sistemas convencionales de secado, donde el calentamiento es producido mediante resistencias eléctricas o combustibles fósiles, en un secador solar la totalidad del proceso térmico depende de la energía electromagnética proveniente del Sol.

Desde el punto de vista físico, la radiación solar determina las condiciones de frontera energéticas del dominio computacional. La energía incidente sobre la cubierta transparente es parcialmente reflejada, parcialmente absorbida y parcialmente transmitida hacia el interior del secador, donde posteriormente es absorbida por las superficies internas, las bandejas y el producto agrícola.

Como consecuencia, la distribución espacial de irradiancia controla directamente la evolución temporal de la temperatura, el desarrollo de la convección natural y la velocidad de evaporación del agua contenida en el café.

Por esta razón, la representación adecuada de la radiación solar constituye uno de los componentes más importantes dentro del modelo físico desarrollado en esta investigación.

%%%%%%%%%%%%%%%%%%%%%%%%%%%%%%%%%%%%%%%%%%%%%%%%%%%%%%%%%%%%%%%
\subsection{Irradiancia global y directa}

La energía solar que alcanza una superficie terrestre no proviene exclusivamente del disco solar. Durante su propagación a través de la atmósfera, la radiación experimenta procesos de absorción, dispersión y reflexión producidos por moléculas de aire, vapor de agua, aerosoles y nubes. Como consecuencia, la irradiancia disponible sobre una superficie se compone de diferentes contribuciones.

Siguiendo la metodología propuesta por el \textit{ASHRAE Handbook}, la irradiancia solar incidente puede expresarse como la suma de tres componentes principales,

\begin{equation}
G=
G_b
+
G_d
+
G_r,
\label{eq:globalirradiance}
\end{equation}

donde

\begin{itemize}

\item $G_b$ corresponde a la irradiancia directa (\textit{beam radiation});

\item $G_d$ representa la irradiancia difusa proveniente del cielo;

\item $G_r$ corresponde a la radiación reflejada por el terreno y las superficies circundantes.
\end{itemize}

Esta descomposición constituye la base de numerosos modelos de carga solar utilizados en simulaciones energéticas y de dinámica de fluidos computacional. :contentReference[oaicite:0]{index=0}

%%%%%%%%%%%%%%%%%%%%%%%%%%%%%%%%%%%%%%%%%%%%%%%%%%%%%%%%%%%%%%%
\subsubsection{Radiación solar directa}

La radiación directa corresponde a la fracción de energía que llega desde el disco solar sin experimentar procesos significativos de dispersión atmosférica.

Su magnitud depende principalmente de:

\begin{itemize}

\item posición solar,

\item espesor óptico atmosférico,

\item nubosidad,

\item altitud,

\item hora del día,

\item época del año.
\end{itemize}

Cuando la superficie posee un vector normal $\mathbf n$, la potencia incidente puede aproximarse mediante la ley del coseno de Lambert,

\begin{equation}
G_{b,s}
=
G_b
\cos\theta,
\end{equation}

donde $\theta$ representa el ángulo formado entre la dirección de incidencia solar y la normal de la superficie.

Esta relación evidencia que la máxima energía incidente ocurre cuando la radiación es perpendicular a la superficie, mientras que disminuye progresivamente conforme aumenta el ángulo de incidencia.

%%%%%%%%%%%%%%%%%%%%%%%%%%%%%%%%%%%%%%%%%%%%%%%%%%%%%%%%%%%%%%%
\subsubsection{Radiación difusa}

La radiación difusa corresponde a la energía dispersada por la atmósfera antes de alcanzar la superficie terrestre.

A diferencia de la radiación directa, esta componente proviene prácticamente de toda la bóveda celeste y continúa aportando energía incluso durante condiciones parcialmente nubladas.

Aunque su intensidad suele ser inferior a la radiación directa durante días despejados, su contribución adquiere mayor importancia cuando aumenta la nubosidad o cuando el ángulo solar disminuye.

%%%%%%%%%%%%%%%%%%%%%%%%%%%%%%%%%%%%%%%%%%%%%%%%%%%%%%%%%%%%%%%
\subsubsection{Radiación reflejada}

Una fracción adicional de energía proviene de la reflexión de la radiación solar sobre el suelo y los objetos circundantes.

Su contribución puede estimarse mediante

\begin{equation}
G_r
=
\rho_g G_h
\left(
\frac{1-\cos\beta}{2}
\right),
\end{equation}

donde

\begin{itemize}

\item $\rho_g$ representa el albedo del terreno,

\item $G_h$ corresponde a la irradiancia horizontal global,

\item $\beta$ es la inclinación de la superficie.
\end{itemize}

Aunque esta componente suele representar una fracción relativamente pequeña de la irradiancia total, puede resultar relevante en superficies inclinadas o cuando existen materiales con elevada reflectancia en las inmediaciones del secador.

%%%%%%%%%%%%%%%%%%%%%%%%%%%%%%%%%%%%%%%%%%%%%%%%%%%%%%%%%%%%%%%
\subsubsection{Balance radiativo sobre la cubierta}

Cuando la radiación solar alcanza la cubierta transparente del secador, la energía incidente se distribuye entre tres mecanismos fundamentales:

\begin{equation}
G=\tau G+\alpha G+\rho G,
\end{equation}

o, de forma equivalente,

\begin{equation}
\tau+\alpha+\rho=1,
\label{eq:radiativebalance}
\end{equation}

donde

\begin{itemize}

\item $\tau$ representa la transmitancia,

\item $\alpha$ la absortancia,

\item $\rho$ la reflectancia.
\end{itemize}

La transmitancia determina la fracción de energía que logra ingresar al interior del secador, mientras que la absortancia controla el calentamiento de la cubierta y la reflectancia representa la energía devuelta hacia el ambiente.

En cubiertas fabricadas con policarbonato o películas plásticas agrícolas, la transmitancia suele constituir el parámetro de mayor interés, ya que condiciona directamente la energía disponible para el calentamiento del aire y del producto.

%%%%%%%%%%%%%%%%%%%%%%%%%%%%%%%%%%%%%%%%%%%%%%%%%%%%%%%%%%%%%%%
\subsubsection{Importancia dentro del secador solar}

La distribución espacial y temporal de la irradiancia constituye la condición de frontera energética que gobierna todo el comportamiento termo-fluidodinámico del secador.

En efecto,

\[
\boxed{
\text{Irradiancia}
\rightarrow
\text{Temperatura}
\rightarrow
\text{Convección}
\rightarrow
\text{Humedad}
\rightarrow
\text{Secado}
}
\]

Por consiguiente, pequeñas variaciones en la irradiancia solar ocasionadas por el paso de nubes, cambios estacionales o modificaciones en la posición solar pueden alterar significativamente el campo de temperaturas, la intensidad de la circulación convectiva y, finalmente, la eficiencia del proceso de secado.

Esta fuerte dependencia justifica la incorporación de modelos detallados de radiación solar dentro del modelo CFD desarrollado en la presente investigación, donde la irradiancia constituye la principal fuente energética del sistema y la condición de frontera más influyente sobre la evolución del flujo térmico. 


%%%%%%%%%%%%%%%%%%%%%%%%%%%%%%%%%%%%%%%%%%%%%%%%%%%%%%%%%%%%%%%
\subsection{Modelo atmosférico}

Antes de alcanzar la superficie terrestre, la radiación emitida por el Sol atraviesa la atmósfera, donde experimenta múltiples procesos de absorción, dispersión y reflexión. Como consecuencia, únicamente una fracción de la irradiancia extraterrestre alcanza finalmente el secador solar.

La atmósfera terrestre actúa como un medio radiativamente participante compuesto por moléculas de nitrógeno, oxígeno, vapor de agua, dióxido de carbono, ozono, aerosoles y partículas en suspensión. Cada uno de estos constituyentes modifica el espectro solar mediante mecanismos de absorción selectiva y dispersión, alterando tanto la intensidad como la dirección de propagación de la radiación incidente.

Desde un punto de vista energético, el balance radiativo puede representarse mediante

\begin{equation}
G_{0}
=
G_{atm}
+
G_{abs}
+
G_{sc},
\end{equation}

donde

\begin{itemize}
\item $G_0$ representa la irradiancia extraterrestre;
\item $G_{atm}$ corresponde a la irradiancia transmitida hasta la superficie;
\item $G_{abs}$ representa la energía absorbida por la atmósfera;
\item $G_{sc}$ corresponde a la energía dispersada hacia otras direcciones.
\end{itemize}

La irradiancia extraterrestre varía ligeramente durante el año debido a la excentricidad de la órbita terrestre y puede aproximarse mediante

\begin{equation}
G_{on}
=
G_{sc}
\left[
1
+
0.033
\cos
\left(
\frac{360n}{365}
\right)
\right],
\end{equation}

donde

\begin{itemize}
\item $G_{sc}=1367\;W/m^2$ es la constante solar,
\item $n$ representa el día juliano.
\end{itemize}

Aunque esta variación anual es relativamente pequeña (aproximadamente $\pm3.3\%$), resulta relevante en simulaciones energéticas de larga duración.

%%%%%%%%%%%%%%%%%%%%%%%%%%%%%%%%%%%%%%%%%%%%%%%%%%%%%%%%%%%%%%%

\subsubsection{Masa óptica atmosférica}

La cantidad de atmósfera atravesada por la radiación solar depende de la posición aparente del Sol respecto al observador.

Este efecto suele caracterizarse mediante la masa óptica atmosférica,

\begin{equation}
m
=
\frac{1}{\cos\theta_z},
\label{eq:airmass}
\end{equation}

donde $\theta_z$ representa el ángulo cenital solar.

Cuando el Sol se encuentra próximo al cenit ($\theta_z\approx0$), la radiación atraviesa una menor cantidad de atmósfera, produciendo pérdidas relativamente pequeñas por absorción y dispersión.

Por el contrario, durante el amanecer y el atardecer la trayectoria atmosférica aumenta considerablemente, incrementando la atenuación de la irradiancia directa.

Para ángulos cenitales elevados suele emplearse la corrección propuesta por Kasten y Young,

\begin{equation}
m=
\frac{1}
{\cos\theta_z
+
0.50572
(96.07995-\theta_z)^{-1.6364}},
\end{equation}

la cual proporciona una mejor aproximación cerca del horizonte.

%%%%%%%%%%%%%%%%%%%%%%%%%%%%%%%%%%%%%%%%%%%%%%%%%%%%%%%%%%%%%%%

\subsubsection{Modelos de cielo despejado}

Con el propósito de estimar la irradiancia disponible bajo condiciones atmosféricas ideales, diversos organismos internacionales han desarrollado modelos de cielo despejado (\textit{clear sky models}).

Entre los más utilizados se encuentran:

\begin{itemize}

\item ASHRAE Clear Sky Model.

\item Bird Clear Sky Model.

\item Hottel Model.

\item REST2.

\end{itemize}

En aplicaciones de ingeniería HVAC y simulación energética, el modelo ASHRAE constituye uno de los métodos más empleados debido a su simplicidad computacional y adecuada precisión para estimar la irradiancia directa normal y la irradiancia difusa.

Estos modelos representan la condición de frontera utilizada posteriormente por los algoritmos de radiación implementados en herramientas CFD.

%%%%%%%%%%%%%%%%%%%%%%%%%%%%%%%%%%%%%%%%%%%%%%%%%%%%%%%%%%%%%%%

\subsubsection{Geometría solar}

La energía recibida por una superficie depende directamente de la geometría Sol-Tierra.

\clearpage

Los principales parámetros utilizados para describir la posición aparente del Sol son:

\begin{table}[h!]
\centering
\caption{Variables fundamentales de la geometría solar.}
\begin{tabular}{ll}
\toprule
Variable & Descripción\\
\midrule

$\phi$ & Latitud\\

$\delta$ & Declinación solar\\

$\omega$ & Ángulo horario\\

$\theta_z$ & Ángulo cenital\\

$\gamma_s$ & Azimut solar\\

$\theta_i$ & Ángulo de incidencia sobre la superficie\\

\bottomrule
\end{tabular}
\end{table}



Estas variables determinan la orientación relativa entre el Sol y el secador, controlando directamente la irradiancia absorbida por la cubierta transparente.

En modelos numéricos modernos, la posición solar suele calcularse mediante algoritmos astronómicos como SOLPOS, ampliamente utilizados en simulaciones energéticas y de dinámica de fluidos computacional. 

\subsection{Acoplamiento entre radiación y convección natural}

La transferencia de calor dentro de un secador solar constituye un problema multifísico en el que la radiación térmica y la convección natural interactúan de manera simultánea. Aunque ambos mecanismos suelen analizarse de forma independiente, en realidad se encuentran fuertemente acoplados a través de la ecuación de energía. La radiación modifica la distribución espacial de temperaturas sobre las superficies del dominio, mientras que dichas distribuciones controlan las fuerzas de flotación responsables del movimiento convectivo del aire.

Desde el punto de vista del balance energético, el flujo total de calor puede expresarse como

\begin{equation}
q_{tot}=q_{conv}+q_{rad},
\label{eq:qtotal}
\end{equation}

donde

\begin{equation}
q_{conv}=h_c(T_s-T_f)
\end{equation}

representa la contribución convectiva y

\begin{equation}
q_{rad}=\varepsilon\sigma\left(T_s^4-T_{sur}^4\right)
\label{eq:qrad}
\end{equation}

corresponde al intercambio radiativo entre una superficie a temperatura $T_s$ y el entorno radiativo efectivo a temperatura $T_{sur}$.

Debido al carácter no lineal del término radiativo, pequeñas variaciones de temperatura pueden producir incrementos significativos en la energía intercambiada cuando las temperaturas superficiales son elevadas.

%%%%%%%%%%%%%%%%%%%%%%%%%%%%%%%%%%%%%%%%%%%%%%%%%%%%%%%%%%%%%%%

\subsubsection{Coeficiente radiativo equivalente}

Para facilitar la comparación entre la transferencia convectiva y radiativa resulta conveniente linealizar la ecuación de Stefan--Boltzmann alrededor de una temperatura media.

Aplicando una expansión de Taylor se obtiene

\begin{equation}
q_{rad}
=
h_r
(T_s-T_{sur}),
\end{equation}

donde

\begin{equation}
h_r=
\varepsilon\sigma
(T_s+T_{sur})
(T_s^2+T_{sur}^2)
\label{eq:hr}
\end{equation}

corresponde al coeficiente radiativo equivalente.

De forma análoga al coeficiente convectivo, este parámetro permite cuantificar la intensidad del intercambio radiativo mediante una formulación lineal equivalente.

%%%%%%%%%%%%%%%%%%%%%%%%%%%%%%%%%%%%%%%%%%%%%%%%%%%%%%%%%%%%%%%

\subsubsection{Número de Nusselt equivalente}

El número de Nusselt constituye una medida adimensional de la intensidad del transporte de calor respecto al mecanismo puramente conductivo.

Para la transferencia convectiva,

\begin{equation}
Nu_c=\frac{h_cL}{k},
\label{eq:nuc}
\end{equation}

mientras que utilizando el coeficiente radiativo equivalente puede definirse

\begin{equation}
Nu_r=\frac{h_rL}{k},
\label{eq:nur}
\end{equation}

donde $L$ representa una longitud característica y $k$ corresponde a la conductividad térmica del fluido.

En consecuencia, el número de Nusselt total puede interpretarse como

\begin{equation}
Nu_{tot}=Nu_c+Nu_r.
\label{eq:nutotal}
\end{equation}

Esta formulación permite evaluar cuantitativamente la contribución relativa de la radiación y la convección al intercambio térmico global dentro del secador.

%%%%%%%%%%%%%%%%%%%%%%%%%%%%%%%%%%%%%%%%%%%%%%%%%%%%%%%%%%%%%%%

\subsubsection{Interacción física entre radiación y convección}

A diferencia de lo que intuitivamente podría suponerse, un incremento de la transferencia radiativa no implica necesariamente un aumento de la intensidad de la convección natural.

La convección libre depende de las fuerzas de flotación generadas por gradientes de temperatura. Dichas fuerzas aparecen en las ecuaciones de cantidad de movimiento mediante la aproximación de Boussinesq,

\begin{equation}
\mathbf{F}_b
=
\rho_0\beta
(T-T_0)
\mathbf{g},
\end{equation}

por lo que la magnitud del movimiento convectivo depende directamente de la distribución espacial de temperaturas.

Cuando el intercambio radiativo entre superficies aumenta, la energía térmica se redistribuye de forma más uniforme dentro de la cavidad. Como consecuencia, disminuyen los gradientes locales de temperatura, reduciendo las diferencias de densidad que originan las fuerzas de flotación.

Desde una perspectiva matemática,

\begin{equation}
\nabla T
\downarrow
\Longrightarrow
Gr
\downarrow
\Longrightarrow
Ra
\downarrow
\Longrightarrow
Nu_c
\downarrow.
\end{equation}

Sin embargo, el incremento simultáneo del intercambio radiativo produce

\begin{equation}
Nu_r\uparrow,
\end{equation}

de manera que el número de Nusselt global puede continuar aumentando aun cuando la componente convectiva disminuya.

Este comportamiento ha sido observado tanto experimentalmente como mediante simulaciones numéricas de cavidades calentadas diferencialmente, donde la incorporación de modelos de radiación superficial reduce la intensidad de las células convectivas mientras incrementa la transferencia térmica total debido al mayor aporte radiativo. 

%%%%%%%%%%%%%%%%%%%%%%%%%%%%%%%%%%%%%%%%%%%%%%%%%%%%%%%%%%%%%%%

\subsubsection{Implicaciones para el modelado de secadores solares}

El fenómeno descrito posee especial importancia en secadores solares tipo invernadero, donde la radiación constituye la principal fuente energética y la circulación del aire depende casi exclusivamente de la convección natural.

En este contexto, la ecuación de energía, la aproximación de Boussinesq y las ecuaciones de Navier--Stokes forman un sistema fuertemente acoplado, en el que la radiación modifica el campo de temperaturas, éste altera las fuerzas de flotación y, finalmente, el movimiento del aire redistribuye nuevamente la temperatura y la humedad dentro del dominio.

En consecuencia, representar ambos mecanismos mediante modelos independientes puede conducir a estimaciones erróneas del flujo de calor, de la distribución de temperaturas y de la velocidad de secado. Por esta razón, los modelos CFD modernos empleados en cavidades térmicas incorporan explícitamente el acoplamiento entre la transferencia radiativa y la convección natural, resolviendo de manera simultánea las ecuaciones de conservación de masa, cantidad de movimiento, energía y radiación.


%%%%%%%%%%%%%%%%%%%%%%%%%%%%%%%%%%%%%%%%%%%%%%%%%%%%%%%%%%%%%%%%%%%%%%%%%%%%%%%%
\section{Métodos de modelado termo-fluidodinámico}

El comportamiento térmico de un secador solar constituye un problema multifísico en el que interactúan simultáneamente la conducción de calor, la convección natural, la radiación solar y el transporte de humedad. Estos fenómenos no ocurren de manera independiente, sino que se encuentran fuertemente acoplados mediante las ecuaciones de conservación de masa, cantidad de movimiento y energía.

La radiación solar incidente constituye la fuente primaria de energía del sistema. Una fracción de esta energía es absorbida por la cubierta y por las superficies internas del secador, elevando su temperatura. El calentamiento de dichas superficies genera gradientes térmicos que, a su vez, producen diferencias de densidad en el aire interior. Estas diferencias inducen fuerzas de flotación responsables de la circulación natural del aire dentro del secador, favoreciendo finalmente la evaporación de la humedad presente en los granos de café.

Desde el punto de vista matemático, este proceso involucra la resolución simultánea de ecuaciones diferenciales parciales acopladas, cuya solución analítica únicamente es posible bajo hipótesis altamente restrictivas. Por esta razón, la dinámica termo-fluidodinámica del sistema se aborda mediante Dinámica de Fluidos Computacional (Computational Fluid Dynamics, CFD), resolviendo numéricamente las ecuaciones gobernantes sobre un dominio discretizado.

En la presente investigación se adopta un enfoque de \textit{Conjugate Heat Transfer} (CHT), ampliamente utilizado para problemas donde coexisten regiones sólidas y fluidas. Bajo esta formulación, el dominio computacional se divide en regiones físicamente diferenciadas, resolviendo en cada una de ellas las ecuaciones gobernantes correspondientes e imponiendo condiciones de continuidad sobre las interfaces sólido--fluido.

%%%%%%%%%%%%%%%%%%%%%%%%%%%%%%%%%%%%%%%%%%%%%%%%%%%%%%%%%%%%%%%%%%%%%%%%%%%%%%%%
\subsection{Descomposición del dominio computacional}

El comportamiento físico del secador puede representarse mediante tres regiones principales, cada una caracterizada por mecanismos dominantes de transferencia de energía diferentes.

\begin{enumerate}

\item Región sólida.

\item Región fluida interna.

\item Región fluida externa.

\end{enumerate}

La región sólida comprende la estructura del secador, incluyendo la cubierta transparente, las paredes, el aislamiento térmico y los elementos estructurales. En esta región el mecanismo dominante de transporte de energía es la conducción de calor.

La región fluida interna corresponde al volumen de aire contenido dentro del secador. En ella predominan la convección natural, el transporte de energía y el transporte de vapor de agua producido durante el secado del café.

Finalmente, la región fluida externa representa la atmósfera circundante, donde la acción combinada del viento, la radiación solar y la convección exterior determina las condiciones de frontera impuestas sobre la estructura del secador.

Las tres regiones interactúan continuamente mediante intercambio de calor sobre sus interfaces comunes, dando origen a un problema de transferencia conjugada de calor cuyo esquema conceptual se muestra en la Figura XX.

%%%%%%%%%%%%%%%%%%%%%%%%%%%%%%%%%%%%%%%%%%%%%%%%%%%%%%%%%%%%%%%%%%%%%%%%%%%%%%%%
\subsection{Modelo de transferencia de calor en la región sólida}

Dentro de los materiales que conforman la estructura del secador, el mecanismo dominante de transporte energético corresponde a la conducción de calor. En consecuencia, la evolución temporal del campo de temperaturas se encuentra gobernada por la Ley de Fourier y la ecuación general de difusión del calor.

%%%%%%%%%%%%%%%%%%%%%%%%%%%%%%%%%%%%%%%%%%%%%%%%%%%%%%%%%%%%%%%%%%%%%%%%%%%%%%%%
\subsubsection{Ley de Fourier}

La Ley de Fourier establece que el flujo de calor es proporcional al gradiente negativo de temperatura,

\begin{equation}
\label{eqn:fourier}
\mathbf q=-k\nabla T,
\end{equation}

donde

\begin{itemize}

\item $\mathbf q$ representa el vector flujo de calor $\left(W/m^2\right)$;

\item $k$ corresponde a la conductividad térmica del material $\left(W/m\cdot K\right)$;

\item $T$ representa la temperatura.

\end{itemize}

En forma cartesiana,

\begin{equation}
\mathbf q=
-k
\left(
\hat{i}\frac{\partial T}{\partial x}
+
\hat{j}\frac{\partial T}{\partial y}
+
\hat{k}\frac{\partial T}{\partial z}
\right).
\end{equation}

El signo negativo indica que el calor siempre fluye desde regiones de mayor temperatura hacia regiones de menor temperatura, siguiendo la dirección opuesta al gradiente térmico.

Sobre una superficie cualquiera cuya normal unitaria es $\hat n$, la Ley de Fourier puede escribirse como

\begin{equation}
\label{eqn:fouriernormal}
q_n=
-k
\frac{\partial T}{\partial n},
\end{equation}

expresión ampliamente utilizada para formular las condiciones de frontera del problema.

Las ecuaciones anteriores son válidas bajo la hipótesis de materiales isotrópicos y homogéneos, donde la conductividad térmica puede considerarse constante en todas las direcciones.

%%%%%%%%%%%%%%%%%%%%%%%%%%%%%%%%%%%%%%%%%%%%%%%%%%%%%%%%%%%%%%%%%%%%%%%%%%%%%%%%
\subsubsection{Ecuación general de conducción}

Aplicando el principio de conservación de la energía sobre un volumen diferencial de control se obtiene la ecuación general de conducción de calor,

\begin{equation}
\label{eqn:heattransfer}
\rho c_p
\frac{\partial T}{\partial t}
=
\nabla\cdot(k\nabla T)
+
\dot q,
\end{equation}

donde

\begin{itemize}

\item $\rho$ es la densidad del material;

\item $c_p$ representa el calor específico;

\item $\dot q$ corresponde a la generación volumétrica de calor.

\end{itemize}

En la estructura del secador no existen fuentes internas significativas de generación térmica, por lo que puede suponerse

\[
\dot q=0.
\]

Si además se consideran propiedades térmicas constantes, la ecuación anterior se reduce a

\begin{equation}
\label{eqn:transferenciacalor}
\frac{\partial T}{\partial t}
=
\alpha
\nabla^2T,
\end{equation}

donde

\begin{equation}
\alpha=
\frac{k}{\rho c_p}
\end{equation}

corresponde a la difusividad térmica del material.

Esta ecuación gobierna la evolución temporal del campo de temperaturas dentro de todas las regiones sólidas del secador.

%%%%%%%%%%%%%%%%%%%%%%%%%%%%%%%%%%%%%%%%%%%%%%%%%%%%%%%%%%%%%%%%%%%%%%%%%%%%%%%%
\subsubsection{Difusividad térmica}

La difusividad térmica constituye una medida de la rapidez con la que una perturbación de temperatura se propaga dentro de un material.

Valores elevados de $\alpha$ indican que el calor se difunde rápidamente, produciendo una respuesta casi inmediata frente a cambios en las condiciones ambientales. Por el contrario, materiales con baja difusividad presentan una elevada capacidad de almacenamiento energético, amortiguando las variaciones térmicas externas.

Esta propiedad resulta particularmente importante en secadores solares, donde la irradiancia incidente varía continuamente debido al movimiento aparente del Sol y al paso de nubes. Un material con baja difusividad permite reducir las fluctuaciones rápidas de temperatura dentro del secador, favoreciendo un proceso de secado más uniforme.

%%%%%%%%%%%%%%%%%%%%%%%%%%%%%%%%%%%%%%%%%%%%%%%%%%%%%%%%%%%%%%%%%%%%%%%%%%%%%%%%
\subsubsection{Discretización espacial}

La resolución analítica de la ecuación (\ref{eqn:transferenciacalor}) únicamente resulta posible para geometrías y condiciones de frontera relativamente simples. En problemas reales, como el modelado de un secador solar, se requiere recurrir a métodos numéricos.

Considerando una discretización unidimensional mediante diferencias finitas explícitas,

\begin{equation}
T_i^{n+1}
=
T_i^n
+
\frac{\alpha\Delta t}{(\Delta x)^2}
\left(
T_{i+1}^n
-
2T_i^n
+
T_{i-1}^n
\right),
\end{equation}

donde

\begin{itemize}

\item $\Delta x$ representa el tamaño espacial de la malla;

\item $\Delta t$ corresponde al paso temporal.

\end{itemize}

Aunque la implementación desarrollada en esta investigación emplea el Método de Volúmenes Finitos (FVM) mediante OpenFOAM, la ecuación anterior ilustra claramente el proceso mediante el cual la ecuación diferencial continua puede transformarse en un sistema algebraico susceptible de resolverse numéricamente.

Asimismo, evidencia que la temperatura en cada nodo depende simultáneamente de las temperaturas de los nodos vecinos y del estado térmico correspondiente al instante temporal anterior, constituyendo la base de todos los métodos modernos de simulación numérica de transferencia de calor.

%%%%%%%%%%%%%%%%%%%%%%%%%%%%%%%%%%%%%%%%%%%%%%%%%%%%%%%%%%%%%%%%%%%%%%%%%%%%%%%%
\subsection{Condiciones de frontera y acoplamiento sólido--fluido}

La solución de la ecuación de conducción presentada en la sección anterior requiere la especificación de condiciones de frontera sobre todas las superficies que delimitan el dominio sólido.

En un secador solar, dichas superficies se encuentran sometidas simultáneamente a la incidencia de radiación solar, al intercambio de calor por convección con el aire circundante y al intercambio radiativo con el entorno. En consecuencia, las condiciones de frontera no pueden formularse únicamente en términos de temperatura prescrita (condiciones de Dirichlet) o flujo impuesto (condiciones de Neumann), sino mediante balances energéticos que consideran la acción conjunta de todos los mecanismos de transferencia de calor.

Desde el punto de vista físico, el flujo conducido hacia el interior del sólido debe ser igual al flujo neto intercambiado con el medio exterior,

\begin{equation}
q_{\mathrm{cond}}
=
q_{\mathrm{solar}}
-
q_{\mathrm{conv}}
-
q_{\mathrm{rad}}.
\label{eq:balance}
\end{equation}

Esta expresión constituye la condición de frontera conjugada empleada para acoplar el dominio sólido con los dominios fluidos.

%%%%%%%%%%%%%%%%%%%%%%%%%%%%%%%%%%%%%%%%%%%%%%%%%%%%%%%%%%%%%%%%%%%%%%%%%%%%%%%%
\subsubsection{Flujo de calor por conducción}

El flujo de calor conducido a través de la superficie viene dado por la Ley de Fourier,

\begin{equation}
q_{\mathrm{cond}}
=
-k
\frac{\partial T}{\partial n},
\label{eq:qcond}
\end{equation}

donde

\begin{itemize}

\item $k$ representa la conductividad térmica del material,

\item $\partial T/\partial n$ corresponde al gradiente normal de temperatura.

\end{itemize}

Este término representa la energía que abandona o ingresa al sólido como consecuencia de los gradientes térmicos existentes en su interior.

%%%%%%%%%%%%%%%%%%%%%%%%%%%%%%%%%%%%%%%%%%%%%%%%%%%%%%%%%%%%%%%%%%%%%%%%%%%%%%%%
\subsubsection{Radiación solar incidente}

La energía absorbida por la superficie del secador depende de la irradiancia incidente, de las propiedades ópticas del material y del ángulo de incidencia de los rayos solares.

El flujo absorbido puede escribirse como

\begin{equation}
q_{\mathrm{solar}}
=
\gamma
I_0
(\hat n_s\cdot\hat r_s),
\label{eq:qsolar}
\end{equation}

donde

\begin{itemize}

\item $\gamma$ representa la absortividad de la superficie;

\item $I_0$ corresponde a la irradiancia solar incidente;

\item $\hat n_s$ es el vector normal unitario a la superficie;

\item $\hat r_s$ representa la dirección de propagación de la radiación solar.

\end{itemize}

El producto escalar

\[
\hat n_s\cdot\hat r_s=\cos\theta
\]

introduce automáticamente la dependencia angular descrita por la Ley del Coseno de Lambert.

%%%%%%%%%%%%%%%%%%%%%%%%%%%%%%%%%%%%%%%%%%%%%%%%%%%%%%%%%%%%%%%%%%%%%%%%%%%%%%%%
\subsubsection{Transferencia de calor por convección}

El intercambio convectivo entre la superficie del secador y el aire circundante se modela mediante la Ley de Enfriamiento de Newton,

\begin{equation}
q_{\mathrm{conv}}
=
h
(T_s-T_\infty),
\label{eq:qconv}
\end{equation}

donde

\begin{itemize}

\item $h$ representa el coeficiente convectivo,

\item $T_s$ corresponde a la temperatura superficial,

\item $T_\infty$ representa la temperatura del fluido alejado de la superficie.

\end{itemize}

Este mecanismo depende tanto de las propiedades del fluido como del régimen de flujo presente alrededor del secador.

%%%%%%%%%%%%%%%%%%%%%%%%%%%%%%%%%%%%%%%%%%%%%%%%%%%%%%%%%%%%%%%%%%%%%%%%%%%%%%%%
\subsubsection{Intercambio radiativo}

Además de la radiación solar incidente, las superficies intercambian continuamente energía térmica mediante emisión y absorción de radiación infrarroja.

Su contribución se describe mediante la Ley de Stefan--Boltzmann,

\begin{equation}
q_{\mathrm{rad}}
=
\epsilon\sigma
(T_s^4-T_\infty^4),
\label{eq:qrad}
\end{equation}

donde

\begin{itemize}

\item $\epsilon$ representa la emisividad superficial;

\item $\sigma=5.670\times10^{-8}\;W/m^2K^4$ corresponde a la constante de Stefan--Boltzmann.

\end{itemize}

A diferencia de la convección, el intercambio radiativo depende de la cuarta potencia de la temperatura absoluta, lo que introduce una importante no linealidad en el problema.

%%%%%%%%%%%%%%%%%%%%%%%%%%%%%%%%%%%%%%%%%%%%%%%%%%%%%%%%%%%%%%%%%%%%%%%%%%%%%%%%
\subsubsection{Condición de frontera exterior}

Sustituyendo las ecuaciones
(\ref{eq:qcond})--(\ref{eq:qrad})
en el balance energético superficial
(\ref{eq:balance}),
se obtiene la condición de frontera aplicada sobre la superficie exterior del secador,

\begin{equation}
\boxed{
-k
\left(
\frac{\partial T}{\partial n}
\right)_{\mathrm{Ext}}
=
\gamma I_0
(\hat n_s\cdot\hat r_s)
-
h_s(T_s-T_\infty)
-
\epsilon\sigma
(T_s^4-T_\infty^4)
}
\label{eqn:neumannexterno}
\end{equation}

Esta ecuación expresa que el flujo conducido hacia el interior del sólido es consecuencia del balance entre la energía solar absorbida y las pérdidas de calor por convección y radiación hacia el ambiente.

%%%%%%%%%%%%%%%%%%%%%%%%%%%%%%%%%%%%%%%%%%%%%%%%%%%%%%%%%%%%%%%%%%%%%%%%%%%%%%%%
\subsubsection{Condición de frontera interior}

En la superficie interna del aislamiento desaparece el término asociado a la radiación solar directa. En este caso, el intercambio energético ocurre exclusivamente entre la pared y el aire contenido dentro del secador.

La condición de frontera correspondiente adopta la forma

\begin{equation}
\boxed{
-k
\left(
\frac{\partial T}{\partial n}
\right)_{\mathrm{Int}}
=
-
\epsilon\sigma
(T_{si}^4-T_i^4)
-
h(T_{si}-T_i)
}
\label{eqn:neumanninterno}
\end{equation}

donde

\begin{itemize}

\item $T_{si}$ representa la temperatura de la superficie interna del sólido;

\item $T_i$ corresponde a la temperatura del aire interno.

\end{itemize}

Esta condición constituye el mecanismo de acoplamiento entre la ecuación de conducción resuelta en la región sólida y las ecuaciones de Navier--Stokes resueltas en el dominio fluido.

%%%%%%%%%%%%%%%%%%%%%%%%%%%%%%%%%%%%%%%%%%%%%%%%%%%%%%%%%%%%%%%%%%%%%%%%%%%%%%%%
\subsubsection{Condiciones iniciales}

Para resolver el problema transitorio es necesario especificar la distribución inicial de temperaturas tanto en el sólido como en el aire.

En el presente trabajo se consideran las siguientes condiciones iniciales,

\begin{equation}
T(0,\mathrm{Ext})=T_\infty,
\label{eqn:inicialexterna}
\end{equation}

\begin{equation}
T(0,\mathrm{Int})=T_i.
\label{eqn:inicialinterna}
\end{equation}

Estas condiciones representan el estado térmico inicial del sistema antes del inicio del proceso de calentamiento por radiación solar.

%%%%%%%%%%%%%%%%%%%%%%%%%%%%%%%%%%%%%%%%%%%%%%%%%%%%%%%%%%%%%%%%%%%%%%%%%%%%%%%%
\subsubsection{Interpretación física}

Las ecuaciones anteriores muestran que la temperatura de la estructura del secador no depende exclusivamente de la irradiancia solar incidente, sino del equilibrio dinámico entre tres mecanismos de transferencia de energía.

Durante las horas de mayor insolación, el término asociado a la radiación solar incrementa la temperatura de las superficies externas. Parte de esta energía se propaga por conducción hacia el interior del material, mientras que otra fracción se pierde por convección hacia el ambiente y por emisión de radiación térmica.

Simultáneamente, el calor conducido a través del sólido es transferido al aire interno mediante convección natural y radiación superficial, elevando la temperatura del fluido y generando las diferencias de densidad responsables de la circulación convectiva dentro del secador.

En consecuencia, las condiciones de frontera constituyen el principal mecanismo de acoplamiento entre la conducción de calor en la estructura y la dinámica termo-fluidodinámica del aire, permitiendo representar de forma consistente la interacción entre los dominios sólido y fluido bajo el enfoque de transferencia conjugada de calor.


%%%%%%%%%%%%%%%%%%%%%%%%%%%%%%%%%%%%%%%%%%%%%%%%%%%%%%%%%%%%%%%%%%%%%%%%%%%%%%%%
\subsection{Estado estacionario y equilibrio térmico}

El comportamiento térmico de un secador solar puede analizarse tanto bajo condiciones transitorias como en régimen estacionario. Mientras el régimen transitorio describe la evolución temporal de la temperatura durante el calentamiento y enfriamiento del sistema, el estado estacionario representa una condición de equilibrio energético en la cual las variables macroscópicas permanecen constantes en el tiempo.

Aunque durante la operación real de un secador solar la irradiancia, la temperatura ambiente y la velocidad del viento varían continuamente, el análisis estacionario constituye una herramienta de gran utilidad para comprender la distribución espacial de temperaturas bajo condiciones ambientales aproximadamente constantes.

%%%%%%%%%%%%%%%%%%%%%%%%%%%%%%%%%%%%%%%%%%%%%%%%%%%%%%%%%%%%%%%%%%%%%%%%%%%%%%%%
\subsubsection{Ecuación gobernante en estado estacionario}

Partiendo de la ecuación general de conducción,

\begin{equation}
\rho c_p
\frac{\partial T}{\partial t}
=
\nabla\cdot(k\nabla T)
+
\dot q,
\end{equation}

el régimen estacionario implica

\begin{equation}
\frac{\partial T}{\partial t}=0.
\end{equation}

En ausencia de generación interna de calor,

\[
\dot q=0,
\]

la ecuación gobernante se reduce a

\begin{equation}
\nabla\cdot(k\nabla T)=0.
\end{equation}

Si además el material puede considerarse homogéneo e isotrópico, la conductividad térmica permanece constante y la ecuación adopta la conocida forma de Laplace,

\begin{equation}
\boxed{
\nabla^2T=0.
}
\label{eq:laplace}
\end{equation}

Esta ecuación establece que la temperatura en cualquier punto del sólido depende únicamente de las condiciones de frontera impuestas sobre el dominio.

%%%%%%%%%%%%%%%%%%%%%%%%%%%%%%%%%%%%%%%%%%%%%%%%%%%%%%%%%%%%%%%%%%%%%%%%%%%%%%%%
\subsubsection{Hipótesis de equilibrio térmico}

Con el propósito de estudiar el comportamiento estacionario del secador se consideran las siguientes hipótesis:

\begin{enumerate}

\item No existe acumulación de energía dentro del sistema.

\item Las propiedades termofísicas permanecen constantes.

\item No existe generación interna de calor.

\item La irradiancia solar permanece constante durante el intervalo analizado.

\item La temperatura ambiente es uniforme.

\item La velocidad del viento permanece aproximadamente constante.

\item La base del secador se considera adiabática.

\end{enumerate}

Estas hipótesis permiten aislar el efecto de la geometría y de las condiciones de frontera sobre la distribución espacial de temperaturas.

%%%%%%%%%%%%%%%%%%%%%%%%%%%%%%%%%%%%%%%%%%%%%%%%%%%%%%%%%%%%%%%%%%%%%%%%%%%%%%%%
\subsubsection{Balance integral de energía}

En régimen estacionario, el Primer Principio de la Termodinámica establece que la energía que ingresa al sistema debe ser exactamente igual a la energía que abandona el mismo.

Aplicando conservación de energía sobre toda la superficie del secador,

\begin{equation}
\oint_S
\mathbf q
\cdot
d\mathbf A
=
0.
\end{equation}

Esta expresión indica que el flujo neto de calor a través de toda la superficie cerrada es nulo.

Separando los aportes energéticos,

\begin{equation}
\oint_S
\left(
q_{solar}
-
q_{conv}
-
q_{rad}
\right)
\,dA
=
0.
\label{eq:globalbalance}
\end{equation}

Sustituyendo las expresiones obtenidas para cada mecanismo de transferencia,

\begin{align}
\oint_S
\Big[
&
\gamma I_0
(\hat n\cdot\hat r)
-
h(T_s-T_\infty)
\nonumber\\
&
-
\epsilon\sigma
(T_s^4-T_\infty^4)
\Big]
dA
=
0.
\label{eq:integralbalance}
\end{align}

La ecuación anterior representa el balance energético global del secador bajo condiciones estacionarias.

%%%%%%%%%%%%%%%%%%%%%%%%%%%%%%%%%%%%%%%%%%%%%%%%%%%%%%%%%%%%%%%%%%%%%%%%%%%%%%%%
\subsubsection{Forma discretizada}

Para resolver numéricamente el balance superficial, la superficie del secador se discretiza en un conjunto de elementos finitos de área $\Delta A_i$.

La integral superficial puede aproximarse mediante

\begin{equation}
\sum_{i=1}^{N}
\left[
\gamma I_0
\cos\theta_i
-
h_i(T_{s,i}-T_\infty)
-
\epsilon\sigma
(T_{s,i}^4-T_\infty^4)
\right]
\Delta A_i
=
0.
\end{equation}

Si la discretización superficial utiliza elementos aproximadamente uniformes,

\[
\Delta A_i
=
\Delta A,
\]

puede escribirse

\begin{equation}
\boxed{
\sum_{i=1}^{N}
\left[
\gamma I_0
\cos\theta_i
-
h_i(T_{s,i}-T_\infty)
-
\epsilon\sigma
(T_{s,i}^4-T_\infty^4)
\right]
=
0.
}
\end{equation}

Esta expresión constituye la forma algebraica utilizada para evaluar el equilibrio térmico sobre la superficie del secador.

%%%%%%%%%%%%%%%%%%%%%%%%%%%%%%%%%%%%%%%%%%%%%%%%%%%%%%%%%%%%%%%%%%%%%%%%%%%%%%%%
\subsubsection{Interpretación física}

El balance energético obtenido muestra que la distribución de temperaturas en estado estacionario no depende exclusivamente de la irradiancia incidente.

La temperatura superficial resulta del equilibrio dinámico entre tres mecanismos principales:

\begin{enumerate}

\item absorción de radiación solar;

\item pérdidas por convección;

\item pérdidas por radiación térmica.

\end{enumerate}

En consecuencia, diferentes regiones del secador pueden alcanzar temperaturas distintas aun cuando reciban la misma irradiancia solar. Dicho comportamiento depende de factores como la orientación de cada superficie respecto al Sol, el coeficiente local de convección, las propiedades ópticas del material y el intercambio radiativo con el entorno.

%%%%%%%%%%%%%%%%%%%%%%%%%%%%%%%%%%%%%%%%%%%%%%%%%%%%%%%%%%%%%%%%%%%%%%%%%%%%%%%%
\subsubsection{Limitaciones del análisis estacionario}

Aunque el estado estacionario proporciona información valiosa sobre la distribución espacial de temperaturas, la operación real de un secador solar constituye un fenómeno esencialmente transitorio.

La irradiancia solar varía continuamente debido al movimiento aparente del Sol, la presencia de nubosidad y los cambios atmosféricos. Simultáneamente, las temperaturas del aire, las velocidades del viento y la humedad relativa evolucionan durante el día.

Por esta razón, el análisis estacionario debe interpretarse como una aproximación al comportamiento promedio del sistema. La descripción completa del proceso requiere resolver las ecuaciones transitorias de conservación mediante simulaciones CFD dependientes del tiempo, las cuales permiten representar la evolución dinámica del secador bajo condiciones ambientales variables.

%%%%%%%%%%%%%%%%%%%%%%%%%%%%%%%%%%%%%%%%%%%%%%%%%%%%%%%%%%%%%%%%%%%%%%%%%%%%%%%%
\subsection{Modelo fluidodinámico del aire}

El comportamiento del aire, tanto en el interior como en el exterior del secador solar, se modela mediante las ecuaciones de Navier--Stokes para un fluido newtoniano incompresible. Estas ecuaciones describen simultáneamente la conservación de la masa, la cantidad de movimiento y la energía, permitiendo obtener los campos tridimensionales de velocidad, presión y temperatura responsables de la dinámica termo-fluidodinámica del sistema.

En el presente trabajo se distinguen dos dominios fluidos con características físicas diferentes:

\begin{enumerate}

\item El dominio atmosférico exterior, donde el flujo se encuentra dominado por el viento y la turbulencia atmosférica.

\item El dominio interior del secador, donde el movimiento del aire es generado principalmente por convección natural como consecuencia del calentamiento producido por la radiación solar.

\end{enumerate}

Aunque ambos dominios se describen mediante las mismas ecuaciones fundamentales, las hipótesis físicas y los modelos de turbulencia empleados difieren debido a la naturaleza del flujo presente en cada región.

%%%%%%%%%%%%%%%%%%%%%%%%%%%%%%%%%%%%%%%%%%%%%%%%%%%%%%%%%%%%%%%%%%%%%%%%%%%%%%%%
\subsubsection{Hipótesis del modelo}

Con el propósito de reducir la complejidad matemática sin comprometer la precisión de las simulaciones, se adoptan las siguientes hipótesis:

\begin{enumerate}

\item El aire se comporta como un fluido newtoniano.

\item El flujo es tridimensional.

\item El número de Mach es suficientemente pequeño para considerar flujo incompresible.

\item Las propiedades termofísicas permanecen aproximadamente constantes.

\item Los efectos gravitacionales se consideran mediante la aproximación de Boussinesq.

\item La transferencia de calor por radiación se incorpora como condición de frontera sobre las superficies del dominio.

\end{enumerate}

Bajo estas hipótesis, las ecuaciones gobernantes desarrolladas en el Capítulo XX describen adecuadamente la circulación del aire dentro del secador.

%%%%%%%%%%%%%%%%%%%%%%%%%%%%%%%%%%%%%%%%%%%%%%%%%%%%%%%%%%%%%%%%%%%%%%%%%%%%%%%%
\subsubsection{Modelo del aire exterior}

El dominio exterior representa el ambiente atmosférico que rodea al secador.

En esta región predominan fenómenos asociados al viento incidente, desprendimiento de capa límite y formación de estructuras turbulentas alrededor de la geometría del secador.

Para representar este comportamiento se resuelven las ecuaciones de Navier--Stokes acopladas con un modelo de turbulencia \textit{Spalart--Allmaras}, ampliamente utilizado en simulaciones aerodinámicas debido a su bajo costo computacional y buen desempeño para flujos externos.

El modelo exterior proporciona como resultado:

\begin{itemize}

\item distribución de velocidades;

\item distribución de presión;

\item coeficientes locales de convección;

\item intercambio térmico entre el secador y el ambiente.

\end{itemize}

Estos resultados constituyen las condiciones de frontera necesarias para el análisis térmico de la estructura.

%%%%%%%%%%%%%%%%%%%%%%%%%%%%%%%%%%%%%%%%%%%%%%%%%%%%%%%%%%%%%%%%%%%%%%%%%%%%%%%%
\subsubsection{Modelo del aire interior}

Dentro del secador el flujo posee una naturaleza completamente distinta.

La circulación del aire no es impuesta por ventiladores ni por diferencias de presión externas, sino que es consecuencia directa de las fuerzas de flotación originadas por diferencias de temperatura.

Estas fuerzas aparecen mediante la aproximación de Boussinesq,

\begin{equation}
\mathbf{F}_b
=
\rho_0
\beta
(T-T_0)
\mathbf g,
\end{equation}

donde

\begin{itemize}

\item $\beta$ representa el coeficiente de expansión térmica;

\item $T_0$ corresponde a la temperatura de referencia;

\item $\mathbf g$ representa el vector gravitacional.

\end{itemize}

Esta aproximación permite representar la variación de densidad únicamente dentro del término gravitacional, manteniendo constante la densidad en el resto de las ecuaciones de Navier--Stokes.

Como consecuencia, el aire caliente asciende hacia la parte superior del secador mientras que el aire relativamente frío desciende, originando células de recirculación responsables del transporte de calor y humedad.

%%%%%%%%%%%%%%%%%%%%%%%%%%%%%%%%%%%%%%%%%%%%%%%%%%%%%%%%%%%%%%%%%%%%%%%%%%%%%%%%
\subsubsection{Modelo de turbulencia}

Aunque las velocidades desarrolladas dentro del secador son relativamente bajas, la presencia de gradientes térmicos pronunciados, obstáculos internos y recirculaciones puede producir regímenes transicionales e incluso parcialmente turbulentos.

Con el propósito de representar adecuadamente dichos fenómenos se emplea el modelo Realizable $k-\varepsilon$, ampliamente utilizado para cavidades con convección natural debido a su adecuado compromiso entre precisión y costo computacional.

Este modelo resuelve dos ecuaciones adicionales correspondientes a la energía cinética turbulenta $k$ y a su tasa de disipación $\varepsilon$, permitiendo estimar la viscosidad turbulenta mediante

\begin{equation}
\mu_t
=
\rho
C_\mu
\frac{k^2}{\varepsilon}.
\end{equation}

La incorporación de este modelo mejora significativamente la predicción de los campos de velocidad y temperatura dentro del secador, especialmente en regiones donde existen zonas de recirculación y mezclado.

%%%%%%%%%%%%%%%%%%%%%%%%%%%%%%%%%%%%%%%%%%%%%%%%%%%%%%%%%%%%%%%%%%%%%%%%%%%%%%%%
\subsection{Acoplamiento sólido--fluido}

El problema físico descrito anteriormente corresponde a un caso clásico de transferencia conjugada de calor, donde la conducción térmica en los sólidos y la dinámica del aire deben resolverse simultáneamente.

El acoplamiento entre ambos dominios ocurre sobre cada una de las interfaces sólido--fluido mediante dos condiciones fundamentales.

La primera corresponde a la continuidad de temperatura,

\begin{equation}
T_s=T_f,
\end{equation}

mientras que la segunda impone la continuidad del flujo de calor,

\begin{equation}
-k_s
\frac{\partial T_s}{\partial n}
=
-k_f
\frac{\partial T_f}{\partial n}.
\end{equation}

Estas condiciones garantizan la conservación de la energía sobre la interfaz y permiten que el calor transferido por conducción a través de la estructura sea exactamente igual al calor recibido por el fluido mediante convección y radiación.

En consecuencia, la evolución térmica del secador no puede determinarse resolviendo independientemente las regiones sólida y fluida. Por el contrario, ambas deben resolverse simultáneamente mediante un esquema iterativo capaz de intercambiar información sobre temperatura y flujo de calor en cada paso temporal.

%%%%%%%%%%%%%%%%%%%%%%%%%%%%%%%%%%%%%%%%%%%%%%%%%%%%%%%%%%%%%%%%%%%%%%%%%%%%%%%%
\subsection{Modelo termo-fluidodinámico integrado}

El modelo desarrollado en esta investigación integra todos los mecanismos físicos descritos en las secciones anteriores dentro de un único sistema acoplado de ecuaciones diferenciales parciales.

La secuencia física del fenómeno puede resumirse como

\[
\boxed{
\text{Radiación Solar}\\

\longrightarrow
\text{Conducción}\\
\longrightarrow
\text{Convección Natural}\\
\longrightarrow
\text{Movimiento del Aire}\\
\longrightarrow
\text{Transporte de Humedad}\\
\longrightarrow
\text{Secado del Café}\\
}
\]

En consecuencia, la distribución final de temperaturas y velocidades observada dentro del secador constituye el resultado de la interacción simultánea entre la radiación solar, la conducción de calor en la estructura, la flotación térmica del aire y la transferencia conjugada de calor en las interfaces sólido--fluido.

Este sistema de ecuaciones constituye la base física sobre la cual posteriormente se implementa el modelo numérico mediante el Método de Volúmenes Finitos y el solucionador \texttt{chtMultiRegionFoam} de OpenFOAM, permitiendo reproducir la evolución termo-fluidodinámica del secador bajo diferentes condiciones ambientales.




%%%%%%%%%%%%%%%%%%%%%%%%%%%%%%%%%%%%%%%%%%%%%%%%%%%%%%%%%%%%%%%%%%%%%%%%%%%%%%%%
\subsection{Redes térmicas}

Aunque los modelos de Dinámica de Fluidos Computacional (CFD) permiten obtener una descripción espacial detallada de los campos de velocidad, temperatura y humedad dentro del secador solar, durante las primeras etapas de diseño resulta conveniente emplear modelos simplificados capaces de capturar los principales mecanismos físicos con un costo computacional significativamente menor.

Una de las aproximaciones más utilizadas corresponde a los modelos de parámetros concentrados (\textit{Lumped Parameter Models}), también conocidos como \textbf{redes térmicas}. En este enfoque, cada componente del sistema se representa mediante un nodo térmico caracterizado por una temperatura uniforme, mientras que los mecanismos de transferencia de calor entre dichos nodos son representados mediante resistencias térmicas equivalentes.

La filosofía de este método es completamente análoga al análisis de circuitos eléctricos. En una red eléctrica, una diferencia de potencial produce una corriente eléctrica limitada por una resistencia. De manera similar, en una red térmica una diferencia de temperatura genera un flujo de calor cuya magnitud depende de la resistencia térmica existente entre ambos nodos.

Esta analogía puede resumirse mediante las siguientes correspondencias:

\begin{table}[H]
\centering
\caption{Analogía entre circuitos eléctricos y redes térmicas.}
\begin{tabular}{ll}
\toprule
Circuito eléctrico & Red térmica\\
\midrule
Voltaje $V$ & Temperatura $T$\\
Corriente $I$ & Flujo de calor $\dot Q$\\
Resistencia eléctrica $R$ & Resistencia térmica $R_t$\\
Capacitancia $C$ & Capacidad térmica $\rho cV$\\
Ley de Ohm & Ley de Fourier/Newton\\
\bottomrule
\end{tabular}
\end{table}

Bajo esta formulación, la evolución térmica del secador puede interpretarse como un circuito dinámico en el cual la energía solar constituye la fuente de excitación del sistema, mientras que la conducción, la convección y la radiación representan los distintos mecanismos mediante los cuales dicha energía es distribuida y finalmente disipada.

%%%%%%%%%%%%%%%%%%%%%%%%%%%%%%%%%%%%%%%%%%%%%%%%%%%%%%%%%%%%%%%%%%%%%%%%%%%%%%%%
\subsubsection{Hipótesis del modelo}

Con el propósito de obtener un modelo matemático de complejidad moderada que conserve los fenómenos físicos dominantes, se adoptan las siguientes hipótesis:

\begin{table}[h!]
\centering
\caption{Hipótesis adoptadas para la formulación de la red térmica.}
\begin{tabular}{p{0.28\linewidth}p{0.62\linewidth}}
\toprule
Hipótesis & Justificación física\\
\midrule

Temperatura uniforme en cada componente
&
Cada elemento del secador se representa mediante un único nodo térmico.

\\[0.3cm]

Propiedades constantes
&
La conductividad térmica, densidad y calor específico permanecen aproximadamente constantes dentro del rango de temperaturas de operación.

\\[0.3cm]

Espesores uniformes
&
Las paredes, cubierta y placa absorbedora presentan espesor constante.

\\[0.3cm]

Radiación linealizada
&
Los intercambios radiativos se representan mediante coeficientes equivalentes de transferencia de calor.

\\[0.3cm]

Aire perfectamente mezclado
&
La temperatura y humedad del aire interior se consideran uniformes para cada instante.

\\[0.3cm]

Transferencia unidimensional
&
Los gradientes térmicos predominantes ocurren en la dirección normal a las superficies.

\\

\bottomrule
\end{tabular}
\end{table}

Estas hipótesis permiten obtener un sistema de ecuaciones diferenciales ordinarias acopladas que describen la evolución temporal de las temperaturas de los principales componentes del secador.

%%%%%%%%%%%%%%%%%%%%%%%%%%%%%%%%%%%%%%%%%%%%%%%%%%%%%%%%%%%%%%%%%%%%%%%%%%%%%%%%
\subsubsection{Representación nodal}

La red térmica desarrollada representa cada uno de los componentes principales del secador mediante un nodo de temperatura concentrada.

Los nodos considerados en el modelo son:

\begin{enumerate}

\item Cubierta transparente ($T_{cov}$).

\item Placa absorbedora ($T_{abs}$).

\item Producto agrícola ($T_{pad}$).

\item Aire interior ($T_{air}$).

\item Cielo efectivo ($T_{sky}$).

\item Ambiente exterior ($T_{amb}$).

\end{enumerate}

Cada uno de estos nodos intercambia energía con los demás mediante distintos mecanismos físicos. La cubierta absorbe una fracción de la radiación incidente y la intercambia con el ambiente, el aire interno y el cielo. La placa absorbedora recibe la radiación transmitida por la cubierta y la distribuye hacia el aire y el producto. Finalmente, el producto intercambia calor con el aire y consume parte de esta energía durante el proceso de evaporación del agua contenida en el café.

La Figura~\ref{fig:redtermica} presenta el esquema conceptual de la red térmica utilizada.

\begin{figure}[h!]
\centering

\fbox{
\begin{minipage}{0.9\linewidth}
\centering

\Large

Sol

$\downarrow$

Cubierta ($T_{cov}$)

$\downarrow$

Absorbedor ($T_{abs}$)

$\downarrow$

Producto ($T_{pad}$)

$\downarrow$

Aire ($T_{air}$)

\vspace{0.3cm}

Intercambios laterales:

Convección --- Radiación --- Conducción

\end{minipage}
}

\caption{Representación conceptual de la red térmica empleada para el secador solar. Posteriormente esta figura puede sustituirse por un diagrama elaborado en TikZ o Inkscape.}

\label{fig:redtermica}

\end{figure}

%%%%%%%%%%%%%%%%%%%%%%%%%%%%%%%%%%%%%%%%%%%%%%%%%%%%%%%%%%%%%%%%%%%%%%%%%%%%%%%%
\subsubsection{Balance energético general}

Todas las ecuaciones que gobiernan la red térmica se obtienen aplicando el Primer Principio de la Termodinámica sobre cada uno de los nodos.

En términos generales, el balance energético puede escribirse como

\begin{equation}
\boxed{
\text{Acumulación}
=
\text{Entradas}
-
\text{Salidas}
}
\end{equation}

Matemáticamente,

\begin{equation}
\boxed{
C_i
\frac{dT_i}{dt}
=
\sum \dot Q_{entrada}
-
\sum \dot Q_{salida}
}
\label{eq:generalbalance}
\end{equation}

donde

\[
C_i=\rho_i c_iV_i
\]

corresponde a la capacidad térmica equivalente del nodo $i$.

La ecuación (\ref{eq:generalbalance}) constituye la base de todo el modelo desarrollado en este trabajo. Cada componente del secador posee una capacidad térmica distinta y recibe o cede energía mediante mecanismos específicos de conducción, convección, radiación y evaporación.

En consecuencia, las ecuaciones particulares que se presentan en las siguientes subsecciones no representan modelos independientes, sino aplicaciones directas del balance energético general a cada uno de los nodos de la red térmica.

%%%%%%%%%%%%%%%%%%%%%%%%%%%%%%%%%%%%%%%%%%%%%%%%%%%%%%%%%%%%%%%%%%%%%%%%%%%%%%%%
\subsubsection{Variables de estado}

El sistema dinámico queda completamente descrito mediante las siguientes variables de estado:

\begin{table}
\centering
\caption{Variables de estado del modelo de red térmica.}
\begin{tabular}{lll}
\toprule
Variable & Descripción & Unidad\\
\midrule
$T_{cov}$ & Temperatura de la cubierta & K\\
$T_{abs}$ & Temperatura del absorbedor & K\\
$T_{pad}$ & Temperatura del producto & K\\
$T_{air}$ & Temperatura del aire interior & K\\
$H$ & Humedad específica del aire & kg/kg\\
$M$ & Contenido de humedad del producto & kg agua/kg materia seca\\
\bottomrule
\end{tabular}
\end{table}

Estas variables evolucionan simultáneamente durante el proceso de secado y se encuentran fuertemente acopladas. Por ejemplo, un incremento en la temperatura del absorbedor aumenta la temperatura del aire por convección, lo que acelera la evaporación del agua contenida en el producto. A su vez, la evaporación modifica tanto la humedad específica del aire como el calor latente requerido por el proceso, alterando nuevamente la distribución de temperaturas del sistema.

Por tanto, el modelo de red térmica constituye un sistema dinámico no lineal de ecuaciones diferenciales ordinarias acopladas cuya solución permite predecir la evolución temporal del proceso de secado antes de recurrir a simulaciones CFD de alta resolución.


%%%%%%%%%%%%%%%%%%%%%%%%%%%%%%%%%%%%%%%%%%%%%%%%%%%%%%%%%%%%%%%%%%%%%%%%%%%%%%%%
\subsubsection{Balance energético de los nodos térmicos}

Aplicando la ecuación general de conservación de la energía a cada uno de los nodos definidos en la red térmica, es posible obtener un sistema de ecuaciones diferenciales ordinarias acopladas que describen la evolución temporal del secador.

Cada nodo posee una capacidad térmica propia y recibe o cede energía mediante diferentes mecanismos de transferencia de calor. En consecuencia, aunque todas las ecuaciones poseen la misma estructura matemática,

\[
C_i\frac{dT_i}{dt}
=
\sum \dot Q_{entrada}
-
\sum \dot Q_{salida},
\]

los términos que aparecen en cada una dependen de la posición física del componente dentro del secador.

%%%%%%%%%%%%%%%%%%%%%%%%%%%%%%%%%%%%%%%%%%%%%%%%%%%%%%%%%%%%%%%%%%%%%%%%%%%%%%%%
\subsubsection{Nodo térmico de la cubierta transparente}

La cubierta transparente constituye el primer elemento que interactúa con la radiación solar. Su función principal consiste en permitir el paso de la mayor parte de la radiación de onda corta hacia el interior del secador mientras reduce las pérdidas de energía por convección y radiación.

Desde el punto de vista energético, la cubierta presenta cuatro mecanismos principales de intercambio de calor:

\begin{enumerate}

\item absorción directa de radiación solar;

\item convección con el ambiente exterior;

\item convección con el aire interno;

\item intercambio radiativo tanto con el cielo como con la placa absorbedora.

\end{enumerate}

El balance energético correspondiente puede visualizarse esquemáticamente como

\[
\boxed{
\text{Radiación solar}
\Longrightarrow
\boxed{\text{Cubierta}}
\Longrightarrow
\left\{
\begin{array}{c}
\text{Convección exterior}\\
\text{Convección interior}\\
\text{Radiación al cielo}\\
\text{Radiación al absorbedor}
\end{array}
\right.
}
\]

Aplicando el Primer Principio de la Termodinámica,

\[
\text{Acumulación}
=
\text{Radiación absorbida}
-
\text{Pérdidas convectivas}
-
\text{Pérdidas radiativas}.
\]

La capacidad térmica superficial de la cubierta viene dada por

\[
C_{cov}
=
\rho_{cov}\delta_{cov}c_{cov},
\]

de manera que el balance energético conduce a

\begin{equation}
\boxed{
\begin{aligned}
\rho_{cov}\delta_{cov}c_{cov}
\frac{\partial T_{cov}}{\partial t}
=
&
\alpha_{cov}(1+\tau_{cov}\rho_{abs})I_t
\\
&
-
\Big[
h_w(T_{cov}-T_{amb})
+
h_{c,cov-air}(T_{cov}-T_{air})
\\
&
+
h_{r,cov-sky}(T_{cov}-T_{sky})
+
h_{r,abs-cov}(T_{cov}-T_{abs})
\Big].
\end{aligned}
}
\end{equation}

\begin{tcolorbox}[title=Interpretación física,colback=blue!3,colframe=blue!60!black]

El primer término representa la energía solar absorbida por la cubierta.

Los cuatro términos restantes corresponden a mecanismos disipativos que reducen la temperatura de la cubierta:

\begin{itemize}

\item convección hacia el ambiente,

\item convección hacia el aire interno,

\item emisión radiativa hacia el cielo,

\item intercambio radiativo con la placa absorbedora.

\end{itemize}

La temperatura de la cubierta resulta del equilibrio dinámico entre estos cinco mecanismos.

\end{tcolorbox}

%%%%%%%%%%%%%%%%%%%%%%%%%%%%%%%%%%%%%%%%%%%%%%%%%%%%%%%%%%%%%%%%%%%%%%%%%%%%%%%%
\subsubsection{Nodo térmico de la placa absorbedora}

La placa absorbedora constituye el principal captador de energía del secador.

La mayor parte de la radiación transmitida por la cubierta es absorbida por esta superficie, transformándose posteriormente en energía térmica. Debido a su elevada absortividad, el absorbedor alcanza generalmente la mayor temperatura de todo el sistema, actuando como la principal fuente de calor para el aire y el producto.

Los mecanismos energéticos presentes son:

\begin{enumerate}

\item radiación solar transmitida por la cubierta;

\item convección hacia el aire interno;

\item radiación hacia la cubierta;

\item radiación hacia el cielo;

\item conducción hacia la estructura del secador.

\end{enumerate}

El esquema energético puede resumirse mediante

\[
\boxed{
\text{Radiación}
\Longrightarrow
\boxed{\text{Absorbedor}}
\Longrightarrow
\left\{
\begin{array}{c}
\text{Convección}\\
\text{Radiación}\\
\text{Conducción}
\end{array}
\right.
}
\]

Su capacidad térmica superficial está dada por

\[
C_{abs}
=
\rho_{abs}\delta_{abs}c_{abs}.
\]

Aplicando nuevamente conservación de energía,

\begin{equation}
\boxed{
\begin{aligned}
\rho_{abs}\delta_{abs}c_{abs}
\frac{\partial T_{abs}}{\partial t}
=
&
\frac{\alpha_{abs}\tau_{cov}}
{1-(1-\alpha_{abs})\rho_{abs}}
I_t
\\
&
-
\Big[
h_c(T_{abs}-T_{air})
+
h_{r,sky-abs}(T_{abs}-T_{sky})
\\
&
+
h_{r,abs-cov}(T_{abs}-T_{cov})
+
K_{abs-asp}(T_{abs}-T_{amb})
\Big].
\end{aligned}
}
\end{equation}

\begin{tcolorbox}[title=Observación,colback=green!3,colframe=green!50!black]

La placa absorbedora constituye el elemento encargado de transformar la energía radiante en energía térmica.

Por esta razón, pequeñas variaciones en la absortividad superficial ($\alpha_{abs}$) producen cambios significativos en la eficiencia global del secador.

\end{tcolorbox}

%%%%%%%%%%%%%%%%%%%%%%%%%%%%%%%%%%%%%%%%%%%%%%%%%%%%%%%%%%%%%%%%%%%%%%%%%%%%%%%%
\subsubsection{Nodo térmico del producto agrícola}

El producto constituye el nodo térmico más complejo de toda la red.

A diferencia de la cubierta y del absorbedor, el café no únicamente almacena energía sensible, sino que además consume una fracción importante del calor recibido durante el proceso de evaporación del agua contenida en su interior.

Por esta razón, la capacidad térmica efectiva del producto depende de su contenido instantáneo de humedad.

El balance energético incluye

\begin{enumerate}

\item radiación absorbida;

\item conducción desde el absorbedor;

\item convección con el aire;

\item intercambio radiativo;

\item consumo de calor latente debido a la evaporación.

\end{enumerate}

La capacidad térmica efectiva puede escribirse como

\[
C_{pad}
=
\rho_{pad}\delta_{pad}
(c_{pad}+c_{liq}M),
\]

mostrando explícitamente que la inercia térmica del producto disminuye conforme avanza el secado.

Aplicando el balance energético,

\begin{equation}
\boxed{
\begin{aligned}
&
\rho_{pad}\delta_{pad}
(c_{pad}+c_{liq}M)
\frac{\partial T_{pad}}{\partial t}
=
\\
&
\frac{\alpha_{pad}\tau_{cov}}
{1-(1-\alpha_{pad})\rho_{pad}}
I_t
\\
&
+
\rho_{pad}\delta_{pad}
\left(
L_{pad}
+
(c_{vap}-c_{liq})T_{pad}
\right)
\frac{\partial M}{\partial t}
\\
&
+
\Big[
h_{c,pad-air}(T_{air}-T_{pad})
+
h_{r,pad-sky}(T_{sky}-T_{pad})
\\
&
+
h_{r,pad-cov}(T_{cov}-T_{pad})
+
K_{pad-abs}(T_{abs}-T_{pad})
\Big].
\end{aligned}
}
\end{equation}

\begin{tcolorbox}[title=Aspecto clave del modelo,colback=yellow!6,colframe=orange!70!black]

Esta ecuación diferencia al producto del resto de los nodos térmicos.

Mientras la cubierta y el absorbedor únicamente almacenan y transfieren calor, el producto consume energía durante el cambio de fase del agua.

En consecuencia, una fracción importante de la energía suministrada por el absorbedor no incrementa la temperatura del café, sino que es utilizada como calor latente de evaporación.

Este mecanismo constituye el fundamento físico del proceso de secado.

\end{tcolorbox}




%%%%%%%%%%%%%%%%%%%%%%%%%%%%%%%%%%%%%%%%%%%%%%%%%%%%%%%%%%%%%%%%%%%%%%%%%%%%%%%%
\subsubsection{Modelo higrotérmico del aire}

Hasta este punto se han obtenido los balances energéticos correspondientes a los elementos sólidos que conforman el secador solar. Sin embargo, el secado del café no depende únicamente de la temperatura de estos componentes, sino también del comportamiento del aire que circula dentro de la cámara.

El aire constituye el medio encargado de transportar simultáneamente energía térmica y vapor de agua. Conforme atraviesa el lecho de café, recibe calor desde las superficies calientes y absorbe la humedad evaporada del producto. En consecuencia, la temperatura y la humedad específica del aire evolucionan de forma acoplada durante todo el proceso de secado.

Desde el punto de vista físico, el dominio fluido puede describirse mediante dos variables de estado:

\begin{itemize}

\item la temperatura del aire, $T_{air}$;

\item la humedad específica del aire, $H$.

\end{itemize}

La primera gobierna el transporte de energía, mientras que la segunda describe el transporte de masa asociado al vapor de agua.

%%%%%%%%%%%%%%%%%%%%%%%%%%%%%%%%%%%%%%%%%%%%%%%%%%%%%%%%%%%%%%%%%%%%%%%%%%%%%%%%
\subsubsection{Balance energético del aire bajo convección natural}

Cuando el secador opera sin ventiladores, el movimiento del aire es producido exclusivamente por las fuerzas de flotación generadas por diferencias de temperatura. Este fenómeno, conocido como convección natural, origina células de recirculación que transportan el calor desde la placa absorbedora hacia el producto.

Bajo la hipótesis de aire perfectamente mezclado, la temperatura del aire puede suponerse uniforme dentro del volumen de control.

El balance energético puede representarse mediante

\[
\boxed{
\text{Convección desde cubierta}
+
\text{Convección desde absorbedor}
+
\text{Convección desde producto}
=
\text{Acumulación de energía del aire}
}
\]

La capacidad térmica efectiva del aire húmedo viene dada por

\[
C_{air}
=
\rho_{air}\delta_{air}
(c_{air}+c_{vap}H),
\]

donde el segundo término representa el efecto del vapor de agua contenido en el aire.

Aplicando el Primer Principio de la Termodinámica se obtiene

\begin{equation}
\boxed{
\begin{aligned}
\rho_{air}\delta_{air}
(c_{air}+c_{vap}H)
\frac{\partial T_{air}}{\partial t}
=
&
h_{c,cov-air}(T_{cov}-T_{air})
\\
&
+
h_{c,abs-air}(T_{abs}-T_{air})
\\
&
+
h_{c,pad-air}(T_{pad}-T_{air}).
\end{aligned}
}
\end{equation}

\begin{tcolorbox}[title=Interpretación física,colback=blue!3,colframe=blue!60!black]

El aire no genera calor por sí mismo.

Su temperatura aumenta únicamente debido al intercambio convectivo con los demás componentes del secador.

En consecuencia, la evolución térmica del aire depende completamente de la temperatura alcanzada por la cubierta, el absorbedor y el producto.

\end{tcolorbox}

%%%%%%%%%%%%%%%%%%%%%%%%%%%%%%%%%%%%%%%%%%%%%%%%%%%%%%%%%%%%%%%%%%%%%%%%%%%%%%%%
\subsubsection{Balance de masa del vapor de agua}

El segundo fenómeno que gobierna la dinámica del aire corresponde al transporte de humedad.

A medida que el café pierde agua por evaporación, el vapor producido es incorporado al aire interior, incrementando progresivamente su humedad específica.

Aplicando conservación de masa al vapor contenido dentro del secador,

\begin{equation}
\boxed{
\rho_{air}\delta_{air}
\frac{\partial H}{\partial t}
=
-
\rho_{pad}
V_{pad}
\frac{\partial M}{\partial t}
}
\end{equation}

donde

\begin{itemize}

\item $H$ representa la humedad específica del aire;

\item $M$ corresponde al contenido de humedad del café.

\end{itemize}

El signo negativo indica que la pérdida de humedad del producto produce un incremento equivalente de humedad en el aire.

%%%%%%%%%%%%%%%%%%%%%%%%%%%%%%%%%%%%%%%%%%%%%%%%%%%%%%%%%%%%%%%%%%%%%%%%%%%%%%%%
\subsubsection{Modelo bajo convección forzada}

Cuando el secador incorpora ventilación mecánica, el transporte energético deja de estar gobernado exclusivamente por la flotación térmica.

En este caso aparece un término advectivo asociado a la velocidad promedio del aire.

El balance energético se transforma en

\begin{equation}
\boxed{
\begin{aligned}
\rho_{air}\delta_{air}
v_{air}
(c_{air}+c_{vap}H)
\frac{\partial T_{air}}{\partial x}
=
&
h_{c,cov-air}(T_{cov}-T_{air})
\\
&
+
h_{c,pad-air}(T_{pad}-T_{air}).
\end{aligned}
}
\end{equation}

De forma análoga, el balance de humedad adopta la forma

\begin{equation}
\boxed{
\rho_{air}\delta_{air}
v_{air}
\frac{\partial H}{\partial x}
=
-
\rho_{pad}
V_{pad}
\frac{\partial M}{\partial t}.
}
\end{equation}

Puede observarse que la velocidad del aire aparece multiplicando los gradientes espaciales de temperatura y humedad, representando el transporte convectivo producido por la ventilación.

%%%%%%%%%%%%%%%%%%%%%%%%%%%%%%%%%%%%%%%%%%%%%%%%%%%%%%%%%%%%%%%%%%%%%%%%%%%%%%%%
\subsubsection{Cinética de secado}

El proceso de evaporación del agua contenida en el café se modela mediante una cinética de primer orden,

\begin{equation}
\boxed{
\frac{\partial M}{\partial t}
=
-k(M-M_e)
}
\end{equation}

donde

\begin{itemize}

\item $M$ representa el contenido instantáneo de humedad;

\item $M_e$ corresponde al contenido de humedad de equilibrio;

\item $k$ es la constante cinética de secado.

\end{itemize}

Esta ecuación indica que la velocidad de secado es proporcional a la diferencia entre la humedad actual del producto y la humedad de equilibrio.

Cuando

\[
M=M_e,
\]

el producto alcanza el equilibrio higroscópico con el ambiente y el proceso de secado cesa naturalmente.

\begin{tcolorbox}[title=Interpretación física,colback=yellow!5,colframe=orange!70!black]

La ecuación anterior constituye el mecanismo de acoplamiento entre el modelo térmico y el modelo de transporte de humedad.

A medida que aumenta la temperatura del aire, disminuye la humedad de equilibrio del café, incrementando la velocidad de evaporación y acelerando el proceso de secado.

\end{tcolorbox}

%%%%%%%%%%%%%%%%%%%%%%%%%%%%%%%%%%%%%%%%%%%%%%%%%%%%%%%%%%%%%%%%%%%%%%%%%%%%%%%%
\subsubsection{Humedad de equilibrio}

La humedad de equilibrio depende de las condiciones psicrométricas del aire y puede estimarse mediante la correlación

\begin{equation}
M_e=
-\frac{1}{0.179}
\ln
\left[
\left(
-\frac{T_{air}-273.15+16.912}{277.091}
\right)
\ln(\psi)
\right].
\end{equation}

Esta expresión relaciona el estado higroscópico del producto con la humedad relativa del aire, permitiendo determinar el límite teórico al cual puede llegar el proceso de secado bajo unas condiciones ambientales determinadas.

%%%%%%%%%%%%%%%%%%%%%%%%%%%%%%%%%%%%%%%%%%%%%%%%%%%%%%%%%%%%%%%%%%%%%%%%%%%%%%%%
\subsubsection{Constante cinética de secado}

La constante de secado depende tanto del producto agrícola como de las condiciones de operación del secador.

Una aproximación preliminar utilizada en este trabajo corresponde a

\begin{equation}
\boxed{
k=
4758
\exp
\left(
-\frac{2987}{T_{air}}
\right)
}
\end{equation}

\begin{tcolorbox}[title=Importante,colback=red!2,colframe=red!70!black]

La correlación anterior constituye una aproximación basada en resultados experimentales reportados en la literatura.

Para aplicaciones de alta precisión, la constante cinética deberá calibrarse experimentalmente para el café panameño considerando la variedad, el grado de madurez, el contenido inicial de humedad y las condiciones específicas de operación del secador.

\end{tcolorbox}


%%%%%%%%%%%%%%%%%%%%%%%%%%%%%%%%%%%%%%%%%%%%%%%%%%%%%%%%%%%%%%%%%%%%%%%%%%%%%%%%
\subsubsection{Modelo radiativo y psicrométrico}

Hasta este punto se ha desarrollado el modelo energético correspondiente a los distintos nodos térmicos del secador. Sin embargo, para cerrar completamente el sistema de ecuaciones es necesario establecer modelos constitutivos capaces de describir los mecanismos de intercambio radiativo, las propiedades psicrométricas del aire y las condiciones atmosféricas que gobiernan el proceso de secado.

En esta sección se presentan las correlaciones utilizadas para calcular los coeficientes equivalentes de radiación, la temperatura efectiva del cielo, las propiedades higrométricas del aire húmedo y el calor latente asociado a la evaporación del agua contenida en el producto agrícola.

%%%%%%%%%%%%%%%%%%%%%%%%%%%%%%%%%%%%%%%%%%%%%%%%%%%%%%%%%%%%%%%%%%%%%%%%%%%%%%%%
\subsubsection{Linealización del intercambio radiativo}

El intercambio de energía mediante radiación constituye un fenómeno inherentemente no lineal debido a la dependencia con la cuarta potencia de la temperatura absoluta, descrita por la Ley de Stefan--Boltzmann,

\begin{equation}
q_{rad}
=
\varepsilon
\sigma
\left(
T_1^4-T_2^4
\right),
\end{equation}

donde

\begin{itemize}

\item $\sigma$ es la constante de Stefan--Boltzmann,

\item $\varepsilon$ representa la emisividad superficial.

\end{itemize}

Aunque esta expresión describe rigurosamente el intercambio radiativo, su incorporación directa en modelos de parámetros concentrados dificulta considerablemente la resolución numérica.

Por esta razón, resulta conveniente linealizar el flujo radiativo definiendo un coeficiente equivalente de transferencia de calor,

\begin{equation}
q_{rad}
=
h_r
(T_1-T_2),
\end{equation}

donde

\begin{equation}
h_r
=
\varepsilon\sigma
(T_1^2+T_2^2)
(T_1+T_2).
\end{equation}

Esta aproximación permite tratar el intercambio radiativo de forma análoga a la convección, simplificando el acoplamiento entre todos los nodos térmicos.

%%%%%%%%%%%%%%%%%%%%%%%%%%%%%%%%%%%%%%%%%%%%%%%%%%%%%%%%%%%%%%%%%%%%%%%%%%%%%%%%
\subsubsection{Coeficientes radiativos equivalentes}

Considerando la geometría del secador y las propiedades ópticas de las distintas superficies, los coeficientes radiativos utilizados en el presente modelo se obtienen mediante las siguientes expresiones.

\paragraph{Radiación entre el absorbedor y la cubierta}

\begin{equation}
h_{r,abs-cov}
=
\frac{
\sigma
(T_{abs}^2+T_{cov}^2)
(T_{abs}+T_{cov})
}
{
\dfrac1{\epsilon_{abs}}
+
\dfrac1{\tau_{abs}}
+
\dfrac1{1-\tau_{cov}}
-
2
}.
\end{equation}

\paragraph{Radiación entre el producto y la cubierta}

\begin{equation}
h_{r,pad-cov}
=
\frac{
\sigma
(T_{pad}^2+T_{cov}^2)
(T_{pad}+T_{cov})
}
{
\dfrac1{\epsilon_{pad}}
+
\dfrac1{\tau_{abs}}
+
\dfrac1{1-\tau_{cov}}
-
2
}.
\end{equation}

\paragraph{Radiación entre la cubierta y el cielo}

\begin{equation}
h_{r,cov-sky}
=
\frac{
\sigma
(T_{cov}^2+T_{sky}^2)
(T_{cov}+T_{sky})
}
{
\dfrac1{\epsilon_{cov}}
+
\dfrac1{\tau_{cov}}
-
1
}.
\end{equation}

En todos los casos anteriores, los coeficientes representan una linealización local del intercambio radiativo entre superficies grises difusas.

\begin{tcolorbox}[title=Observación,colback=green!5,colframe=green!60!black]

Aunque los coeficientes radiativos poseen las mismas unidades que los coeficientes convectivos ($W/m^2K$), ambos representan mecanismos físicos completamente diferentes.

Mientras los coeficientes convectivos dependen del movimiento del fluido, los coeficientes radiativos dependen exclusivamente de las propiedades ópticas de las superficies y de sus temperaturas absolutas.

\end{tcolorbox}

%%%%%%%%%%%%%%%%%%%%%%%%%%%%%%%%%%%%%%%%%%%%%%%%%%%%%%%%%%%%%%%%%%%%%%%%%%%%%%%%
\subsubsection{Temperatura efectiva del cielo}

Durante la noche e incluso en condiciones diurnas despejadas, las superficies del secador intercambian energía mediante radiación con la bóveda celeste.

Este intercambio puede representarse mediante una temperatura radiativa equivalente denominada temperatura efectiva del cielo ($T_{sky}$), la cual generalmente es inferior a la temperatura ambiente.

En este trabajo se adopta la correlación

\begin{equation}
T_{sky}
=
T_{air}
\left[
0.8+
\frac{T_{dp}-273}{250}
\right]^{1/4}.
\end{equation}

Esta expresión incorpora indirectamente el efecto de la humedad atmosférica mediante la temperatura del punto de rocío.

%%%%%%%%%%%%%%%%%%%%%%%%%%%%%%%%%%%%%%%%%%%%%%%%%%%%%%%%%%%%%%%%%%%%%%%%%%%%%%%%
\subsubsection{Temperatura del punto de rocío}

La temperatura del punto de rocío se estima utilizando la correlación propuesta por ASHRAE,

\begin{equation}
T_{dp}
=
6.54
+
14.526\xi
+
0.7389\xi^2
+
0.09486\xi^3
+
0.4569(\phi p_s)^{0.1984},
\end{equation}

donde

\[
\xi
=
\ln(\phi p_s).
\]

Esta temperatura representa el valor al cual el vapor de agua comienza a condensarse para una presión dada.

%%%%%%%%%%%%%%%%%%%%%%%%%%%%%%%%%%%%%%%%%%%%%%%%%%%%%%%%%%%%%%%%%%%%%%%%%%%%%%%%
\subsubsection{Presión de saturación del vapor}

La presión de saturación necesaria para el cálculo psicrométrico se obtiene mediante

\begin{equation}
p_s
=
\frac
{6\times10^{25}}
{(273.15+T_{air})^5}
\exp
\left(
-\frac{6800}{273.15+T_{air}}
\right).
\end{equation}

Esta propiedad depende únicamente de la temperatura del aire y constituye una de las variables fundamentales en el análisis psicrométrico.

%%%%%%%%%%%%%%%%%%%%%%%%%%%%%%%%%%%%%%%%%%%%%%%%%%%%%%%%%%%%%%%%%%%%%%%%%%%%%%%%
\subsubsection{Humedad específica del aire}

La humedad específica representa la masa de vapor de agua contenida por unidad de masa de aire seco.

Se calcula mediante

\begin{equation}
H
=
\frac
{0.622\phi p_s}
{p_{atm}-\phi p_s}.
\end{equation}

Esta variable constituye el principal indicador del estado higrométrico del aire y gobierna la capacidad de éste para continuar absorbiendo humedad proveniente del producto.

%%%%%%%%%%%%%%%%%%%%%%%%%%%%%%%%%%%%%%%%%%%%%%%%%%%%%%%%%%%%%%%%%%%%%%%%%%%%%%%%
\subsubsection{Calor latente de evaporación}

Durante el proceso de secado una fracción importante de la energía absorbida por el café no incrementa su temperatura sino que es utilizada para evaporar el agua contenida en su interior.

El calor latente efectivo empleado en el modelo viene dado por

\begin{equation}
L_{pad}
=
2501.61
\left(
1+23e^{-0.4M}
\right).
\end{equation}

Puede observarse que el calor latente depende del contenido instantáneo de humedad del producto, permitiendo representar la evolución energética del proceso de secado conforme disminuye la cantidad de agua presente en el café.

%%%%%%%%%%%%%%%%%%%%%%%%%%%%%%%%%%%%%%%%%%%%%%%%%%%%%%%%%%%%%%%%%%%%%%%%%%%%%%%%
\subsubsection{Acoplamiento completo del modelo}

Las correlaciones desarrolladas en esta sección permiten cerrar completamente el sistema de ecuaciones diferenciales presentado anteriormente.

En conjunto, el modelo queda constituido por:

\begin{itemize}

\item Tres balances energéticos correspondientes a la cubierta, la placa absorbedora y el producto.

\item Un balance energético del aire.

\item Un balance de masa del vapor de agua.

\item Una ecuación cinética para el contenido de humedad del producto.

\item Correlaciones psicrométricas para el cálculo de la humedad específica.

\item Modelos radiativos para el intercambio energético entre superficies.

\end{itemize}

Todos estos modelos se encuentran acoplados mediante las variables de estado

\[
T_{cov},
\quad
T_{abs},
\quad
T_{pad},
\quad
T_{air},
\quad
H,
\quad
M,
\]

formando un sistema dinámico no lineal cuya solución permite predecir simultáneamente la evolución térmica e higrométrica del secador solar.

\begin{tcolorbox}[title=Conclusión del modelo,colback=blue!4,colframe=blue!70!black]

El modelo de red térmica desarrollado constituye una representación de parámetros concentrados del secador solar. Aunque simplifica la distribución espacial de temperaturas y velocidades, conserva los principales mecanismos físicos responsables del proceso de secado: absorción de radiación solar, conducción, convección, intercambio radiativo y evaporación del agua contenida en el producto.

Este modelo proporciona una herramienta de bajo costo computacional para el análisis preliminar, la validación conceptual y la evaluación de estrategias de diseño, además de servir como referencia para contrastar los resultados obtenidos posteriormente mediante simulaciones CFD tridimensionales.

\end{tcolorbox}


%\subsection{Modelos reducidos en 2D}

%%%%%%%%%%%%%%%%%%%%%%%%%%%%%%%%%%%%%%%%%%%%%%%%%%%%%%%%%%%%%%%%%%%%%%%%%%%%%%%%
\section{Machine Learning aplicado a procesos térmicos}
\label{sec:ml_procesos_termicos}

Los procesos térmicos que tienen lugar dentro de un secador solar se encuentran gobernados por la interacción entre radiación, conducción, convección, circulación de aire y transferencia de humedad. Estos mecanismos producen una respuesta dinámica, no lineal y dependiente de condiciones ambientales que varían continuamente durante el día.

El comportamiento del secador puede describirse mediante modelos físicos, como las redes térmicas y las ecuaciones de conservación implementadas mediante Dinámica de Fluidos Computacional. Sin embargo, la aplicación de estos modelos puede presentar limitaciones cuando se requiere realizar predicciones rápidas, analizar grandes cantidades de datos experimentales o evaluar múltiples configuraciones de operación.

En este contexto, el aprendizaje automático constituye una herramienta complementaria para identificar relaciones entre las condiciones ambientales, las variables de operación y la respuesta térmica del sistema. Su aplicación permite construir modelos predictivos capaces de estimar temperaturas, humedad, tasas de secado o indicadores de desempeño a partir de observaciones experimentales o resultados de simulación.

De manera general, un modelo basado en datos puede expresarse como

\begin{equation}
\hat{\mathbf{y}}
=
f_{\boldsymbol{\theta}}
\left(
\mathbf{x}
\right),
\label{eq:ml_general}
\end{equation}

donde $\mathbf{x}$ representa el conjunto de variables de entrada, $\hat{\mathbf{y}}$ contiene las variables predichas y $\boldsymbol{\theta}$ corresponde a los parámetros aprendidos durante el entrenamiento.

Para el secador solar, el vector de entrada puede incluir variables meteorológicas, térmicas y operacionales:

\begin{equation}
\mathbf{x}
=
\begin{bmatrix}
I_t &
T_{amb} &
\phi_{amb} &
v_w &
T_{air} &
T_{abs} &
u_f
\end{bmatrix}^{T},
\end{equation}

donde $I_t$ es la irradiancia solar, $T_{amb}$ la temperatura ambiente, $\phi_{amb}$ la humedad relativa exterior, $v_w$ la velocidad del viento y $u_f$ una variable asociada al funcionamiento del sistema de ventilación.

Las variables de salida pueden estar constituidas por

\begin{equation}
\mathbf{y}
=
\begin{bmatrix}
T_{pad} &
T_{air,int} &
\phi_{int} &
M &
\dot{M}
\end{bmatrix}^{T},
\end{equation}

donde $T_{pad}$ representa la temperatura del producto, $\phi_{int}$ la humedad relativa interior, $M$ el contenido de humedad y $\dot{M}$ la tasa instantánea de secado.

En esta investigación se consideran principalmente dos familias de métodos. La primera corresponde a los árboles de decisión y los algoritmos de \textit{boosting}, los cuales permiten construir modelos predictivos no lineales a partir de datos tabulares. La segunda comprende las redes neuronales informadas por la física, conocidas como \textit{Physics-Informed Neural Networks}, que incorporan las ecuaciones diferenciales gobernantes dentro del proceso de entrenamiento.

\begin{tcolorbox}[
title=Enfoque adoptado,
colback=blue!4,
colframe=blue!65!black
]
El aprendizaje automático no se plantea como sustituto de la termodinámica ni de la Dinámica de Fluidos Computacional. Su función es complementar estos modelos mediante tres capacidades principales:

\begin{enumerate}
    \item estimar variables térmicas a partir de mediciones experimentales;
    \item construir modelos sustitutos de simulaciones numéricas;
    \item integrar datos y ecuaciones físicas dentro de una formulación híbrida.
\end{enumerate}
\end{tcolorbox}

La selección de los métodos debe responder a la naturaleza del problema y a la disponibilidad de datos. Los métodos basados en árboles resultan especialmente adecuados para bases de datos experimentales estructuradas, mientras que las PINNs son apropiadas cuando se dispone de una formulación matemática conocida que puede utilizarse como restricción física durante el entrenamiento.



%%%%%%%%%%%%%%%%%%%%%%%%%%%%%%%%%%%%%%%%%%%%%%%%%%%%%%%%%%%%%%%%%%%%%%%%%%%%%%%%
\subsection{Árboles de decisión y métodos de boosting}
\label{subsec:arboles_boosting}

Los árboles de decisión constituyen una familia de algoritmos de aprendizaje supervisado capaces de representar relaciones no lineales mediante una sucesión de reglas jerárquicas. A diferencia de los modelos de regresión lineal, no requieren establecer previamente una forma funcional entre las variables de entrada y la variable objetivo. En su lugar, dividen recursivamente el espacio de características en regiones cada vez más homogéneas.

Esta propiedad resulta particularmente conveniente para el modelado de procesos térmicos, cuya respuesta depende de interacciones no lineales entre irradiancia solar, temperatura ambiente, humedad relativa, velocidad del aire, propiedades del producto y condiciones de operación.

Para un problema de regresión, el conjunto de datos puede expresarse como

\begin{equation}
\mathcal{D}
=
\left\{
\left(
\mathbf{x}_i,y_i
\right)
\right\}_{i=1}^{N},
\end{equation}

donde $\mathbf{x}_i\in\mathbb{R}^{p}$ representa el vector de variables de entrada de la observación $i$, mientras que $y_i$ corresponde a la variable térmica o higrométrica que se desea predecir.

En el caso del secador solar, un posible vector de entrada está dado por

\begin{equation}
\mathbf{x}_i
=
\begin{bmatrix}
I_t &
T_{amb} &
\phi_{amb} &
v_w &
T_{air} &
T_{abs} &
u_f &
t
\end{bmatrix}^{T},
\end{equation}

donde $I_t$ es la irradiancia solar, $T_{amb}$ la temperatura ambiente, $\phi_{amb}$ la humedad relativa exterior, $v_w$ la velocidad del viento, $T_{air}$ la temperatura del aire interior, $T_{abs}$ la temperatura del absorbedor, $u_f$ la condición de operación del sistema de ventilación y $t$ el tiempo transcurrido.

La variable objetivo puede corresponder, por ejemplo, a la temperatura futura del producto,

\begin{equation}
y_i
=
T_{pad}(t_i+\Delta t),
\end{equation}

al contenido de humedad,

\begin{equation}
y_i=M(t_i+\Delta t),
\end{equation}

o al tiempo restante para alcanzar una humedad de almacenamiento determinada.

%%%%%%%%%%%%%%%%%%%%%%%%%%%%%%%%%%%%%%%%%%%%%%%%%%%%%%%%%%%%%%%%%%%%%%%%%%%%%%%%
\subsubsection{Estructura de un árbol de regresión}

Un árbol de regresión divide el espacio de entrada mediante condiciones del tipo
\begin{equation}
x_j\leq s,
\end{equation}
donde $x_j$ representa una característica y $s$ es un valor umbral seleccionado durante el entrenamiento.

Cada división separa una región padre $R$ en dos regiones hijas:

\begin{equation}
R_{\mathrm{izq}}(j,s)
=
\left\{
\mathbf{x}\mid x_j\leq s
\right\},
\end{equation}

\begin{equation}
R_{\mathrm{der}}(j,s)
=
\left\{
\mathbf{x}\mid x_j>s
\right\}.
\end{equation}

Por ejemplo, un árbol entrenado para predecir la temperatura del producto podría generar reglas conceptuales como
\[
I_t>600~\mathrm{W\,m^{-2}},
\]
\[
T_{abs}>45~^{\circ}\mathrm{C},
\]
\[
\phi_{amb}<70\,\%.
\]

Estas reglas subdividen las condiciones de funcionamiento en regímenes térmicos diferentes. De esta forma, el algoritmo puede distinguir automáticamente entre periodos de baja radiación, calentamiento intenso, elevada humedad ambiental o ventilación forzada.

Una representación simplificada sería

\begin{center}
\begin{minipage}{0.82\linewidth}
\centering
\[
\boxed{I_t\leq s_1}
\]

\[
\begin{array}{ccc}
\swarrow && \searrow\\[-0.15cm]
T_{amb}\leq s_2 && T_{abs}\leq s_3\\
\swarrow\quad\searrow && \swarrow\quad\searrow\\
\hat{T}_1\qquad\hat{T}_2 &&
\hat{T}_3\qquad\hat{T}_4
\end{array}
\]
\end{minipage}
\end{center}

Los nodos intermedios contienen las reglas de decisión, mientras que los nodos terminales u hojas almacenan las predicciones finales.

%%%%%%%%%%%%%%%%%%%%%%%%%%%%%%%%%%%%%%%%%%%%%%%%%%%%%%%%%%%%%%%%%%%%%%%%%%%%%%%%
\subsubsection{Criterio de partición}

Durante el entrenamiento, el algoritmo debe seleccionar tanto la característica $x_j$ como el umbral $s$ que produzcan la mayor reducción posible del error.

Para una región $R$, la predicción constante que minimiza el error cuadrático es la media de las observaciones contenidas en dicha región:

\begin{equation}
\hat{c}_{R}
=
\frac{1}{N_R}
\sum_{\mathbf{x}_i\in R}
y_i,
\end{equation}

donde $N_R$ es el número de observaciones pertenecientes a $R$.

El error asociado a la región se calcula mediante la suma de cuadrados residual:

\begin{equation}
\operatorname{SSE}(R)
=
\sum_{\mathbf{x}_i\in R}
\left(
y_i-\hat{c}_{R}
\right)^2.
\end{equation}

La partición óptima se obtiene resolviendo

\begin{equation}
\boxed{
(j^{*},s^{*})
=
\underset{j,s}{\arg\min}
\left[
\sum_{\mathbf{x}_i\in R_{\mathrm{izq}}(j,s)}
\left(
y_i-\hat{c}_{\mathrm{izq}}
\right)^2
+
\sum_{\mathbf{x}_i\in R_{\mathrm{der}}(j,s)}
\left(
y_i-\hat{c}_{\mathrm{der}}
\right)^2
\right]
}
\label{eq:particion_arbol}
\end{equation}

Este procedimiento se aplica de forma recursiva hasta que se satisface algún criterio de detención, como una profundidad máxima, un número mínimo de observaciones por hoja o una reducción mínima del error.

La predicción general de un árbol con $J$ hojas puede escribirse como

\begin{equation}
\boxed{
\hat{y}(\mathbf{x})
=
\sum_{m=1}^{J}
c_m
\mathbb{I}
\left(
\mathbf{x}\in R_m
\right)
}
\label{eq:prediccion_arbol}
\end{equation}

donde $R_m$ representa la región correspondiente a la hoja $m$, $c_m$ es su valor de predicción y $\mathbb{I}(\cdot)$ es la función indicadora.

\begin{tcolorbox}[
title=Interpretación física,
colback=blue!4,
colframe=blue!65!black
]
Un árbol de regresión aproxima el comportamiento térmico del secador mediante un conjunto de regímenes de operación.

Cada hoja puede interpretarse como una combinación específica de irradiancia, temperatura, humedad y ventilación para la cual la variable objetivo presenta un comportamiento aproximadamente homogéneo.
\end{tcolorbox}

%%%%%%%%%%%%%%%%%%%%%%%%%%%%%%%%%%%%%%%%%%%%%%%%%%%%%%%%%%%%%%%%%%%%%%%%%%%%%%%%
\subsubsection{Complejidad y sobreajuste}

Si el árbol continúa dividiendo los datos hasta producir hojas con muy pocas observaciones, puede ajustar no solamente la dinámica física del proceso, sino también el ruido, los errores instrumentales y las fluctuaciones particulares del conjunto de entrenamiento.

Este fenómeno se conoce como sobreajuste y se caracteriza por un error reducido sobre los datos utilizados para entrenar el modelo, pero un desempeño deficiente ante nuevas observaciones.

La complejidad del árbol puede regularse mediante hiperparámetros como

\begin{itemize}
    \item profundidad máxima;
    \item número mínimo de observaciones necesarias para dividir un nodo;
    \item número mínimo de observaciones por hoja;
    \item reducción mínima del error requerida para generar una partición;
    \item número máximo de hojas terminales.
\end{itemize}

También puede aplicarse un procedimiento de poda, en el cual se minimiza una función que combina el error de ajuste y el tamaño del árbol:

\begin{equation}
\boxed{
R_{\alpha}(T)
=
R(T)
+
\alpha |T|
}
\label{eq:poda_arbol}
\end{equation}

donde $R(T)$ representa el error del árbol, $|T|$ es el número de hojas y $\alpha$ controla la penalización asociada a la complejidad.

Un valor reducido de $\alpha$ permite árboles más complejos, mientras que un valor elevado favorece modelos más pequeños y fáciles de interpretar.

%%%%%%%%%%%%%%%%%%%%%%%%%%%%%%%%%%%%%%%%%%%%%%%%%%%%%%%%%%%%%%%%%%%%%%%%%%%%%%%%
\subsubsection{Métodos de ensamble}

Aunque un árbol individual es fácilmente interpretable, puede presentar una elevada sensibilidad frente a pequeñas modificaciones del conjunto de entrenamiento. Los métodos de ensamble reducen esta inestabilidad mediante la combinación de múltiples árboles.

La predicción colectiva puede expresarse de forma general como

\begin{equation}
\hat{y}
=
\mathcal{A}
\left[
f_1(\mathbf{x}),
f_2(\mathbf{x}),
\ldots,
f_B(\mathbf{x})
\right],
\end{equation}

donde $f_b(\mathbf{x})$ es la predicción del árbol $b$ y $\mathcal{A}$ representa una regla de agregación.

Los principales enfoques son:

\begin{enumerate}
    \item \textit{bagging}, en el cual los árboles se construyen de manera independiente y sus resultados se promedian;
    \item \textit{boosting}, en el cual los árboles se construyen secuencialmente para corregir los errores cometidos por los modelos anteriores.
\end{enumerate}

%%%%%%%%%%%%%%%%%%%%%%%%%%%%%%%%%%%%%%%%%%%%%%%%%%%%%%%%%%%%%%%%%%%%%%%%%%%%%%%%
\subsubsection{Bosques aleatorios como modelo de referencia}

Los bosques aleatorios, o \textit{Random Forests}, combinan múltiples árboles entrenados a partir de muestras diferentes del conjunto de datos. Además, en cada división se evalúa únicamente un subconjunto aleatorio de características.

Para un problema de regresión, la predicción final corresponde al promedio de los $B$ árboles:

\begin{equation}
\hat{y}_{RF}(\mathbf{x})
=
\frac{1}{B}
\sum_{b=1}^{B}
f_b(\mathbf{x}).
\end{equation}

El promedio reduce la varianza y produce predicciones más estables que las obtenidas mediante un único árbol.

Aunque los bosques aleatorios no pertenecen estrictamente a la familia de \textit{boosting}, constituyen una referencia útil para comparar el desempeño de diferentes métodos basados en árboles.

%%%%%%%%%%%%%%%%%%%%%%%%%%%%%%%%%%%%%%%%%%%%%%%%%%%%%%%%%%%%%%%%%%%%%%%%%%%%%%%%
\subsubsection{Gradient Boosting}

El método de \textit{Gradient Boosting} construye un modelo aditivo mediante una sucesión de árboles débiles. Cada nuevo árbol se entrena para reducir el error residual del conjunto formado por los árboles anteriores.

El modelo comienza con una predicción inicial constante:

\begin{equation}
F_0(\mathbf{x})
=
\underset{c}{\arg\min}
\sum_{i=1}^{N}
\mathcal{L}(y_i,c),
\end{equation}

donde $\mathcal{L}$ es la función de pérdida.

Para el error cuadrático,

\begin{equation}
\mathcal{L}(y_i,F(\mathbf{x}_i))
=
\frac{1}{2}
\left[
y_i-F(\mathbf{x}_i)
\right]^2,
\end{equation}

la predicción inicial corresponde a la media de las observaciones.

En cada iteración $m$, se calculan los pseudo-residuos:

\begin{equation}
r_{im}
=
-
\left[
\frac{
\partial
\mathcal{L}
\left(
y_i,F(\mathbf{x}_i)
\right)
}
{
\partial F(\mathbf{x}_i)
}
\right]_{F=F_{m-1}}.
\end{equation}

Para la pérdida cuadrática, estos pseudo-residuos coinciden con los residuos ordinarios:

\begin{equation}
r_{im}
=
y_i-F_{m-1}(\mathbf{x}_i).
\end{equation}

Posteriormente, se entrena un árbol $h_m(\mathbf{x})$ para aproximar dichos residuos y se actualiza el modelo:

\begin{equation}
\boxed{
F_m(\mathbf{x})
=
F_{m-1}(\mathbf{x})
+
\eta h_m(\mathbf{x})
}
\label{eq:gradient_boosting}
\end{equation}

donde $\eta$ es la tasa de aprendizaje.

Después de $M$ iteraciones, la predicción final es

\begin{equation}
\boxed{
\hat{y}(\mathbf{x})
=
F_M(\mathbf{x})
=
F_0(\mathbf{x})
+
\eta
\sum_{m=1}^{M}
h_m(\mathbf{x})
}
\end{equation}

\begin{tcolorbox}[
title=Idea fundamental del boosting,
colback=green!4,
colframe=green!55!black
]
Cada árbol individual realiza una corrección pequeña. El modelo final surge de la suma de todas estas correcciones:

\[
\text{Predicción inicial}
+
\text{corrección}_{1}
+
\text{corrección}_{2}
+
\cdots
+
\text{corrección}_{M}.
\]

De esta manera, el algoritmo puede aproximar progresivamente relaciones térmicas de elevada complejidad.
\end{tcolorbox}

%%%%%%%%%%%%%%%%%%%%%%%%%%%%%%%%%%%%%%%%%%%%%%%%%%%%%%%%%%%%%%%%%%%%%%%%%%%%%%%%
\subsubsection{Regularización del modelo}

El desempeño del \textit{boosting} depende del equilibrio entre la capacidad de ajuste y la generalización.

Los principales hiperparámetros son:

\begin{itemize}
    \item número de árboles $M$;
    \item tasa de aprendizaje $\eta$;
    \item profundidad máxima de cada árbol;
    \item número mínimo de observaciones por hoja;
    \item fracción de observaciones utilizada en cada iteración;
    \item fracción de características evaluada por árbol.
\end{itemize}

Existe una relación directa entre el número de árboles y la tasa de aprendizaje. Una tasa pequeña suele requerir un mayor número de iteraciones, pero permite que las correcciones se incorporen gradualmente:

\begin{equation}
\eta\downarrow
\qquad\Longrightarrow\qquad
M\uparrow.
\end{equation}

Los parámetros deben seleccionarse mediante validación cruzada respetando la estructura temporal de las mediciones.

Para datos de secado no resulta apropiado mezclar aleatoriamente observaciones pertenecientes al pasado y al futuro, debido a que esto puede producir fuga de información. La validación debe conservar el orden cronológico:

\begin{equation}
\underbrace{
\mathcal{D}_{train}
}_{\text{periodo inicial}}
\longrightarrow
\underbrace{
\mathcal{D}_{val}
}_{\text{periodo intermedio}}
\longrightarrow
\underbrace{
\mathcal{D}_{test}
}_{\text{periodo final}}.
\end{equation}

%%%%%%%%%%%%%%%%%%%%%%%%%%%%%%%%%%%%%%%%%%%%%%%%%%%%%%%%%%%%%%%%%%%%%%%%%%%%%%%%
\subsubsection{XGBoost}

\textit{Extreme Gradient Boosting}, comúnmente denominado XGBoost, corresponde a una implementación regularizada y computacionalmente eficiente del método de \textit{gradient boosting}.

Su función objetivo combina el error predictivo y una penalización sobre la complejidad de los árboles:

\begin{equation}
\mathcal{J}^{(m)}
=
\sum_{i=1}^{N}
\mathcal{L}
\left[
y_i,
\hat{y}_i^{(m-1)}
+
f_m(\mathbf{x}_i)
\right]
+
\Omega(f_m),
\end{equation}

donde la regularización puede escribirse como

\begin{equation}
\Omega(f)
=
\gamma J
+
\frac{\lambda}{2}
\sum_{j=1}^{J}
w_j^2.
\end{equation}

En esta expresión,

\begin{itemize}
    \item $J$ es el número de hojas del árbol;
    \item $w_j$ es el valor asociado a la hoja $j$;
    \item $\gamma$ penaliza la generación de nuevas hojas;
    \item $\lambda$ penaliza valores excesivamente elevados en las hojas.
\end{itemize}

La regularización explícita permite reducir el sobreajuste y controlar la complejidad del modelo, lo que resulta especialmente importante cuando la cantidad de datos experimentales es limitada.

%%%%%%%%%%%%%%%%%%%%%%%%%%%%%%%%%%%%%%%%%%%%%%%%%%%%%%%%%%%%%%%%%%%%%%%%%%%%%%%%
\subsubsection{Aplicación al modelado térmico del secador}

Los métodos de \textit{boosting} pueden emplearse para construir distintos modelos predictivos dentro del secador solar.

Una primera aplicación consiste en estimar la temperatura futura del aire interior:

\begin{equation}
\hat{T}_{air}(t+\Delta t)
=
f_{\theta}
\left[
I_t(t),
T_{amb}(t),
\phi_{amb}(t),
v_w(t),
T_{air}(t),
T_{abs}(t),
u_f(t)
\right].
\end{equation}

Para representar la inercia térmica pueden incluirse valores retardados:

\begin{equation}
\begin{aligned}
\hat{T}_{air}(t+\Delta t)
=
f_{\theta}
\Big[
&
T_{air}(t),
T_{air}(t-\Delta t),
T_{air}(t-2\Delta t),
\\
&
I_t(t),
I_t(t-\Delta t),
T_{amb}(t),
\phi_{amb}(t),
u_f(t)
\Big].
\end{aligned}
\end{equation}

Otra aplicación consiste en estimar el contenido de humedad del producto:

\begin{equation}
\hat{M}(t+\Delta t)
=
f_{\theta}
\left[
M(t),
T_{pad}(t),
T_{air}(t),
\phi_{air}(t),
v_{air}(t),
I_t(t)
\right].
\end{equation}

También puede predecirse la tasa instantánea de secado:

\begin{equation}
\widehat{\dot{M}}(t)
=
f_{\theta}
\left[
T_{pad}(t),
T_{air}(t),
\phi_{air}(t),
M(t),
v_{air}(t)
\right].
\end{equation}

Finalmente, el algoritmo puede utilizarse como modelo sustituto de simulaciones CFD:

\begin{equation}
\hat{\mathbf{y}}_{CFD}
=
f_{\theta}
\left(
\mathbf{p}
\right),
\end{equation}

donde $\mathbf{p}$ contiene las condiciones de simulación y $\hat{\mathbf{y}}_{CFD}$ representa indicadores obtenidos originalmente mediante el modelo numérico.

Por ejemplo,

\begin{equation}
\mathbf{p}
=
\begin{bmatrix}
I_t &
v_w &
u_f &
A_{vent} &
M_0 &
m_{pad}
\end{bmatrix}^{T},
\end{equation}

y

\begin{equation}
\hat{\mathbf{y}}_{CFD}
=
\begin{bmatrix}
\overline{T}_{air} &
T_{max} &
\overline{v}_{air} &
\Delta p &
U_T
\end{bmatrix}^{T},
\end{equation}

donde $U_T$ puede representar un indicador de uniformidad térmica.

%%%%%%%%%%%%%%%%%%%%%%%%%%%%%%%%%%%%%%%%%%%%%%%%%%%%%%%%%%%%%%%%%%%%%%%%%%%%%%%%
\subsubsection{Ingeniería de características basada en la física}

Aunque los árboles pueden detectar automáticamente relaciones no lineales, su desempeño puede mejorarse mediante variables derivadas del conocimiento físico del proceso.

Entre las características relevantes se encuentran las diferencias de temperatura

\begin{equation}
\Delta T_{abs-air}
=
T_{abs}-T_{air},
\end{equation}

\begin{equation}
\Delta T_{air-amb}
=
T_{air}-T_{amb},
\end{equation}

y el déficit de humedad del aire,

\begin{equation}
\Delta H
=
H_{sat}(T_{air})-H_{air}.
\end{equation}

También puede incorporarse la potencia solar absorbida:

\begin{equation}
\dot{Q}_{solar}
=
\alpha_{eff}A_{cap}I_t,
\end{equation}

la tasa de calentamiento aproximada,

\begin{equation}
\dot{T}_{air}(t)
\approx
\frac{
T_{air}(t)-T_{air}(t-\Delta t)
}
{\Delta t},
\end{equation}

y la energía solar acumulada:

\begin{equation}
E_{solar}(t)
=
\int_{t_0}^{t}
I_t(\tau)A_{cap}\,d\tau.
\end{equation}

Estas variables proporcionan al algoritmo información relacionada directamente con los mecanismos de transferencia de calor y humedad.

\begin{tcolorbox}[
title=Datos informados por la física,
colback=yellow!5,
colframe=orange!70!black
]
La incorporación de características como diferencias de temperatura, déficit de humedad, energía solar acumulada y valores retardados no convierte al algoritmo en un modelo físico.

Sin embargo, permite que el aprendizaje se realice sobre variables que poseen una interpretación termodinámica clara, mejorando potencialmente la precisión, la estabilidad y la capacidad de generalización.
\end{tcolorbox}

%%%%%%%%%%%%%%%%%%%%%%%%%%%%%%%%%%%%%%%%%%%%%%%%%%%%%%%%%%%%%%%%%%%%%%%%%%%%%%%%
\subsubsection{Interpretabilidad}

Una ventaja importante de los métodos basados en árboles es que permiten analizar la contribución de las variables de entrada a las predicciones.

La importancia de una característica puede estimarse a partir de la reducción acumulada del error obtenida en las divisiones donde participa dicha variable:

\begin{equation}
\operatorname{Imp}(x_j)
=
\sum_{m\in\mathcal{S}_j}
\Delta R_m,
\end{equation}

donde $\mathcal{S}_j$ representa el conjunto de divisiones realizadas con la característica $x_j$ y $\Delta R_m$ es la reducción del error producida en el nodo $m$.

También pueden utilizarse valores SHAP para expresar cada predicción como

\begin{equation}
\boxed{
\hat{y}_i
=
\phi_0
+
\sum_{j=1}^{p}
\phi_{ij}
}
\end{equation}

donde $\phi_0$ representa la predicción promedio del modelo y $\phi_{ij}$ cuantifica la contribución de la característica $j$ a la observación $i$.

Esta formulación permite responder preguntas como:

\begin{itemize}
    \item ¿cuánto contribuyó la irradiancia solar al incremento de temperatura?;
    \item ¿qué efecto tuvo la humedad ambiental sobre la tasa de secado?;
    \item ¿la ventilación redujo o incrementó la temperatura del producto?;
    \item ¿qué variables controlan principalmente la humedad final?;
    \item ¿cuáles condiciones explican una predicción anómala?
\end{itemize}

%%%%%%%%%%%%%%%%%%%%%%%%%%%%%%%%%%%%%%%%%%%%%%%%%%%%%%%%%%%%%%%%%%%%%%%%%%%%%%%%
\subsubsection{Evaluación del desempeño}

El desempeño predictivo debe evaluarse sobre datos no utilizados durante el entrenamiento.

Entre las métricas aplicables se encuentra el error absoluto medio:

\begin{equation}
\operatorname{MAE}
=
\frac{1}{N}
\sum_{i=1}^{N}
\left|
y_i-\hat{y}_i
\right|,
\end{equation}

la raíz del error cuadrático medio:

\begin{equation}
\operatorname{RMSE}
=
\sqrt{
\frac{1}{N}
\sum_{i=1}^{N}
\left(
y_i-\hat{y}_i
\right)^2
},
\end{equation}

y el coeficiente de determinación:

\begin{equation}
R^2
=
1-
\frac{
\sum_{i=1}^{N}
\left(
y_i-\hat{y}_i
\right)^2
}{
\sum_{i=1}^{N}
\left(
y_i-\bar{y}
\right)^2
}.
\end{equation}

Para la predicción del contenido de humedad también puede emplearse el error relativo:

\begin{equation}
\operatorname{MAPE}
=
\frac{100}{N}
\sum_{i=1}^{N}
\left|
\frac{
y_i-\hat{y}_i
}{
y_i
}
\right|,
\end{equation}

siempre que los valores de referencia no se encuentren próximos a cero.

La evaluación no debe limitarse a una única métrica. El MAE proporciona una interpretación directa del error promedio, mientras que el RMSE penaliza con mayor intensidad las desviaciones elevadas.

%%%%%%%%%%%%%%%%%%%%%%%%%%%%%%%%%%%%%%%%%%%%%%%%%%%%%%%%%%%%%%%%%%%%%%%%%%%%%%%%
\subsubsection{Síntesis del método}

El procedimiento general para aplicar árboles y métodos de \textit{boosting} al secador puede resumirse mediante la siguiente secuencia:

\begin{equation}
\boxed{
\begin{aligned}
&\text{Adquisición de datos}
\\
&\qquad\Downarrow
\\
&\text{Depuración y sincronización}
\\
&\qquad\Downarrow
\\
&\text{Construcción de características físicas}
\\
&\qquad\Downarrow
\\
&\text{División temporal de los datos}
\\
&\qquad\Downarrow
\\
&\text{Entrenamiento y ajuste de hiperparámetros}
\\
&\qquad\Downarrow
\\
&\text{Evaluación sobre datos no observados}
\\
&\qquad\Downarrow
\\
&\text{Interpretación de las predicciones}
\end{aligned}
}
\end{equation}

\begin{table}[H]
\centering
\caption{Comparación conceptual de métodos basados en árboles.}
\label{tab:comparacion_arboles}
\begin{tabular}{
p{0.20\linewidth}
p{0.20\linewidth}
p{0.22\linewidth}
p{0.25\linewidth}
}
\toprule
Método
&
Estrategia
&
Ventaja principal
&
Aplicación potencial
\\
\midrule

Árbol individual
&
Divisiones jerárquicas
&
Elevada interpretabilidad
&
Identificación de regímenes térmicos
\\[0.2cm]

Random Forest
&
Promedio de árboles independientes
&
Estabilidad frente al ruido
&
Modelo de referencia para datos experimentales
\\[0.2cm]

Gradient Boosting
&
Corrección secuencial de residuos
&
Elevada precisión predictiva
&
Predicción de temperatura y humedad
\\[0.2cm]

XGBoost
&
Boosting regularizado
&
Control del sobreajuste y eficiencia computacional
&
Modelos sustitutos y predicción operacional
\\
\bottomrule
\end{tabular}
\end{table}

En consecuencia, los árboles de decisión y los métodos de \textit{boosting} constituyen herramientas apropiadas para el análisis de datos térmicos tabulares, especialmente cuando el número de observaciones es moderado y se requiere capturar relaciones no lineales sin renunciar completamente a la interpretabilidad.

Su aplicación al secador permite predecir variables de difícil medición, identificar los factores que dominan el comportamiento térmico y construir aproximaciones rápidas de modelos numéricos de mayor costo computacional.



%%%%%%%%%%%%%%%%%%%%%%%%%%%%%%%%%%%%%%%%%%%%%%%%%%%%%%%%%%%%%%%%%%%%%%%%%%%%%%%%
\subsection{Physics-Informed Neural Networks (PINNs)}
\label{subsec:pinn}

Las redes neuronales artificiales convencionales aprenden una relación funcional entre un conjunto de variables de entrada y una variable objetivo utilizando exclusivamente información contenida en los datos de entrenamiento. Aunque este enfoque ha demostrado excelentes resultados en numerosos problemas de regresión y clasificación, su desempeño puede deteriorarse cuando la cantidad de datos disponibles es limitada o cuando las predicciones deben respetar leyes físicas conocidas.

En ingeniería térmica, la evolución de la temperatura, la velocidad del fluido y el contenido de humedad no constituye un fenómeno arbitrario, sino que está gobernada por ecuaciones diferenciales derivadas de los principios de conservación de masa, cantidad de movimiento y energía. En consecuencia, un modelo puramente basado en datos podría producir soluciones que ajusten adecuadamente las observaciones experimentales, pero que violen restricciones físicas fundamentales.

Con el objetivo de incorporar conocimiento físico dentro del proceso de aprendizaje, Raissi, Perdikaris y Karniadakis introdujeron en 2019 las denominadas \textit{Physics-Informed Neural Networks} (PINNs). Estas redes integran las ecuaciones diferenciales gobernantes directamente en el entrenamiento, permitiendo obtener aproximaciones que satisfacen simultáneamente los datos observados y las leyes físicas del sistema.

En lugar de aprender únicamente una función de aproximación

\[
\hat{u}=f_{\theta}(\mathbf{x}),
\]

las PINNs buscan una solución que además satisfaga una ecuación diferencial de la forma

\begin{equation}
\boxed{
\mathcal{N}[u(\mathbf{x},t)]=0,
}
\label{eq:pde_general}
\end{equation}

donde $\mathcal{N}$ representa un operador diferencial no lineal y $u(\mathbf{x},t)$ corresponde al campo físico de interés.

Dependiendo del problema, $u$ puede representar la temperatura, la presión, la velocidad del fluido, la humedad o cualquier otra variable gobernada por una ecuación diferencial parcial.

%%%%%%%%%%%%%%%%%%%%%%%%%%%%%%%%%%%%%%%%%%%%%%%%%%%%%%%%%%%%%%%%%%%%%%%%%%%%%%%%
\subsubsection{Fundamentos de las PINNs}

La idea central consiste en utilizar una red neuronal para aproximar directamente la solución de una ecuación diferencial.

Sea

\begin{equation}
u(\mathbf{x},t)
\approx
u_{\theta}(\mathbf{x},t),
\end{equation}

donde $u_{\theta}$ representa la salida de una red neuronal parametrizada por el conjunto de pesos y sesgos $\theta$.

A diferencia de una red neuronal convencional, las derivadas espaciales y temporales necesarias para evaluar la ecuación diferencial no se calculan mediante aproximaciones numéricas, sino utilizando diferenciación automática (\textit{Automatic Differentiation}).

Esta técnica permite obtener derivadas exactas de la función aproximada respecto de cualquiera de sus variables de entrada, preservando la precisión de la representación analítica implementada por la red neuronal.

Como consecuencia, el residuo físico puede escribirse como

\begin{equation}
\boxed{
r(\mathbf{x},t)
=
\mathcal{N}
\left[
u_{\theta}(\mathbf{x},t)
\right],
}
\label{eq:residual}
\end{equation}

el cual mide el grado de incumplimiento de la ecuación diferencial en cada punto del dominio.

Cuando el residuo es igual a cero,

\[
r(\mathbf{x},t)=0,
\]

la solución aproximada satisface exactamente el modelo físico considerado.

%%%%%%%%%%%%%%%%%%%%%%%%%%%%%%%%%%%%%%%%%%%%%%%%%%%%%%%%%%%%%%%%%%%%%%%%%%%%%%%%
\subsubsection{Función de pérdida informada por la física}

Una diferencia fundamental entre una red neuronal convencional y una PINN radica en la construcción de la función objetivo utilizada durante el entrenamiento.

Mientras que en el aprendizaje supervisado clásico únicamente se minimiza el error respecto de los datos experimentales,

\[
\mathcal{L}_{data},
\]

las PINNs incorporan además un término asociado al cumplimiento de las ecuaciones diferenciales.

En términos generales, la función de pérdida puede expresarse como

\begin{equation}
\boxed{
\mathcal{L}
=
\mathcal{L}_{data}
+
\lambda_f
\mathcal{L}_{physics}
+
\lambda_b
\mathcal{L}_{boundary}
+
\lambda_i
\mathcal{L}_{initial},
}
\label{eq:loss_total}
\end{equation}

donde

\begin{itemize}

\item $\mathcal{L}_{data}$ representa el error asociado a las observaciones experimentales;

\item $\mathcal{L}_{physics}$ cuantifica el residuo de la ecuación diferencial;

\item $\mathcal{L}_{boundary}$ corresponde al incumplimiento de las condiciones de frontera;

\item $\mathcal{L}_{initial}$ representa el error sobre las condiciones iniciales.

\end{itemize}

Los parámetros $\lambda_f$, $\lambda_b$ y $\lambda_i$ permiten ajustar la importancia relativa de cada contribución durante el entrenamiento.

En consecuencia, el proceso de aprendizaje no sólo busca aproximar los datos disponibles, sino también satisfacer las restricciones físicas impuestas por el problema.

%%%%%%%%%%%%%%%%%%%%%%%%%%%%%%%%%%%%%%%%%%%%%%%%%%%%%%%%%%%%%%%%%%%%%%%%%%%%%%%%
\subsubsection{Aplicación a procesos térmicos}

Las PINNs resultan especialmente atractivas para problemas de transferencia de calor, debido a que las ecuaciones gobernantes son conocidas y derivan directamente de principios físicos bien establecidos.

Por ejemplo, la evolución de la temperatura en un sólido puede describirse mediante la ecuación de conducción

\begin{equation}
\rho c_p
\frac{\partial T}{\partial t}
=
\nabla\cdot
(k\nabla T)
+\dot q,
\end{equation}

mientras que el transporte de calor en un fluido puede representarse mediante la ecuación de energía

\begin{equation}
\rho c_p
\left(
\frac{\partial T}{\partial t}
+\mathbf{u}\cdot\nabla T
\right)
=
\nabla\cdot(k\nabla T)
+\dot q.
\end{equation}

En lugar de resolver estas ecuaciones mediante discretización espacial, una PINN aproxima directamente el campo de temperatura

\[
T(\mathbf{x},t)
\]

y utiliza el residuo de las ecuaciones anteriores como parte de la función objetivo.

De esta manera, el modelo incorpora simultáneamente información experimental y conocimiento físico.

%%%%%%%%%%%%%%%%%%%%%%%%%%%%%%%%%%%%%%%%%%%%%%%%%%%%%%%%%%%%%%%%%%%%%%%%%%%%%%%%
\subsubsection{PINNs y Dinámica de Fluidos Computacional}

Las PINNs no pretenden reemplazar completamente los métodos tradicionales de Dinámica de Fluidos Computacional. En la actualidad, ambas metodologías deben considerarse herramientas complementarias.

Los métodos CFD clásicos, como el Método de Volúmenes Finitos, obtienen la solución mediante la discretización explícita del dominio espacial. Por el contrario, las PINNs aproximan directamente la solución continua utilizando una red neuronal y evalúan el cumplimiento de las ecuaciones mediante diferenciación automática.

Como consecuencia, las PINNs pueden resultar particularmente útiles cuando:

\begin{itemize}

\item existen pocos datos experimentales;

\item se requiere incorporar restricciones físicas conocidas;

\item el problema presenta geometrías complejas o dominios parcialmente observados;

\item se desea combinar simulaciones numéricas con mediciones reales.

\end{itemize}

En aplicaciones de transferencia de calor y secado solar, las PINNs representan una alternativa prometedora para integrar modelos físicos, datos experimentales y simulaciones CFD dentro de un mismo marco de aprendizaje, favoreciendo el desarrollo de modelos híbridos con mayor capacidad de generalización.





%\subsection{Ventajas y limitaciones}

%%%%%%%%%%%%%%%%%%%%%%%%%%%%%%%%%%%%%%%%%%%%%%%%%%%%%%%%%%%%%%%%%%%%%%%%%%%%%%%%
\section{Control de procesos térmicos y de humedad}
\label{sec:control_termohigrometrico}

El control de un proceso térmico consiste en modificar determinadas variables de operación con el propósito de mantener la temperatura, la humedad u otros indicadores de desempeño dentro de intervalos previamente definidos. En sistemas de secado solar, esta tarea resulta particularmente compleja debido a la interacción entre la radiación solar, la transferencia de calor, el movimiento del aire, la evaporación de humedad y las perturbaciones meteorológicas.

La dinámica del secador puede representarse, de manera general, mediante un sistema no lineal de la forma

\begin{equation}
\dot{\mathbf{x}}(t)
=
\mathbf{f}
\left[
\mathbf{x}(t),
\mathbf{u}(t),
\mathbf{d}(t),
\boldsymbol{\theta}
\right],
\label{eq:modelo_general_control}
\end{equation}

\begin{equation}
\mathbf{y}(t)
=
\mathbf{h}
\left[
\mathbf{x}(t),
\mathbf{u}(t),
\mathbf{d}(t)
\right],
\label{eq:modelo_salida_control}
\end{equation}

donde

\begin{itemize}
    \item $\mathbf{x}(t)$ representa el vector de estados térmicos e higrométricos;
    \item $\mathbf{u}(t)$ contiene las variables manipuladas;
    \item $\mathbf{d}(t)$ agrupa las perturbaciones externas;
    \item $\boldsymbol{\theta}$ representa los parámetros físicos del sistema;
    \item $\mathbf{y}(t)$ corresponde a las variables observadas o controladas.
\end{itemize}

Para un sistema de secado, el vector de estado puede comprender temperaturas de los componentes sólidos, temperatura del aire, contenido de humedad del aire y contenido de humedad del producto:

\begin{equation}
\mathbf{x}
=
\begin{bmatrix}
T_{cov} &
T_{abs} &
T_{pad} &
T_{air} &
H &
M
\end{bmatrix}^{T}.
\end{equation}

Las variables de control pueden estar asociadas al movimiento y renovación del aire, mientras que la irradiancia solar, la temperatura ambiente, la humedad exterior y la velocidad del viento actúan principalmente como perturbaciones:

\begin{equation}
\mathbf{u}
=
\begin{bmatrix}
u_1 & u_2 & \cdots & u_m
\end{bmatrix}^{T},
\end{equation}

\begin{equation}
\mathbf{d}
=
\begin{bmatrix}
I_t &
T_{amb} &
H_{amb} &
v_w
\end{bmatrix}^{T}.
\end{equation}

El objetivo general del sistema de control puede expresarse como la reducción del error entre la salida medida y una referencia deseada:

\begin{equation}
\mathbf{e}(t)
=
\mathbf{r}(t)-\mathbf{y}(t),
\label{eq:error_control}
\end{equation}

donde $\mathbf{r}(t)$ representa el conjunto de referencias de temperatura, humedad o desempeño del proceso.

Sin embargo, en un secador solar no basta necesariamente con controlar valores promedio. También debe considerarse la uniformidad espacial de las variables térmicas e higrométricas, debido a que distribuciones no uniformes pueden producir diferencias en la velocidad de secado y en el contenido final de humedad del producto.

Por tanto, el problema de control puede involucrar simultáneamente

\begin{equation}
\overline{T}(t)
\rightarrow
T_{ref}(t),
\end{equation}

\begin{equation}
\overline{H}(t)
\rightarrow
H_{ref}(t),
\end{equation}

y la minimización de la variabilidad espacial:

\begin{equation}
\sigma_T^2(t)
=
\frac{1}{V}
\int_{\Omega}
\left[
T(\mathbf{x},t)-\overline{T}(t)
\right]^2
\,dV,
\end{equation}

\begin{equation}
\sigma_H^2(t)
=
\frac{1}{V}
\int_{\Omega}
\left[
H(\mathbf{x},t)-\overline{H}(t)
\right]^2
\,dV.
\end{equation}

Desde esta perspectiva, el control termo-higrométrico puede abordarse mediante estrategias clásicas, técnicas predictivas basadas en modelos dinámicos y arquitecturas híbridas que combinan conocimiento físico y modelos aprendidos a partir de datos.

%%%%%%%%%%%%%%%%%%%%%%%%%%%%%%%%%%%%%%%%%%%%%%%%%%%%%%%%%%%%%%%%%%%%%%%%%%%%%%%%
\subsection{Control clásico}
\label{subsec:control_clasico}

El control clásico comprende un conjunto de estrategias en las cuales la acción de control se determina principalmente a partir del error presente y pasado entre la variable medida y su referencia. Estas técnicas se desarrollaron originalmente para sistemas de una entrada y una salida, aunque también pueden extenderse a configuraciones multivariables.

La estructura fundamental de un sistema de control en lazo cerrado puede representarse como

\begin{equation}
\boxed{
\text{Referencia}
\longrightarrow
\text{Controlador}
\longrightarrow
\text{Proceso}
\longrightarrow
\text{Salida}
}
\end{equation}

con una realimentación desde la salida hacia el controlador:

\begin{equation}
u(t)
=
\mathcal{C}
\left[
e(t)
\right],
\end{equation}

donde $\mathcal{C}$ representa la ley de control.

La realimentación permite corregir desviaciones producidas por perturbaciones, incertidumbres paramétricas o cambios en las condiciones de operación. Cuando la salida se aleja de la referencia, el controlador modifica la entrada manipulada con el propósito de reducir el error.

%%%%%%%%%%%%%%%%%%%%%%%%%%%%%%%%%%%%%%%%%%%%%%%%%%%%%%%%%%%%%%%%%%%%%%%%%%%%%%%%
\subsubsection{Control proporcional}

La forma más simple de control realimentado es el control proporcional:

\begin{equation}
u(t)
=
K_p e(t),
\label{eq:control_p}
\end{equation}

donde $K_p$ es la ganancia proporcional.

La acción de control aumenta en proporción directa con la magnitud del error. Si la temperatura interior se encuentra por encima de su referencia, por ejemplo, el controlador puede incrementar una acción asociada a la ventilación.

Una ganancia proporcional elevada produce una respuesta más intensa, pero puede incrementar las oscilaciones o conducir a inestabilidad. Una ganancia reducida produce respuestas más suaves, aunque puede ser insuficiente para compensar perturbaciones significativas.

El control proporcional presenta además la posibilidad de mantener un error permanente en estado estacionario. Este fenómeno ocurre cuando el sistema requiere una entrada de control distinta de cero para sostener la salida en la referencia.

%%%%%%%%%%%%%%%%%%%%%%%%%%%%%%%%%%%%%%%%%%%%%%%%%%%%%%%%%%%%%%%%%%%%%%%%%%%%%%%%
\subsubsection{Control integral}

La acción integral depende de la acumulación temporal del error:

\begin{equation}
u_I(t)
=
K_i
\int_{0}^{t}
e(\tau)\,d\tau,
\label{eq:control_i}
\end{equation}

donde $K_i$ es la ganancia integral.

Si existe un error pequeño pero persistente, su integral aumenta progresivamente y modifica la acción de control hasta eliminar la desviación estacionaria.

Esta propiedad resulta útil en procesos térmicos, donde las pérdidas de calor, las variaciones ambientales o las incertidumbres del modelo pueden producir diferencias sostenidas entre la temperatura real y la referencia.

Sin embargo, una acción integral excesiva puede provocar sobreimpulso, oscilaciones y acumulación del término integral cuando el actuador alcanza sus límites físicos. Este último fenómeno se conoce como saturación o \textit{integral windup}.

%%%%%%%%%%%%%%%%%%%%%%%%%%%%%%%%%%%%%%%%%%%%%%%%%%%%%%%%%%%%%%%%%%%%%%%%%%%%%%%%
\subsubsection{Control derivativo}

La acción derivativa responde a la velocidad de cambio del error:

\begin{equation}
u_D(t)
=
K_d
\frac{de(t)}{dt},
\label{eq:control_d}
\end{equation}

donde $K_d$ es la ganancia derivativa.

Este término introduce una acción anticipativa: si el error cambia rápidamente, el controlador reacciona antes de que alcance valores elevados. En sistemas térmicos, la acción derivativa puede contribuir a reducir el sobreimpulso asociado con la inercia térmica.

No obstante, la diferenciación amplifica las componentes de alta frecuencia y el ruido de medición, por lo que normalmente requiere filtrado.

%%%%%%%%%%%%%%%%%%%%%%%%%%%%%%%%%%%%%%%%%%%%%%%%%%%%%%%%%%%%%%%%%%%%%%%%%%%%%%%%
\subsubsection{Control PID}

El controlador proporcional-integral-derivativo combina las tres acciones anteriores:

\begin{equation}
\boxed{
u(t)
=
K_p e(t)
+
K_i
\int_{0}^{t}
e(\tau)\,d\tau
+
K_d
\frac{de(t)}{dt}
}
\label{eq:control_pid}
\end{equation}

Cada término cumple una función diferente:

\begin{itemize}
    \item la acción proporcional responde al error presente;
    \item la acción integral responde a la acumulación del error pasado;
    \item la acción derivativa responde a la tendencia futura aproximada del error.
\end{itemize}

En el dominio de Laplace, la función de transferencia ideal del controlador PID se expresa como

\begin{equation}
C(s)
=
K_p
+
\frac{K_i}{s}
+
K_d s.
\end{equation}

También puede escribirse mediante el tiempo integral $T_i$ y el tiempo derivativo $T_d$:

\begin{equation}
C(s)
=
K_p
\left(
1+
\frac{1}{T_i s}
+
T_d s
\right),
\end{equation}

con

\begin{equation}
K_i=\frac{K_p}{T_i},
\qquad
K_d=K_pT_d.
\end{equation}

\begin{tcolorbox}[
title=Interpretación del control PID,
colback=blue!4,
colframe=blue!65!black
]
El controlador PID no necesita representar explícitamente todos los mecanismos de transferencia de calor y humedad. Su acción se determina principalmente a partir del error medido.

Esta simplicidad explica su amplia utilización industrial, pero también limita su capacidad para anticipar perturbaciones meteorológicas, restricciones de operación e interacciones multivariables.
\end{tcolorbox}

%%%%%%%%%%%%%%%%%%%%%%%%%%%%%%%%%%%%%%%%%%%%%%%%%%%%%%%%%%%%%%%%%%%%%%%%%%%%%%%%
\subsubsection{Control de temperatura y humedad}

En los procesos de secado, la temperatura y la humedad se encuentran físicamente acopladas. Un incremento de la temperatura modifica la presión de saturación del vapor, mientras que la ventilación altera simultáneamente la temperatura y el contenido de humedad del aire.

De forma conceptual,

\begin{equation}
u
\longrightarrow
\begin{cases}
T_{air},\\
H_{air},\\
\dot{M}.
\end{cases}
\end{equation}

Por tanto, una misma acción de control puede producir efectos distintos e incluso opuestos sobre varias variables.

Por ejemplo, un incremento de la ventilación puede

\begin{itemize}
    \item aumentar la remoción de vapor;
    \item reducir la humedad interior;
    \item incrementar la pérdida de calor;
    \item disminuir la temperatura del aire;
    \item modificar la velocidad de secado.
\end{itemize}

Esta interacción dificulta el diseño de controladores independientes para temperatura y humedad. Cuando dos lazos de control actúan sobre variables fuertemente acopladas, la acción de uno puede perturbar el desempeño del otro.

%%%%%%%%%%%%%%%%%%%%%%%%%%%%%%%%%%%%%%%%%%%%%%%%%%%%%%%%%%%%%%%%%%%%%%%%%%%%%%%%
\subsubsection{Limitaciones del control clásico}

El control clásico resulta atractivo por su simplicidad, reducido costo computacional y facilidad de implementación. Sin embargo, presenta limitaciones importantes en sistemas termo-higrométricos no lineales:

\begin{itemize}
    \item responde principalmente después de que aparece el error;
    \item no incorpora de manera explícita predicciones meteorológicas;
    \item presenta dificultades para manejar múltiples variables acopladas;
    \item no considera naturalmente restricciones sobre actuadores y estados;
    \item su desempeño puede cambiar entre distintos regímenes de operación;
    \item no utiliza explícitamente un modelo distribuido del campo térmico.
\end{itemize}

Estas limitaciones motivan el empleo de estrategias capaces de predecir la evolución futura del sistema y optimizar las acciones de control antes de que se produzcan desviaciones significativas.

%%%%%%%%%%%%%%%%%%%%%%%%%%%%%%%%%%%%%%%%%%%%%%%%%%%%%%%%%%%%%%%%%%%%%%%%%%%%%%%%
\subsection{Control predictivo}
\label{subsec:control_predictivo}

El control predictivo basado en modelos, conocido como \textit{Model Predictive Control} (MPC), es una estrategia de control avanzado que utiliza un modelo dinámico para predecir la evolución futura del sistema dentro de un horizonte temporal.

A diferencia del control clásico, que determina la acción principalmente a partir del error actual, el control predictivo resuelve un problema de optimización en cada instante de muestreo.

Su principio de funcionamiento puede representarse mediante la secuencia

\begin{equation}
\boxed{
\text{Estado actual}
\rightarrow
\text{Predicción futura}
\rightarrow
\text{Optimización}
\rightarrow
\text{Acción de control}
}
\end{equation}

Una vez calculada la secuencia óptima de acciones, únicamente se aplica la primera. En el siguiente instante se actualiza el estado del sistema, se desplaza el horizonte y se resuelve nuevamente el problema. Este procedimiento se denomina estrategia de horizonte deslizante o \textit{receding horizon}.

%%%%%%%%%%%%%%%%%%%%%%%%%%%%%%%%%%%%%%%%%%%%%%%%%%%%%%%%%%%%%%%%%%%%%%%%%%%%%%%%
\subsubsection{Modelo de predicción}

Para un sistema lineal discreto, el modelo puede escribirse como

\begin{equation}
\mathbf{x}_{k+1}
=
\mathbf{A}\mathbf{x}_k
+
\mathbf{B}\mathbf{u}_k
+
\mathbf{E}\mathbf{d}_k,
\label{eq:modelo_estado_mpc}
\end{equation}

\begin{equation}
\mathbf{y}_k
=
\mathbf{C}\mathbf{x}_k
+
\mathbf{D}\mathbf{u}_k.
\label{eq:modelo_salida_mpc}
\end{equation}

La matriz $\mathbf{A}$ describe la evolución interna de los estados, $\mathbf{B}$ la influencia de las acciones de control, $\mathbf{E}$ el efecto de las perturbaciones y $\mathbf{C}$ la relación entre los estados y las salidas observadas.

A partir del estado actual, el modelo permite predecir la salida futura:

\begin{equation}
\hat{\mathbf{y}}_{k+j|k},
\qquad
j=1,\ldots,N_p,
\end{equation}

donde $N_p$ es el horizonte de predicción y la notación $k+j|k$ representa una predicción para el instante $k+j$ calculada utilizando la información disponible en $k$.

Para procesos no lineales puede emplearse una formulación general:

\begin{equation}
\mathbf{x}_{k+1}
=
\mathbf{f}
\left(
\mathbf{x}_k,
\mathbf{u}_k,
\mathbf{d}_k
\right),
\end{equation}

\begin{equation}
\mathbf{y}_k
=
\mathbf{h}
\left(
\mathbf{x}_k,
\mathbf{u}_k
\right).
\end{equation}

Esta formulación conduce al control predictivo no lineal, o \textit{Nonlinear Model Predictive Control}.

%%%%%%%%%%%%%%%%%%%%%%%%%%%%%%%%%%%%%%%%%%%%%%%%%%%%%%%%%%%%%%%%%%%%%%%%%%%%%%%%
\subsubsection{Función objetivo}

La secuencia de control se determina mediante la minimización de una función de costo. Una formulación cuadrática típica es

\begin{equation}
\boxed{
J
=
\sum_{j=1}^{N_p}
\left\|
\hat{\mathbf{y}}_{k+j|k}
-
\mathbf{r}_{k+j}
\right\|_{\mathbf{Q}}^2
+
\sum_{j=0}^{N_c-1}
\left\|
\Delta\mathbf{u}_{k+j|k}
\right\|_{\mathbf{R}}^2
}
\label{eq:costo_mpc}
\end{equation}

donde

\begin{itemize}
    \item $N_p$ es el horizonte de predicción;
    \item $N_c$ es el horizonte de control;
    \item $\mathbf{Q}$ pondera el error respecto de la referencia;
    \item $\mathbf{R}$ penaliza cambios excesivos en los actuadores;
    \item $\Delta\mathbf{u}_k=\mathbf{u}_k-\mathbf{u}_{k-1}$.
\end{itemize}

La notación cuadrática se define como

\begin{equation}
\|\mathbf{z}\|_{\mathbf{Q}}^2
=
\mathbf{z}^{T}\mathbf{Q}\mathbf{z}.
\end{equation}

En un proceso de secado, la función objetivo puede incorporar varios criterios:

\begin{equation}
J
=
J_T
+
J_H
+
J_M
+
J_u
+
J_{\Delta u},
\end{equation}

donde los términos pueden representar respectivamente el seguimiento de temperatura, humedad, contenido de humedad del producto, consumo energético y suavidad de las acciones de control.

También puede incluirse la uniformidad espacial:

\begin{equation}
J_{\mathrm{unif}}
=
w_T\sigma_T^2
+
w_H\sigma_H^2.
\end{equation}

De este modo, el controlador puede buscar no solamente una temperatura media determinada, sino también una distribución térmica más uniforme.

%%%%%%%%%%%%%%%%%%%%%%%%%%%%%%%%%%%%%%%%%%%%%%%%%%%%%%%%%%%%%%%%%%%%%%%%%%%%%%%%
\subsubsection{Restricciones}

Una de las principales ventajas del control predictivo es su capacidad para incorporar restricciones físicas y operacionales de manera explícita.

Las acciones de control pueden estar limitadas por

\begin{equation}
\mathbf{u}_{min}
\leq
\mathbf{u}_k
\leq
\mathbf{u}_{max},
\end{equation}

mientras que la velocidad de cambio de los actuadores puede restringirse mediante

\begin{equation}
\Delta\mathbf{u}_{min}
\leq
\Delta\mathbf{u}_k
\leq
\Delta\mathbf{u}_{max}.
\end{equation}

Las variables de estado o salida también pueden mantenerse dentro de intervalos de seguridad:

\begin{equation}
\mathbf{y}_{min}
\leq
\mathbf{y}_k
\leq
\mathbf{y}_{max}.
\end{equation}

En un proceso térmico, estas restricciones pueden representar

\begin{itemize}
    \item temperaturas máximas admisibles;
    \item intervalos de humedad requeridos;
    \item capacidades máximas de ventilación;
    \item limitaciones de potencia;
    \item velocidades máximas de variación;
    \item condiciones de protección del producto.
\end{itemize}

%%%%%%%%%%%%%%%%%%%%%%%%%%%%%%%%%%%%%%%%%%%%%%%%%%%%%%%%%%%%%%%%%%%%%%%%%%%%%%%%
\subsubsection{Incorporación de perturbaciones}

La radiación solar y las condiciones ambientales varían durante la operación. Si se dispone de estimaciones o pronósticos, estas perturbaciones pueden incorporarse al modelo predictivo:

\begin{equation}
\hat{\mathbf{d}}_{k+j|k}
=
\begin{bmatrix}
\hat{I}_{t,k+j|k} &
\hat{T}_{amb,k+j|k} &
\hat{H}_{amb,k+j|k} &
\hat{v}_{w,k+j|k}
\end{bmatrix}^{T}.
\end{equation}

Esto permite anticipar cambios en lugar de reaccionar únicamente después de que la salida se haya desviado.

Por ejemplo, ante una reducción prevista de la irradiancia solar, el controlador puede modificar anticipadamente la ventilación para evitar una disminución excesiva de la temperatura interna.

\begin{tcolorbox}[
title=Carácter anticipativo del MPC,
colback=green!4,
colframe=green!55!black
]
El control clásico corrige principalmente las consecuencias observadas de una perturbación.

El control predictivo puede utilizar un modelo y una estimación de las perturbaciones futuras para modificar la operación antes de que dichas consecuencias se manifiesten completamente.
\end{tcolorbox}

%%%%%%%%%%%%%%%%%%%%%%%%%%%%%%%%%%%%%%%%%%%%%%%%%%%%%%%%%%%%%%%%%%%%%%%%%%%%%%%%
\subsubsection{Estimación de estados}

No todos los estados necesarios para el modelo pueden medirse directamente. En particular, el contenido de humedad interno del producto, determinados campos térmicos o algunas variables distribuidas pueden ser difíciles de observar en tiempo real.

Por esta razón, el control predictivo puede complementarse con un estimador de estados:

\begin{equation}
\hat{\mathbf{x}}_k
=
\mathcal{E}
\left(
\mathbf{y}_{0:k},
\mathbf{u}_{0:k-1}
\right),
\end{equation}

donde $\mathcal{E}$ representa un observador o algoritmo de estimación.

La arquitectura general se convierte entonces en

\begin{equation}
\boxed{
\text{Mediciones}
\rightarrow
\text{Estimador de estados}
\rightarrow
\text{MPC}
\rightarrow
\text{Proceso}
}
\end{equation}

El modelo puede utilizarse para reconstruir estados no medidos a partir de la información disponible.

%%%%%%%%%%%%%%%%%%%%%%%%%%%%%%%%%%%%%%%%%%%%%%%%%%%%%%%%%%%%%%%%%%%%%%%%%%%%%%%%
\subsubsection{Ventajas y limitaciones del control predictivo}

Entre las principales ventajas del MPC se encuentran

\begin{itemize}
    \item capacidad para manejar múltiples entradas y salidas;
    \item consideración explícita de restricciones;
    \item utilización de predicciones de perturbaciones;
    \item optimización simultánea de varios objetivos;
    \item anticipación de la dinámica futura;
    \item incorporación de modelos físicos del proceso.
\end{itemize}

No obstante, su desempeño depende de la calidad del modelo de predicción. Un modelo excesivamente simplificado puede producir decisiones inadecuadas, mientras que un modelo de alta fidelidad puede resultar demasiado costoso para su utilización en tiempo real.

Esta tensión entre fidelidad física y costo computacional motiva el desarrollo de modelos reducidos, sustitutos e híbridos.

%%%%%%%%%%%%%%%%%%%%%%%%%%%%%%%%%%%%%%%%%%%%%%%%%%%%%%%%%%%%%%%%%%%%%%%%%%%%%%%%
\subsection{Control basado en modelos híbridos}
\label{subsec:control_modelos_hibridos}

Los modelos híbridos combinan conocimiento físico y aprendizaje basado en datos dentro de una misma representación dinámica. Su propósito es aprovechar la capacidad explicativa de las ecuaciones gobernantes y, al mismo tiempo, compensar fenómenos no modelados, incertidumbres paramétricas o discrepancias respecto del sistema real.

En términos generales, un modelo híbrido puede representarse como

\begin{equation}
\boxed{
\text{Modelo híbrido}
=
\text{Modelo físico}
+
\text{Corrección basada en datos}
}
\end{equation}

Matemáticamente,

\begin{equation}
\dot{\mathbf{x}}
=
\mathbf{f}_{fis}
\left(
\mathbf{x},
\mathbf{u},
\mathbf{d},
\boldsymbol{\theta}
\right)
+
\mathbf{f}_{ML}
\left(
\mathbf{x},
\mathbf{u},
\mathbf{d};
\boldsymbol{\vartheta}
\right),
\label{eq:modelo_hibrido_dinamico}
\end{equation}

donde $\mathbf{f}_{fis}$ representa la dinámica conocida y $\mathbf{f}_{ML}$ una corrección aprendida a partir de datos.

Esta formulación no implica necesariamente que ambos componentes posean la misma importancia. El modelo físico puede constituir la estructura principal, mientras que el modelo basado en datos se limita a representar discrepancias residuales.

%%%%%%%%%%%%%%%%%%%%%%%%%%%%%%%%%%%%%%%%%%%%%%%%%%%%%%%%%%%%%%%%%%%%%%%%%%%%%%%%
\subsubsection{Modelado residual}

Sea la predicción del modelo físico

\begin{equation}
\hat{\mathbf{y}}_{fis}
=
\mathcal{M}_{fis}
\left(
\mathbf{x},
\mathbf{u},
\mathbf{d}
\right).
\end{equation}

La diferencia entre la observación real y la predicción física es

\begin{equation}
\mathbf{r}
=
\mathbf{y}_{real}
-
\hat{\mathbf{y}}_{fis}.
\label{eq:residuo_modelo_fisico}
\end{equation}

Un modelo basado en datos puede aproximar este residuo:

\begin{equation}
\hat{\mathbf{r}}
=
\mathcal{M}_{ML}
\left(
\mathbf{x},
\mathbf{u},
\mathbf{d}
\right).
\end{equation}

La predicción híbrida resultante es

\begin{equation}
\boxed{
\hat{\mathbf{y}}_{hib}
=
\hat{\mathbf{y}}_{fis}
+
\hat{\mathbf{r}}
}
\label{eq:prediccion_hibrida}
\end{equation}

Esta arquitectura resulta útil cuando el modelo físico reproduce correctamente la tendencia general, pero presenta errores sistemáticos debidos a parámetros inciertos, pérdidas térmicas no representadas o simplificaciones geométricas.

%%%%%%%%%%%%%%%%%%%%%%%%%%%%%%%%%%%%%%%%%%%%%%%%%%%%%%%%%%%%%%%%%%%%%%%%%%%%%%%%
\subsubsection{Modelos físicos con parámetros aprendidos}

Otra forma de hibridación consiste en conservar completamente la estructura física y utilizar técnicas de aprendizaje para estimar parámetros desconocidos:

\begin{equation}
\dot{\mathbf{x}}
=
\mathbf{f}_{fis}
\left[
\mathbf{x},
\mathbf{u},
\mathbf{d},
\hat{\boldsymbol{\theta}}
\right],
\end{equation}

con

\begin{equation}
\hat{\boldsymbol{\theta}}
=
g_{\boldsymbol{\vartheta}}
\left(
\mathbf{x},
\mathbf{u},
\mathbf{d}
\right).
\end{equation}

Los parámetros aprendidos pueden representar, por ejemplo,

\begin{itemize}
    \item coeficientes convectivos efectivos;
    \item propiedades térmicas aparentes;
    \item tasas efectivas de secado;
    \item coeficientes de infiltración;
    \item pérdidas térmicas globales;
    \item parámetros radiativos equivalentes.
\end{itemize}

En este caso, el aprendizaje automático no sustituye la ecuación física, sino que estima cantidades difíciles de determinar directamente.

%%%%%%%%%%%%%%%%%%%%%%%%%%%%%%%%%%%%%%%%%%%%%%%%%%%%%%%%%%%%%%%%%%%%%%%%%%%%%%%%
\subsubsection{Modelos sustitutos}

Los modelos de alta fidelidad, como las simulaciones CFD, pueden proporcionar una representación detallada de los campos de temperatura, velocidad y humedad. Sin embargo, su costo computacional puede impedir su utilización directa dentro de un controlador en tiempo real.

Un modelo sustituto aproxima la relación entrada--salida del modelo de alta fidelidad:

\begin{equation}
\mathbf{y}_{CFD}
=
\mathcal{F}_{CFD}
\left(
\mathbf{x},
\mathbf{u},
\mathbf{d}
\right),
\end{equation}

\begin{equation}
\hat{\mathbf{y}}_{surr}
=
\mathcal{F}_{ML}
\left(
\mathbf{x},
\mathbf{u},
\mathbf{d}
\right).
\end{equation}

El objetivo es obtener

\begin{equation}
\mathcal{F}_{ML}
\approx
\mathcal{F}_{CFD}
\end{equation}

dentro del dominio de operación considerado, pero con un tiempo de evaluación significativamente menor.

La arquitectura conceptual puede expresarse como

\begin{equation}
\boxed{
\text{CFD de alta fidelidad}
\rightarrow
\text{datos de simulación}
\rightarrow
\text{modelo sustituto}
\rightarrow
\text{control predictivo}
}
\end{equation}

No obstante, un modelo sustituto reproduce principalmente el comportamiento del modelo que generó sus datos. Por tanto, si el CFD contiene errores o se utiliza fuera de su dominio validado, dichas limitaciones pueden transmitirse al modelo aprendido.

%%%%%%%%%%%%%%%%%%%%%%%%%%%%%%%%%%%%%%%%%%%%%%%%%%%%%%%%%%%%%%%%%%%%%%%%%%%%%%%%
\subsubsection{PINNs dentro de una arquitectura de control}

Las redes neuronales informadas por la física constituyen otra forma de modelo híbrido. En ellas, las ecuaciones diferenciales, las condiciones iniciales y las condiciones de frontera se incorporan a la función de pérdida.

Una PINN puede aproximar el estado del sistema mediante

\begin{equation}
\hat{\mathbf{x}}
=
\mathbf{x}_{\boldsymbol{\vartheta}}
\left(
\mathbf{r},
t,
\mathbf{u},
\mathbf{d}
\right),
\end{equation}

mientras que el residuo físico se expresa como

\begin{equation}
\mathbf{r}_{fis}
=
\frac{\partial\mathbf{x}_{\boldsymbol{\vartheta}}}{\partial t}
-
\mathbf{f}_{fis}
\left(
\mathbf{x}_{\boldsymbol{\vartheta}},
\mathbf{u},
\mathbf{d}
\right).
\end{equation}

La función de pérdida general incluye información experimental y restricciones físicas:

\begin{equation}
\mathcal{L}
=
\lambda_{data}\mathcal{L}_{data}
+
\lambda_{fis}\mathcal{L}_{fis}
+
\lambda_{bc}\mathcal{L}_{bc}
+
\lambda_{ic}\mathcal{L}_{ic}.
\end{equation}

En una arquitectura de control, una PINN podría utilizarse como

\begin{itemize}
    \item estimador de estados no observados;
    \item modelo de predicción;
    \item identificador de parámetros;
    \item aproximación de campos térmicos;
    \item corrector de un modelo físico simplificado.
\end{itemize}

Sin embargo, la utilización de una PINN dentro del lazo de control exige que su tiempo de evaluación, estabilidad y capacidad de generalización sean compatibles con los requisitos del proceso.

%%%%%%%%%%%%%%%%%%%%%%%%%%%%%%%%%%%%%%%%%%%%%%%%%%%%%%%%%%%%%%%%%%%%%%%%%%%%%%%%
\subsubsection{Control predictivo híbrido}

Un controlador predictivo basado en un modelo híbrido utiliza

\begin{equation}
\mathbf{x}_{k+1}
=
\mathbf{f}_{hib}
\left(
\mathbf{x}_k,
\mathbf{u}_k,
\mathbf{d}_k
\right),
\end{equation}

donde

\begin{equation}
\mathbf{f}_{hib}
=
\mathbf{f}_{fis}
+
\mathbf{f}_{ML}.
\end{equation}

El problema de optimización conserva la estructura general del MPC:

\begin{equation}
\begin{aligned}
\min_{\mathbf{u}_{k:k+N_c-1}}
\quad &
J
\\
\text{sujeto a}
\quad &
\mathbf{x}_{k+j+1}
=
\mathbf{f}_{hib}
\left(
\mathbf{x}_{k+j},
\mathbf{u}_{k+j},
\mathbf{d}_{k+j}
\right),
\\
&
\mathbf{u}_{min}
\leq
\mathbf{u}_{k+j}
\leq
\mathbf{u}_{max},
\\
&
\mathbf{y}_{min}
\leq
\mathbf{y}_{k+j}
\leq
\mathbf{y}_{max}.
\end{aligned}
\label{eq:mpc_hibrido}
\end{equation}

La ventaja potencial consiste en mejorar la precisión predictiva sin abandonar completamente las restricciones físicas.

%%%%%%%%%%%%%%%%%%%%%%%%%%%%%%%%%%%%%%%%%%%%%%%%%%%%%%%%%%%%%%%%%%%%%%%%%%%%%%%%
\subsubsection{Actualización y adaptación}

Los procesos térmicos pueden cambiar debido a variaciones en la carga de producto, envejecimiento de materiales, modificaciones geométricas o cambios estacionales.

Un modelo híbrido puede actualizar algunos de sus parámetros utilizando nuevas observaciones:

\begin{equation}
\boldsymbol{\vartheta}_{k+1}
=
\boldsymbol{\vartheta}_{k}
+
\Delta\boldsymbol{\vartheta}_{k},
\end{equation}

donde $\Delta\boldsymbol{\vartheta}_{k}$ representa una corrección derivada del error entre las predicciones y las mediciones.

Esta capacidad conduce a modelos adaptativos que modifican su representación conforme cambia el proceso.

Sin embargo, la adaptación debe limitarse para evitar que el componente basado en datos produzca comportamientos físicamente inconsistentes o inestables.

%%%%%%%%%%%%%%%%%%%%%%%%%%%%%%%%%%%%%%%%%%%%%%%%%%%%%%%%%%%%%%%%%%%%%%%%%%%%%%%%
\subsubsection{Validez y dominio operacional}

La aplicación de modelos híbridos requiere definir el dominio dentro del cual existe evidencia suficiente para respaldar sus predicciones.

Sea $\mathcal{D}_{val}$ el dominio validado experimentalmente y $\mathcal{D}_{sim}$ el dominio cubierto mediante simulaciones. En general,

\begin{equation}
\mathcal{D}_{val}
\subseteq
\mathcal{D}_{sim}.
\end{equation}

La existencia de simulaciones en $\mathcal{D}_{sim}$ no garantiza automáticamente que el modelo físico sea válido en todo dicho dominio. Por tanto, es necesario distinguir entre

\begin{equation}
\text{interpolación dentro del dominio validado},
\end{equation}

\begin{equation}
\text{predicción dentro del dominio simulado},
\end{equation}

y

\begin{equation}
\text{extrapolación fuera del dominio conocido}.
\end{equation}

La incertidumbre asociada a la predicción debería aumentar al alejarse de las regiones respaldadas por mediciones.

\begin{tcolorbox}[
title=Principio de validez,
colback=yellow!5,
colframe=orange!70!black
]
Un modelo híbrido no adquiere validez física únicamente por incluir ecuaciones gobernantes o datos CFD.

Su confiabilidad depende de la validez de las ecuaciones, la calidad de los datos, la identificación de parámetros y la proximidad de las condiciones evaluadas al dominio experimentalmente respaldado.
\end{tcolorbox}

%%%%%%%%%%%%%%%%%%%%%%%%%%%%%%%%%%%%%%%%%%%%%%%%%%%%%%%%%%%%%%%%%%%%%%%%%%%%%%%%
\subsubsection{Arquitectura conceptual de control híbrido}

La integración de los distintos niveles de modelado puede resumirse mediante

\begin{equation}
\boxed{
\begin{aligned}
&\text{Leyes de conservación}
\\
&\qquad\Downarrow
\\
&\text{Modelo térmico y CFD}
\\
&\qquad\Downarrow
\\
&\text{Validación experimental}
\\
&\qquad\Downarrow
\\
&\text{Modelo reducido o híbrido}
\\
&\qquad\Downarrow
\\
&\text{Estimación de estados}
\\
&\qquad\Downarrow
\\
&\text{Control predictivo}
\\
&\qquad\Downarrow
\\
&\text{Acciones de operación}
\end{aligned}
}
\end{equation}

Esta arquitectura establece una separación entre el modelo de alta fidelidad, utilizado para comprender el fenómeno y generar información, y el modelo reducido, utilizado para realizar predicciones rápidas dentro del sistema de control.

%%%%%%%%%%%%%%%%%%%%%%%%%%%%%%%%%%%%%%%%%%%%%%%%%%%%%%%%%%%%%%%%%%%%%%%%%%%%%%%%
\subsubsection{Síntesis comparativa}

\begin{table}
\centering
\caption{Comparación conceptual de estrategias de control para procesos termo-higrométricos.}
\label{tab:comparacion_control}
\begin{tabular}{
p{0.18\linewidth}
p{0.21\linewidth}
p{0.23\linewidth}
p{0.27\linewidth}
}
\toprule
Estrategia
&
Base de decisión
&
Ventaja principal
&
Limitación principal
\\
\midrule

Control clásico
&
Error actual y acumulado
&
Simplicidad y bajo costo computacional
&
Capacidad limitada para anticipar perturbaciones y manejar restricciones
\\[0.25cm]

Control predictivo
&
Modelo dinámico y optimización
&
Anticipación, manejo multivariable y restricciones explícitas
&
Dependencia de la precisión del modelo y del costo de optimización
\\[0.25cm]

Control híbrido
&
Modelo físico complementado con datos
&
Mayor flexibilidad ante incertidumbres y fenómenos no modelados
&
Requiere validación, control de incertidumbre y supervisión del dominio operacional
\\
\bottomrule
\end{tabular}
\end{table}

En síntesis, el control clásico proporciona una referencia fundamental y una solución de baja complejidad. El control predictivo incorpora anticipación y optimización, mientras que los modelos híbridos permiten combinar la estructura física del proceso con información derivada de datos experimentales o simulaciones.


\begin{table}[h!]
\centering
\caption{Comparación conceptual de estrategias de control aplicadas a procesos termo-higrométricos.}
\label{tab:comparacion_estrategias_control}
\renewcommand{\arraystretch}{1.45}
\setlength{\tabcolsep}{7pt}

\begin{tabular}{
>{\centering\arraybackslash}p{0.16\textwidth}
>{\raggedright\arraybackslash}p{0.22\textwidth}
>{\raggedright\arraybackslash}p{0.24\textwidth}
>{\raggedright\arraybackslash}p{0.27\textwidth}
}
\toprule
\textbf{Estrategia}
&
\textbf{Principio de operación}
&
\textbf{Fortaleza principal}
&
\textbf{Consideración fundamental}
\\
\midrule

\cellcolor{blue!8}
\textbf{Control clásico}
&
Determina la acción de control a partir del error entre la variable medida y su referencia.
&
Proporciona una solución de baja complejidad, fácil implementación y reducido costo computacional.
&
Su capacidad de anticipación es limitada y puede presentar dificultades ante procesos no lineales, multivariables o sometidos a perturbaciones significativas.
\\

\addlinespace[0.15cm]

\cellcolor{green!8}
\textbf{Control predictivo}
&
Utiliza un modelo dinámico para predecir la evolución futura del proceso y resolver un problema de optimización.
&
Permite anticipar perturbaciones, manejar restricciones y coordinar múltiples variables de control.
&
Su desempeño depende de la fidelidad del modelo de predicción y de la viabilidad computacional de la optimización en tiempo real.
\\

\addlinespace[0.15cm]

\cellcolor{orange!10}
\textbf{Control basado en modelos híbridos}
&
Combina ecuaciones físicas con información obtenida de datos experimentales o simulaciones.
&
Integra interpretabilidad física, capacidad de adaptación y corrección de fenómenos no modelados.
&
Requiere validación rigurosa, control de la incertidumbre y delimitación clara del dominio operacional.
\\

\bottomrule
\end{tabular}
\end{table}



%%%%%%%%%%%%%%%%%%%%%%%%%%%%%%%%%%%%%%%%%%%%%%%%%%%%%%%%%%%%%%%%%%%%%%%%%%%%%%%%
\section{Marco conceptual}
\label{sec:marco_conceptual}

El marco conceptual de esta investigación establece una representación unificada del secador solar de café como un sistema termo-higrométrico dinámico, abierto, multivariable y sometido a perturbaciones ambientales. Sobre esta representación se articulan el problema de investigación, los objetivos, las hipótesis de trabajo, la estrategia de modelado y las contribuciones propuestas.

La investigación parte de la premisa de que el desempeño de un secador solar no depende de una única variable ni de un único mecanismo de transferencia, sino de la interacción entre la radiación solar, la acumulación de energía en los componentes sólidos, el movimiento del aire, el transporte de vapor, la evaporación de humedad desde el producto y las acciones de control aplicadas al sistema.

Por tanto, el secador no se interpreta exclusivamente como una estructura pasiva para la captación de energía solar, sino como una planta dinámica cuyo comportamiento puede ser observado, modelado, predicho y modificado.

La arquitectura conceptual general puede expresarse mediante la cadena

\begin{equation}
\boxed{
\begin{aligned}
&\text{Entorno meteorológico}
\\[-0.05cm]
&\qquad\Downarrow
\\[-0.05cm]
&\text{Interacciones termo-higrométricas}
\\[-0.05cm]
&\qquad\Downarrow
\\[-0.05cm]
&\text{Estados internos del secador}
\\[-0.05cm]
&\qquad\Downarrow
\\[-0.05cm]
&\text{Proceso de secado del producto}
\\[-0.05cm]
&\qquad\Downarrow
\\[-0.05cm]
&\text{Medición, predicción y control}
\end{aligned}
}
\label{eq:cadena_conceptual_general}
\end{equation}

Esta secuencia no debe entenderse como una relación estrictamente unidireccional. Las interacciones internas se encuentran acopladas: las temperaturas superficiales modifican la convección, la evaporación altera la temperatura del producto y del aire, la ventilación modifica simultáneamente la extracción de humedad y las pérdidas térmicas, mientras que las acciones de control cambian la evolución futura del sistema.

%%%%%%%%%%%%%%%%%%%%%%%%%%%%%%%%%%%%%%%%%%%%%%%%%%%%%%%%%%%%%%%%%%%%%%%%%%%%%%%%
\subsection{Concepción del sistema físico}
\label{subsec:concepcion_sistema_fisico}

El objeto de estudio se concibe como la integración de tres subsistemas principales:

\begin{enumerate}
    \item el entorno meteorológico;
    \item la planta termo-higrométrica;
    \item el sistema de observación y control.
\end{enumerate}

El entorno meteorológico comprende las condiciones externas que suministran energía al proceso o perturban su funcionamiento. Estas variables incluyen, entre otras, la irradiancia solar, la temperatura ambiente, la humedad exterior y la velocidad del viento.

La planta termo-higrométrica incluye la cubierta, las superficies absorbentes, el aire interior, la estructura sólida, el producto y los mecanismos de intercambio de calor y masa.

El sistema de observación y control comprende los sensores, estimadores, modelos de predicción, algoritmos de decisión y actuadores utilizados para modificar la operación.

En forma compacta,

\begin{equation}
\boxed{
\mathcal{S}
=
\mathcal{E}
\cup
\mathcal{P}
\cup
\mathcal{O}
\cup
\mathcal{C},
}
\label{eq:sistema_conceptual}
\end{equation}

donde

\begin{itemize}
    \item $\mathcal{E}$ representa el entorno;
    \item $\mathcal{P}$ representa la planta física;
    \item $\mathcal{O}$ representa el sistema de observación;
    \item $\mathcal{C}$ representa el sistema de control.
\end{itemize}

La separación anterior es conceptual y funcional. En la práctica, los cuatro componentes interactúan de manera continua.

%%%%%%%%%%%%%%%%%%%%%%%%%%%%%%%%%%%%%%%%%%%%%%%%%%%%%%%%%%%%%%%%%%%%%%%%%%%%%%%%
\subsection{Estado termo-higrométrico del sistema}
\label{subsec:estado_termohigrometrico}

El estado del secador se define como el conjunto mínimo de variables internas necesarias para describir su condición presente y predecir, junto con las entradas y perturbaciones, su evolución futura.

De manera general, el vector de estados puede escribirse como

\begin{equation}
\mathbf{x}_s(t)
=
\begin{bmatrix}
\mathbf{T}_s(t)\\
\mathbf{T}_f(t)\\
\mathbf{H}(t)\\
\mathbf{M}(t)
\end{bmatrix},
\label{eq:estado_secador}
\end{equation}

donde

\begin{itemize}
    \item $\mathbf{T}_s$ representa las temperaturas de los componentes sólidos;
    \item $\mathbf{T}_f$ representa las temperaturas del dominio fluido;
    \item $\mathbf{H}$ representa las variables asociadas con el contenido de vapor en el aire;
    \item $\mathbf{M}$ representa el contenido de humedad del producto.
\end{itemize}

Cuando el comportamiento espacial es relevante, estas variables deben interpretarse como campos:

\begin{equation}
T=T(\mathbf{r},t),
\qquad
H=H(\mathbf{r},t),
\qquad
M=M(\mathbf{r},t).
\end{equation}

Cuando se emplea una representación concentrada, los campos se sustituyen por estados nodales o valores promedio que conservan únicamente los grados de libertad considerados dominantes.

Esta distinción permite establecer dos escalas complementarias de representación:

\begin{equation}
\boxed{
\text{Campos distribuidos}
\longleftrightarrow
\text{Variables agregadas}
}
\end{equation}

El modelo distribuido describe la heterogeneidad espacial, mientras que el modelo agregado reduce la complejidad para facilitar la identificación, predicción y síntesis de control.

%%%%%%%%%%%%%%%%%%%%%%%%%%%%%%%%%%%%%%%%%%%%%%%%%%%%%%%%%%%%%%%%%%%%%%%%%%%%%%%%
\subsection{Entorno como sistema dinámico exógeno}
\label{subsec:entorno_dinamico}

Las condiciones ambientales no se consideran parámetros constantes, sino variables dinámicas que modifican permanentemente la respuesta del secador.

El estado del entorno puede representarse mediante

\begin{equation}
\mathbf{x}_e(t)
=
\begin{bmatrix}
I_t(t) &
T_{amb}(t) &
H_{amb}(t) &
v_w(t)
\end{bmatrix}^{T},
\label{eq:estado_entorno}
\end{equation}

donde $I_t$ es la irradiancia incidente, $T_{amb}$ la temperatura ambiente, $H_{amb}$ la humedad exterior y $v_w$ la velocidad del viento.

La integración conceptual entre la planta y su entorno conduce al estado aumentado

\begin{equation}
\boxed{
\mathbf{z}(t)
=
\begin{bmatrix}
\mathbf{x}_s(t)\\
\mathbf{x}_e(t)
\end{bmatrix}.
}
\label{eq:metaestado}
\end{equation}

Este estado aumentado, denominado en esta investigación \textit{metaestado sistema--entorno}, no implica que el ambiente sea directamente controlable. Su función consiste en reconocer que la evolución de la planta no puede analizarse correctamente si las perturbaciones meteorológicas se consideran únicamente como ruido no estructurado.

La dinámica general se representa mediante

\begin{equation}
\dot{\mathbf{x}}_s
=
\mathbf{f}_s
\left(
\mathbf{x}_s,
\mathbf{x}_e,
\mathbf{u},
\boldsymbol{\theta}
\right),
\label{eq:dinamica_sistema}
\end{equation}

\begin{equation}
\dot{\mathbf{x}}_e
=
\mathbf{f}_e
\left(
\mathbf{x}_e,
\mathbf{w}
\right),
\label{eq:dinamica_entorno}
\end{equation}

donde $\mathbf{u}$ representa las acciones de control, $\boldsymbol{\theta}$ los parámetros físicos y $\mathbf{w}$ las excitaciones externas que gobiernan la evolución meteorológica.

Esta representación permite distinguir entre:

\begin{itemize}
    \item estados internos del proceso;
    \item perturbaciones medibles;
    \item perturbaciones no medibles;
    \item acciones manipulables;
    \item parámetros estructurales.
\end{itemize}

%%%%%%%%%%%%%%%%%%%%%%%%%%%%%%%%%%%%%%%%%%%%%%%%%%%%%%%%%%%%%%%%%%%%%%%%%%%%%%%%
\subsection{Interacciones físicas fundamentales}
\label{subsec:interacciones_fundamentales}

El comportamiento del secador emerge del acoplamiento entre transferencia de calor, flujo de aire y transporte de humedad.

La radiación solar suministra energía a las superficies y al aire interior. La conducción redistribuye esta energía dentro de los componentes sólidos. La convección transfiere energía entre las superficies y el aire. La evaporación consume energía y aumenta la cantidad de vapor presente en el dominio fluido. Finalmente, la ventilación transporta calor y humedad hacia el exterior.

La relación conceptual puede representarse como

\begin{equation}
\boxed{
\begin{aligned}
\text{Radiación solar}
&\longrightarrow
\text{calentamiento de superficies}
\\
&\longrightarrow
\text{convección hacia el aire}
\\
&\longrightarrow
\text{calentamiento del producto}
\\
&\longrightarrow
\text{evaporación de humedad}
\\
&\longrightarrow
\text{transporte de vapor}.
\end{aligned}
}
\label{eq:cadena_fisica_secado}
\end{equation}

No obstante, el proceso contiene múltiples retroalimentaciones físicas. Por ejemplo,

\begin{equation}
\text{evaporación}
\longrightarrow
\text{enfriamiento del producto},
\end{equation}

\begin{equation}
\text{ventilación}
\longrightarrow
\begin{cases}
\text{mayor extracción de humedad},\\
\text{mayores pérdidas térmicas},
\end{cases}
\end{equation}

y

\begin{equation}
\text{radiación superficial}
\longrightarrow
\text{modificación de los gradientes convectivos}.
\end{equation}

En consecuencia, los mecanismos físicos no deben analizarse como procesos independientes cuya contribución se suma al final. Cada mecanismo modifica las condiciones bajo las cuales actúan los demás.

%%%%%%%%%%%%%%%%%%%%%%%%%%%%%%%%%%%%%%%%%%%%%%%%%%%%%%%%%%%%%%%%%%%%%%%%%%%%%%%%
\subsection{Variables conceptuales de la investigación}
\label{subsec:variables_conceptuales}

Para organizar el problema se distinguen cinco clases de variables.

\subsubsection{Parámetros de diseño}

Los parámetros de diseño describen características estructurales o relativamente invariantes del secador:

\begin{equation}
\boldsymbol{\theta}_d
=
\begin{bmatrix}
\text{geometría} &
\text{materiales} &
\text{espesores} &
\text{propiedades ópticas} &
\text{áreas de ventilación}
\end{bmatrix}^{T}.
\end{equation}

Estas variables determinan la capacidad del sistema para captar, almacenar, distribuir y disipar energía.

\subsubsection{Condiciones iniciales}

Las condiciones iniciales representan el estado del sistema al comenzar la operación:

\begin{equation}
\mathbf{x}(t_0)=\mathbf{x}_0.
\end{equation}

Incluyen las temperaturas iniciales, el contenido inicial de humedad del producto y la condición inicial del aire interno.

\subsubsection{Perturbaciones}

Las perturbaciones representan influencias externas que no son determinadas por el controlador:

\begin{equation}
\mathbf{d}(t)
=
\begin{bmatrix}
I_t &
T_{amb} &
H_{amb} &
v_w
\end{bmatrix}^{T}.
\end{equation}

\subsubsection{Variables manipuladas}

Las variables manipuladas corresponden a las acciones que pueden modificar deliberadamente la evolución del sistema:

\begin{equation}
\mathbf{u}(t)
=
\begin{bmatrix}
u_1(t) &
u_2(t) &
\cdots &
u_m(t)
\end{bmatrix}^{T}.
\end{equation}

Estas acciones pueden estar asociadas con ventilación, extracción, apertura de compuertas u otros mecanismos de regulación.

\subsubsection{Variables de desempeño}

Las variables de desempeño permiten evaluar si la operación satisface los objetivos del proceso. Entre ellas se encuentran:

\begin{itemize}
    \item temperatura media;
    \item humedad media;
    \item contenido de humedad del producto;
    \item velocidad de secado;
    \item uniformidad térmica;
    \item uniformidad higrométrica;
    \item consumo energético;
    \item duración del proceso;
    \item exposición a condiciones potencialmente perjudiciales.
\end{itemize}

%%%%%%%%%%%%%%%%%%%%%%%%%%%%%%%%%%%%%%%%%%%%%%%%%%%%%%%%%%%%%%%%%%%%%%%%%%%%%%%%
\subsection{Uniformidad como criterio de desempeño}
\label{subsec:uniformidad_desempeno}

Una de las ideas centrales de esta investigación es que el control de un secador no debe evaluarse únicamente mediante valores promedio.

Una temperatura media adecuada no garantiza que todas las regiones del producto se encuentren sometidas a condiciones semejantes. Pueden existir simultáneamente zonas sobrecalentadas, regiones con secado insuficiente o acumulaciones locales de humedad.

La temperatura promedio se define como

\begin{equation}
\overline{T}(t)
=
\frac{1}{V}
\int_{\Omega}
T(\mathbf{r},t)\,dV,
\end{equation}

mientras que la dispersión espacial puede representarse mediante

\begin{equation}
\sigma_T^2(t)
=
\frac{1}{V}
\int_{\Omega}
\left[
T(\mathbf{r},t)-\overline{T}(t)
\right]^2
\,dV.
\end{equation}

De forma análoga,

\begin{equation}
\sigma_H^2(t)
=
\frac{1}{V}
\int_{\Omega}
\left[
H(\mathbf{r},t)-\overline{H}(t)
\right]^2
\,dV.
\end{equation}

El desempeño termo-higrométrico se interpreta entonces como un compromiso entre seguimiento de referencias y reducción de heterogeneidades espaciales:

\begin{equation}
\boxed{
J_{th}
=
J_{\text{seguimiento}}
+
J_{\text{uniformidad}}.
}
\label{eq:desempeno_termohigrometrico}
\end{equation}

Este planteamiento transforma el problema desde el control de una variable escalar hacia la regulación aproximada de un sistema distribuido.

%%%%%%%%%%%%%%%%%%%%%%%%%%%%%%%%%%%%%%%%%%%%%%%%%%%%%%%%%%%%%%%%%%%%%%%%%%%%%%%%
\subsection{Jerarquía de modelos}
\label{subsec:jerarquia_modelos}

La investigación no adopta un único modelo como representación absoluta del secador. En su lugar, utiliza una jerarquía de modelos con diferentes niveles de fidelidad, interpretabilidad y costo computacional.

\begin{equation}
\boxed{
\begin{aligned}
&\text{Modelo físico distribuido de alta fidelidad}
\\
&\hspace{2cm}\Downarrow
\\
&\text{Modelo físico reducido}
\\
&\hspace{2cm}\Downarrow
\\
&\text{Modelo basado en datos o modelo híbrido}
\\
&\hspace{2cm}\Downarrow
\\
&\text{Modelo operativo para estimación y control}
\end{aligned}
}
\label{eq:jerarquia_modelos}
\end{equation}

Cada nivel cumple una función distinta.

\begin{table}[H]
\centering
\caption{Función conceptual de los diferentes niveles de modelado.}
\label{tab:jerarquia_conceptual_modelos}
\renewcommand{\arraystretch}{1.35}
\begin{tabular}{
p{0.24\textwidth}
p{0.31\textwidth}
p{0.34\textwidth}
}
\toprule
\textbf{Nivel de modelado}
&
\textbf{Propósito principal}
&
\textbf{Característica dominante}
\\
\midrule

Modelo CFD
&
Representar campos espaciales y analizar mecanismos físicos locales.
&
Alta fidelidad espacial y elevado costo computacional.
\\

Modelo concentrado
&
Describir la dinámica dominante mediante un número reducido de estados.
&
Interpretabilidad física y menor complejidad.
\\

Modelo basado en datos
&
Aproximar relaciones no lineales a partir de observaciones o simulaciones.
&
Rapidez de inferencia y dependencia del dominio de entrenamiento.
\\

Modelo híbrido
&
Combinar ecuaciones físicas con correcciones o parámetros aprendidos.
&
Equilibrio entre estructura física y flexibilidad empírica.
\\

Modelo de control
&
Predecir la evolución relevante para seleccionar acciones.
&
Rapidez, estabilidad y compatibilidad con la optimización.
\\
\bottomrule
\end{tabular}
\end{table}

La utilización de varios modelos no constituye una redundancia. Responde a la imposibilidad práctica de maximizar simultáneamente fidelidad, velocidad, interpretabilidad y facilidad de control mediante una única representación.

%%%%%%%%%%%%%%%%%%%%%%%%%%%%%%%%%%%%%%%%%%%%%%%%%%%%%%%%%%%%%%%%%%%%%%%%%%%%%%%%
\subsection{Relación entre simulación, experimento y aprendizaje}
\label{subsec:simulacion_experimento_aprendizaje}

La evidencia empleada en la investigación procede de tres fuentes principales:

\begin{enumerate}
    \item principios físicos;
    \item simulaciones numéricas;
    \item mediciones experimentales.
\end{enumerate}

Estas fuentes no tienen el mismo estatus epistemológico.

Las leyes de conservación proporcionan la estructura física general. Las simulaciones permiten explorar esa estructura bajo condiciones definidas, pero sus resultados dependen de las hipótesis, parámetros, condiciones de frontera y aproximaciones numéricas adoptadas. Las mediciones representan la respuesta del sistema real, aunque se encuentran afectadas por ruido, cobertura espacial limitada e incertidumbre instrumental.

La relación adecuada entre ellas es

\begin{equation}
\boxed{
\text{Física}
\longrightarrow
\text{modelo}
\longrightarrow
\text{simulación}
\longleftrightarrow
\text{experimento}
\longrightarrow
\text{revisión del modelo}.
}
\label{eq:ciclo_validacion}
\end{equation}

Por tanto, la simulación no sustituye la validación experimental y los datos experimentales no eliminan la necesidad de una estructura física.

El aprendizaje automático puede incorporarse después de establecer esta relación:

\begin{equation}
\boxed{
\text{Física}
+
\text{simulación}
+
\text{experimento}
\longrightarrow
\text{modelo aprendido o híbrido}.
}
\label{eq:fuentes_modelo_hibrido}
\end{equation}

%%%%%%%%%%%%%%%%%%%%%%%%%%%%%%%%%%%%%%%%%%%%%%%%%%%%%%%%%%%%%%%%%%%%%%%%%%%%%%%%
\subsection{Dominio operacional y validez}
\label{subsec:dominio_operacional_validez}

Todo modelo utilizado en esta investigación posee un dominio de validez limitado.

Sea $\mathcal{D}_{op}$ el conjunto de condiciones de operación definido por combinaciones de variables ambientales, estados internos, configuraciones de carga y acciones de control:

\begin{equation}
\mathcal{D}_{op}
=
\left\{
\mathbf{x},
\mathbf{u},
\mathbf{d},
\boldsymbol{\theta}
\right\}.
\end{equation}

Dentro de este dominio pueden distinguirse tres regiones:

\begin{equation}
\mathcal{D}_{val}
\subseteq
\mathcal{D}_{sim}
\subseteq
\mathcal{D}_{op},
\end{equation}

donde

\begin{itemize}
    \item $\mathcal{D}_{val}$ es la región respaldada por validación experimental;
    \item $\mathcal{D}_{sim}$ es la región explorada mediante simulación;
    \item $\mathcal{D}_{op}$ es el dominio operacional de interés.
\end{itemize}

Debe distinguirse entre:

\begin{equation}
\text{interpolación},
\qquad
\text{generalización},
\qquad
\text{extrapolación}.
\end{equation}

La interpolación ocurre dentro de una región suficientemente cubierta por datos. La generalización corresponde al desempeño sobre condiciones no utilizadas directamente durante el ajuste, pero pertenecientes a un régimen semejante. La extrapolación ocurre fuera de las regiones respaldadas por evidencia suficiente.

\begin{tcolorbox}[
title=Principio conceptual de validez,
colback=yellow!5,
colframe=orange!70!black,
boxrule=0.7pt,
arc=2mm
]
La existencia de una predicción numérica no implica que dicha predicción sea válida.

La confiabilidad depende de la correspondencia entre el modelo y el fenómeno físico, de la calidad de los datos, de la adecuada definición de las condiciones de frontera y de la proximidad entre la condición evaluada y el dominio respaldado por evidencia.
\end{tcolorbox}

%%%%%%%%%%%%%%%%%%%%%%%%%%%%%%%%%%%%%%%%%%%%%%%%%%%%%%%%%%%%%%%%%%%%%%%%%%%%%%%%
\subsection{Generalización entre configuraciones de secadores}
\label{subsec:generalizacion_secadores}

La respuesta térmica local de un secador depende de su geometría, materiales, ubicación, clima y condiciones de operación. Por ello, no se plantea la existencia de una solución universal que reproduzca directamente todos los campos termo-higrométricos para cualquier secador.

La generalización propuesta se encuentra en un nivel diferente: se plantea que distintas configuraciones pueden describirse mediante una estructura común de variables, balances y grupos de parámetros.

En términos conceptuales,

\begin{equation}
\boxed{
\text{No se generaliza una solución particular;}
\qquad
\text{se generaliza la estructura del problema.}
}
\label{eq:generalizacion_estructura}
\end{equation}

Una familia de secadores puede representarse mediante

\begin{equation}
\mathcal{F}
=
\left\{
\mathcal{M}
\left(
\boldsymbol{\theta}_d,
\mathbf{x}_0,
\mathbf{d},
\mathbf{u}
\right)
\right\},
\end{equation}

donde la estructura funcional $\mathcal{M}$ se conserva, mientras que los parámetros de diseño y operación cambian entre configuraciones.

Esta formulación permite estudiar similitudes, escalabilidad y transferencia de modelos sin asumir que dos geometrías diferentes producen campos idénticos.

%%%%%%%%%%%%%%%%%%%%%%%%%%%%%%%%%%%%%%%%%%%%%%%%%%%%%%%%%%%%%%%%%%%%%%%%%%%%%%%%
\subsection{Observación mínima del proceso}
\label{subsec:observacion_minima}

El control de un sistema requiere información suficiente sobre su estado. Sin embargo, una instrumentación extensa puede resultar económicamente inviable o difícil de mantener.

Por ello, el marco conceptual introduce el problema de determinar un conjunto reducido de observables capaz de preservar información suficiente para la estimación y el control.

Sea

\begin{equation}
\mathcal{S}_y
=
\{y_1,y_2,\ldots,y_p\}
\end{equation}

el conjunto total de mediciones potenciales. Se busca identificar un subconjunto

\begin{equation}
\mathcal{S}_y^{*}
\subseteq
\mathcal{S}_y
\end{equation}

que proporcione información suficiente para reconstruir o aproximar los estados relevantes:

\begin{equation}
\hat{\mathbf{x}}
=
\mathcal{E}
\left(
\mathcal{S}_y^{*},
\mathbf{u},
\mathbf{d}
\right).
\end{equation}

El concepto de observación mínima no significa utilizar la menor cantidad absoluta de sensores, sino alcanzar un equilibrio entre información, costo, redundancia y robustez.

Este equilibrio puede representarse mediante

\begin{equation}
\boxed{
\mathcal{S}_y^{*}
=
\arg\min_{\mathcal{S}_y}
\left[
J_{\mathrm{estimación}}
+
\lambda_y
\left|
\mathcal{S}_y
\right|
\right].
}
\label{eq:seleccion_sensores}
\end{equation}

%%%%%%%%%%%%%%%%%%%%%%%%%%%%%%%%%%%%%%%%%%%%%%%%%%%%%%%%%%%%%%%%%%%%%%%%%%%%%%%%
\subsection{Actuación mínima y controlabilidad práctica}
\label{subsec:actuacion_minima}

De manera análoga, no todas las acciones físicamente posibles resultan necesarias o económicamente justificables.

Sea

\begin{equation}
\mathcal{S}_u
=
\{u_1,u_2,\ldots,u_m\}
\end{equation}

el conjunto de actuadores potenciales. El problema consiste en identificar una arquitectura reducida

\begin{equation}
\mathcal{S}_u^{*}
\subseteq
\mathcal{S}_u
\end{equation}

que permita mantener un desempeño aceptable.

El concepto de controlabilidad utilizado en este contexto es principalmente práctico: se refiere a la capacidad de modificar suficientemente las variables relevantes dentro de los límites físicos, económicos y temporales del proceso.

El problema conjunto de instrumentación y actuación puede expresarse como

\begin{equation}
\boxed{
\min_{\mathcal{S}_y,\mathcal{S}_u}
\left[
J_{\mathrm{proceso}}
+
\lambda_y|\mathcal{S}_y|
+
\lambda_u|\mathcal{S}_u|
\right].
}
\label{eq:arquitectura_minima}
\end{equation}

Esta formulación conecta directamente el análisis físico con la viabilidad tecnológica del sistema.

%%%%%%%%%%%%%%%%%%%%%%%%%%%%%%%%%%%%%%%%%%%%%%%%%%%%%%%%%%%%%%%%%%%%%%%%%%%%%%%%
\subsection{Concepto de control inteligente}
\label{subsec:concepto_control_inteligente}

En esta investigación, el término \textit{control inteligente} no se utiliza como sinónimo general de automatización ni como referencia exclusiva al uso de redes neuronales.

Se define como una estrategia capaz de integrar:

\begin{itemize}
    \item mediciones del proceso;
    \item estimación de estados no observados;
    \item conocimiento físico;
    \item predicción de la evolución futura;
    \item consideración explícita de perturbaciones;
    \item optimización de acciones;
    \item adaptación ante discrepancias del modelo.
\end{itemize}

Por tanto,

\begin{equation}
\boxed{
\text{Control inteligente}
=
\text{observación}
+
\text{modelo}
+
\text{predicción}
+
\text{decisión}
+
\text{adaptación}.
}
\label{eq:definicion_control_inteligente}
\end{equation}

La inteligencia del sistema no reside únicamente en el algoritmo empleado, sino en la coherencia con la que se integran información, física, incertidumbre y objetivos de operación.

%%%%%%%%%%%%%%%%%%%%%%%%%%%%%%%%%%%%%%%%%%%%%%%%%%%%%%%%%%%%%%%%%%%%%%%%%%%%%%%%
\subsection{Integración de modelos físicos y modelos basados en datos}
\label{subsec:integracion_fisica_datos}

El marco conceptual adopta una filosofía híbrida. Los modelos físicos aportan estructura, interpretabilidad y restricciones derivadas de principios de conservación. Los modelos basados en datos aportan capacidad para aproximar relaciones no lineales, corregir discrepancias y realizar inferencias rápidas.

Esta relación puede expresarse mediante

\begin{equation}
\boxed{
\mathbf{f}_{hib}
=
\mathbf{f}_{fis}
+
\mathbf{f}_{data}.
}
\label{eq:modelo_hibrido_conceptual}
\end{equation}

El término basado en datos puede cumplir diferentes funciones:

\begin{enumerate}
    \item estimar parámetros desconocidos;
    \item corregir residuos del modelo físico;
    \item sustituir temporalmente un modelo costoso;
    \item reconstruir estados no medidos;
    \item aproximar campos distribuidos;
    \item acelerar la predicción dentro de un controlador.
\end{enumerate}

La jerarquía conceptual adoptada es

\begin{equation}
\boxed{
\text{La física define la estructura;}
\qquad
\text{los datos ajustan, complementan o corrigen.}
}
\label{eq:filosofia_hibrida}
\end{equation}

Esta filosofía evita dos extremos: confiar exclusivamente en un modelo físico incapaz de representar todas las discrepancias reales, o emplear un modelo puramente empírico sin restricciones de consistencia.

%%%%%%%%%%%%%%%%%%%%%%%%%%%%%%%%%%%%%%%%%%%%%%%%%%%%%%%%%%%%%%%%%%%%%%%%%%%%%%%%
\subsection{Problema de investigación dentro del marco conceptual}
\label{subsec:problema_marco_conceptual}

El problema de investigación surge de la dificultad para regular de manera uniforme y eficiente un proceso de secado impulsado por una fuente energética variable y sometido a condiciones ambientales no controlables.

De manera sintética, el problema puede expresarse como:

\begin{tcolorbox}[
title=Problema central de investigación,
colback=red!3,
colframe=red!55!black,
boxrule=0.8pt,
arc=2mm
]
¿Cómo desarrollar una representación integrada, físicamente fundamentada y computacionalmente utilizable de un secador solar de café que permita estimar su estado termo-higrométrico y seleccionar acciones de control capaces de mantener condiciones de secado adecuadas y uniformes frente a perturbaciones meteorológicas?
\end{tcolorbox}

Este problema contiene cuatro dificultades interrelacionadas:

\begin{enumerate}
    \item la dinámica es no lineal y acoplada;
    \item las variables presentan heterogeneidad espacial;
    \item no todos los estados pueden medirse directamente;
    \item los modelos de alta fidelidad no siempre son utilizables en tiempo real.
\end{enumerate}

%%%%%%%%%%%%%%%%%%%%%%%%%%%%%%%%%%%%%%%%%%%%%%%%%%%%%%%%%%%%%%%%%%%%%%%%%%%%%%%%
\subsection{Objetivo general dentro del marco conceptual}
\label{subsec:objetivo_general_marco}

El objetivo general puede formularse conceptualmente como:

\begin{tcolorbox}[
title=Objetivo general,
colback=blue!3,
colframe=blue!60!black,
boxrule=0.8pt,
arc=2mm
]
Desarrollar un marco integrado de modelado, validación, estimación y control para la caracterización termo-higrométrica de un secador solar de café, mediante la combinación de modelos físicos, simulaciones numéricas, mediciones experimentales y modelos basados en datos.
\end{tcolorbox}

Este objetivo no presupone que todos los niveles de modelado deban tener la misma función. El modelo CFD se orienta principalmente a la caracterización espacial y al análisis físico; los modelos reducidos e híbridos se orientan a la predicción, estimación y control.

%%%%%%%%%%%%%%%%%%%%%%%%%%%%%%%%%%%%%%%%%%%%%%%%%%%%%%%%%%%%%%%%%%%%%%%%%%%%%%%%
\subsection{Objetivos conceptuales específicos}
\label{subsec:objetivos_conceptuales}

Los objetivos específicos pueden organizarse según la secuencia lógica del marco conceptual:

\begin{enumerate}
    \item caracterizar los mecanismos de transferencia de calor y masa dominantes;
    \item construir una representación física del proceso;
    \item estudiar la distribución espacial de temperatura, velocidad y humedad;
    \item obtener evidencia experimental para evaluar la representación numérica;
    \item desarrollar modelos reducidos o sustitutos compatibles con la predicción rápida;
    \item estimar estados que no pueden medirse directamente;
    \item formular estrategias de control para temperatura, humedad y uniformidad;
    \item evaluar arquitecturas reducidas de observación y actuación.
\end{enumerate}

La secuencia conceptual es, por tanto,

\begin{equation}
\boxed{
\text{comprender}
\rightarrow
\text{modelar}
\rightarrow
\text{validar}
\rightarrow
\text{reducir}
\rightarrow
\text{estimar}
\rightarrow
\text{controlar}.
}
\label{eq:secuencia_objetivos}
\end{equation}

%%%%%%%%%%%%%%%%%%%%%%%%%%%%%%%%%%%%%%%%%%%%%%%%%%%%%%%%%%%%%%%%%%%%%%%%%%%%%%%%
\subsection{Postulados orientadores}
\label{subsec:postulados_orientadores}

El marco conceptual se apoya en los siguientes postulados de trabajo.

\subsubsection*{Postulado 1: naturaleza sistémica}

El desempeño del secador es una propiedad emergente de la interacción entre diseño, ambiente, estados internos, carga de producto y acciones de control.

\begin{equation}
\mathcal{P}
=
\mathcal{G}
\left(
\boldsymbol{\theta}_d,
\mathbf{x}_0,
\mathbf{d},
\mathbf{u}
\right),
\end{equation}

donde $\mathcal{P}$ representa el desempeño global.

\subsubsection*{Postulado 2: acoplamiento termo-higrométrico}

La temperatura, la humedad del aire y la humedad del producto no pueden regularse como variables completamente independientes.

\begin{equation}
T
\leftrightarrow
H
\leftrightarrow
M.
\end{equation}

\subsubsection*{Postulado 3: necesidad de representación multiescala}

Ningún modelo individual satisface simultáneamente los requisitos de detalle espacial, bajo costo computacional, interpretabilidad y uso en control.

\subsubsection*{Postulado 4: complementariedad entre física y datos}

Los modelos físicos y los modelos basados en datos no constituyen alternativas excluyentes, sino representaciones complementarias.

\subsubsection*{Postulado 5: validez local condicionada}

La confiabilidad de una predicción está condicionada por el dominio de operación, la calidad de la validación y la incertidumbre de los parámetros.

\subsubsection*{Postulado 6: importancia de la uniformidad}

El desempeño no debe evaluarse únicamente mediante valores promedio, sino también mediante la distribución espacial de las variables relevantes.

\subsubsection*{Postulado 7: suficiencia antes que abundancia}

Una arquitectura eficiente debe buscar el conjunto suficiente de sensores, estados, modelos y actuadores, no necesariamente el conjunto más grande o complejo.

%%%%%%%%%%%%%%%%%%%%%%%%%%%%%%%%%%%%%%%%%%%%%%%%%%%%%%%%%%%%%%%%%%%%%%%%%%%%%%%%
\subsection{Proposiciones de investigación}
\label{subsec:proposiciones_investigacion}

A partir de los postulados anteriores se plantean las siguientes proposiciones.

\begin{enumerate}
    \item Un modelo físico distribuido permite identificar patrones espaciales y mecanismos que no pueden observarse mediante valores promedio.
    
    \item Un modelo reducido físicamente interpretable puede conservar la dinámica dominante necesaria para tareas de estimación y control.
    
    \item La incorporación de datos experimentales puede reducir discrepancias entre la representación teórica y el comportamiento real.
    
    \item Los modelos híbridos pueden alcanzar un equilibrio favorable entre fidelidad y velocidad dentro del dominio respaldado por datos.
    
    \item La anticipación de perturbaciones ambientales puede mejorar el desempeño respecto de estrategias exclusivamente reactivas.
    
    \item La selección adecuada de sensores y actuadores puede reducir la complejidad del sistema sin perder necesariamente la capacidad de regulación requerida.
    
    \item El control simultáneo de valores medios y medidas de uniformidad puede producir una operación más representativa del estado real del producto.
\end{enumerate}

Estas proposiciones constituyen expectativas razonadas que deben ser evaluadas mediante análisis, simulación y experimentación. No se consideran resultados demostrados de antemano.

%%%%%%%%%%%%%%%%%%%%%%%%%%%%%%%%%%%%%%%%%%%%%%%%%%%%%%%%%%%%%%%%%%%%%%%%%%%%%%%%
\subsection{Contribuciones conceptuales propuestas}
\label{subsec:contribuciones_conceptuales}

La contribución de la investigación no reside exclusivamente en el uso individual de CFD, control predictivo, redes neuronales o sensores. Estas herramientas ya se encuentran establecidas en sus respectivos campos.

La contribución propuesta se localiza en la integración coherente de estos elementos dentro de una arquitectura orientada al secado solar de café.

Las principales contribuciones conceptuales son:

\begin{enumerate}
    \item representar conjuntamente la planta y las perturbaciones meteorológicas mediante un estado aumentado sistema--entorno;
    
    \item integrar niveles de modelado con diferentes grados de fidelidad y distintas funciones dentro de la investigación;
    
    \item formular el desempeño en términos de seguimiento y uniformidad espacial;
    
    \item diferenciar explícitamente entre dominio validado, dominio simulado y dominio extrapolado;
    
    \item utilizar modelos híbridos como puente entre la simulación de alta fidelidad y la operación en tiempo real;
    
    \item plantear la observación y la actuación como problemas de suficiencia y viabilidad económica;
    
    \item estructurar el control inteligente como una combinación de observación, predicción, conocimiento físico, optimización y adaptación.
\end{enumerate}

La contribución integrada puede sintetizarse mediante

\begin{tcolorbox}[
title=Contribución conceptual integrada,
colback=green!3,
colframe=green!50!black,
boxrule=0.8pt,
arc=2mm
]
Desarrollar una arquitectura multiescala y físico--data-driven para caracterizar, predecir y controlar la dinámica termo-higrométrica de secadores solares de café, con énfasis en la uniformidad del proceso, la validez dentro del dominio operacional y la reducción de la complejidad de observación y actuación.
\end{tcolorbox}

%%%%%%%%%%%%%%%%%%%%%%%%%%%%%%%%%%%%%%%%%%%%%%%%%%%%%%%%%%%%%%%%%%%%%%%%%%%%%%%%
\subsection{Filosofía de modelado}
\label{subsec:filosofia_modelado}

La filosofía de esta investigación se fundamenta en cinco principios.

\subsubsection*{Primero: comenzar por la física}

La selección de variables, modelos y algoritmos debe partir de los mecanismos de conservación y transporte, no únicamente de la disponibilidad de datos.

\subsubsection*{Segundo: utilizar el nivel de fidelidad adecuado}

La mayor complejidad no implica necesariamente una mejor representación para todas las tareas. El modelo debe ser adecuado al propósito para el cual se utiliza.

\subsubsection*{Tercero: distinguir predicción de explicación}

Un modelo puede predecir correctamente sin representar explícitamente todos los mecanismos físicos. De manera inversa, un modelo físicamente interpretable puede presentar errores cuantitativos. Ambas propiedades deben evaluarse por separado.

\subsubsection*{Cuarto: validar antes de generalizar}

Las afirmaciones sobre desempeño deben restringirse a las regiones respaldadas por evidencia suficiente. La extrapolación debe identificarse explícitamente como tal.

\subsubsection*{Quinto: diseñar para la implementación real}

La utilidad del modelo debe evaluarse no sólo por su precisión, sino también por sus requisitos de instrumentación, costo computacional, robustez y capacidad de operación.

Esta filosofía puede resumirse como

\begin{equation}
\boxed{
\text{Rigor físico}
+
\text{evidencia experimental}
+
\text{viabilidad computacional}
+
\text{utilidad operativa}.
}
\label{eq:filosofia_general}
\end{equation}

%%%%%%%%%%%%%%%%%%%%%%%%%%%%%%%%%%%%%%%%%%%%%%%%%%%%%%%%%%%%%%%%%%%%%%%%%%%%%%%%
\subsection{Framework conceptual unificado}
\label{subsec:framework_unificado}

La totalidad del marco conceptual puede sintetizarse mediante la siguiente arquitectura:

\begin{equation}
\boxed{
\begin{aligned}
&\textbf{ENTORNO}
\\
&\left[
I_t,\,
T_{amb},\,
H_{amb},\,
v_w
\right]
\\[0.15cm]
&\qquad\Downarrow
\\[0.15cm]
&\textbf{PLANTA TERMO-HIGROMÉTRICA}
\\
&\left[
\text{radiación}
+
\text{conducción}
+
\text{convección}
+
\text{evaporación}
+
\text{transporte de vapor}
\right]
\\[0.15cm]
&\qquad\Downarrow
\\[0.15cm]
&\textbf{ESTADOS}
\\
&\left[
T,\,
H,\,
M,\,
\mathbf{v},\,
p
\right]
\\[0.15cm]
&\qquad\Downarrow
\\[0.15cm]
&\textbf{OBSERVACIÓN}
\\
&\left[
\text{sensores}
+
\text{estimadores}
\right]
\\[0.15cm]
&\qquad\Downarrow
\\[0.15cm]
&\textbf{JERARQUÍA DE MODELOS}
\\
&\left[
\text{CFD}
+
\text{modelo reducido}
+
\text{modelo basado en datos}
+
\text{modelo híbrido}
\right]
\\[0.15cm]
&\qquad\Downarrow
\\[0.15cm]
&\textbf{PREDICCIÓN Y CONTROL}
\\
&\left[
\text{control clásico}
+
\text{MPC}
+
\text{control híbrido}
\right]
\\[0.15cm]
&\qquad\Downarrow
\\[0.15cm]
&\textbf{DESEMPEÑO}
\\
&\left[
\text{seguimiento}
+
\text{uniformidad}
+
\text{eficiencia}
+
\text{calidad}
\right].
\end{aligned}
}
\label{eq:framework_conceptual_unificado}
\end{equation}

La realimentación cierra esta arquitectura:

\begin{equation}
\boxed{
\text{Desempeño observado}
\longrightarrow
\text{actualización del estado y del modelo}
\longrightarrow
\text{nueva acción de control}.
}
\end{equation}

El sistema completo se interpreta, por tanto, como un ciclo continuo de observación, predicción, decisión y corrección.

%%%%%%%%%%%%%%%%%%%%%%%%%%%%%%%%%%%%%%%%%%%%%%%%%%%%%%%%%%%%%%%%%%%%%%%%%%%%%%%%
\subsection{Síntesis del marco conceptual}
\label{subsec:sintesis_marco_conceptual}

El marco conceptual desarrollado establece que el secado solar de café es un problema de dinámica termo-higrométrica acoplada y no únicamente de calentamiento del aire.

La radiación solar constituye la fuente energética primaria, pero su aprovechamiento depende de la geometría, las propiedades de los materiales, el movimiento del aire, la carga de producto y las condiciones atmosféricas. El estado del proceso no puede caracterizarse completamente mediante una única medición, y su control requiere combinar observación, estimación y predicción.

La investigación adopta una jerarquía de modelos. Los modelos físicos de alta fidelidad permiten comprender los campos y mecanismos locales. Los modelos reducidos facilitan la interpretación dinámica. Los modelos basados en datos y los modelos híbridos proporcionan aproximaciones rápidas para predicción y control. La experimentación establece el vínculo entre estas representaciones y el sistema real.

La filosofía central puede expresarse finalmente como

\begin{equation}
\boxed{
\begin{aligned}
\text{Comprender la física}
&\longrightarrow
\text{construir modelos}
\\
&\longrightarrow
\text{contrastar con evidencia}
\\
&\longrightarrow
\text{reducir la complejidad}
\\
&\longrightarrow
\text{predecir el comportamiento}
\\
&\longrightarrow
\text{seleccionar acciones}
\\
&\longrightarrow
\text{mejorar el proceso}.
\end{aligned}
}
\label{eq:cierre_marco_conceptual}
\end{equation}

Sobre esta estructura conceptual se fundamenta la metodología de la investigación, en la cual se especifican los modelos seleccionados, la instrumentación empleada, los procedimientos de simulación, los métodos de validación y las estrategias de control evaluadas.